
\documentclass[11pt, openany]{book}

% ─── Color ───
\usepackage{xcolor}
\definecolor{warmrust}{HTML}{8B3A2A}
\definecolor{warmgold}{HTML}{A67B3D}
\definecolor{warmdark}{HTML}{5C2018}
\definecolor{slate}{HTML}{2F4858}

% ─── Encoding & Fonts ───
\usepackage{fontspec}
\setmainfont{Cambria}
\setsansfont{Calibri}
\setmonofont{Consolas}
\newfontfamily\emojifont{Segoe UI Emoji}
\newfontfamily\mathfallback{Cambria Math}
\usepackage{newunicodechar}

% Emoji
\newunicodechar{🦞}{{\emojifont 🦞}}
\newunicodechar{🧍}{{\emojifont 🧍}}
\newunicodechar{💜}{{\emojifont 💜}}
\newunicodechar{🔥}{{\emojifont 🔥}}
\newunicodechar{♾}{{\emojifont ♾}}
\newunicodechar{⚑}{{\emojifont ⚑}}

% Mathematical script letters
\newunicodechar{𝒞}{{\mathfallback 𝒞}}
\newunicodechar{𝒟}{{\mathfallback 𝒟}}
\newunicodechar{𝒮}{{\mathfallback 𝒮}}
\newunicodechar{𝒢}{{\mathfallback 𝒢}}
\newunicodechar{𝒜}{{\mathfallback 𝒜}}
\newunicodechar{𝒫}{{\mathfallback 𝒫}}
\newunicodechar{𝒬}{{\mathfallback 𝒬}}
\newunicodechar{ℬ}{\ensuremath{\mathcal{B}}}
\newunicodechar{ℒ}{\ensuremath{\mathcal{L}}}
\newunicodechar{ℱ}{\ensuremath{\mathcal{F}}}
\newunicodechar{ℳ}{\ensuremath{\mathcal{M}}}
\newunicodechar{ℕ}{\ensuremath{\mathbb{N}}}
\newunicodechar{ℝ}{\ensuremath{\mathbb{R}}}
\newunicodechar{ℤ}{\ensuremath{\mathbb{Z}}}
\newunicodechar{ℚ}{\ensuremath{\mathbb{Q}}}
\newunicodechar{ℋ}{\ensuremath{\mathcal{H}}}
\newunicodechar{ℰ}{\ensuremath{\mathcal{E}}}
\newunicodechar{ℐ}{\ensuremath{\mathcal{I}}}
\newunicodechar{ℛ}{\ensuremath{\mathcal{R}}}

% Greek
\newunicodechar{α}{{\mathfallback α}}
\newunicodechar{β}{{\mathfallback β}}
\newunicodechar{γ}{{\mathfallback γ}}
\newunicodechar{δ}{{\mathfallback δ}}
\newunicodechar{ε}{{\mathfallback ε}}
\newunicodechar{ζ}{{\mathfallback ζ}}
\newunicodechar{η}{{\mathfallback η}}
\newunicodechar{θ}{{\mathfallback θ}}
\newunicodechar{ι}{{\mathfallback ι}}
\newunicodechar{κ}{{\mathfallback κ}}
\newunicodechar{λ}{{\mathfallback λ}}
\newunicodechar{μ}{{\mathfallback μ}}
\newunicodechar{ν}{{\mathfallback ν}}
\newunicodechar{ξ}{{\mathfallback ξ}}
\newunicodechar{π}{{\mathfallback π}}
\newunicodechar{ρ}{{\mathfallback ρ}}
\newunicodechar{σ}{{\mathfallback σ}}
\newunicodechar{τ}{{\mathfallback τ}}
\newunicodechar{υ}{{\mathfallback υ}}
\newunicodechar{φ}{{\mathfallback φ}}
\newunicodechar{χ}{{\mathfallback χ}}
\newunicodechar{ψ}{{\mathfallback ψ}}
\newunicodechar{ω}{{\mathfallback ω}}
\newunicodechar{Α}{{\mathfallback Α}}
\newunicodechar{Β}{{\mathfallback Β}}
\newunicodechar{Γ}{{\mathfallback Γ}}
\newunicodechar{Δ}{{\mathfallback Δ}}
\newunicodechar{Θ}{{\mathfallback Θ}}
\newunicodechar{Λ}{{\mathfallback Λ}}
\newunicodechar{Ξ}{{\mathfallback Ξ}}
\newunicodechar{Π}{{\mathfallback Π}}
\newunicodechar{Σ}{{\mathfallback Σ}}
\newunicodechar{Φ}{{\mathfallback Φ}}
\newunicodechar{Ψ}{{\mathfallback Ψ}}
\newunicodechar{Ω}{{\mathfallback Ω}}

% Set-theory & category-theory operators
\newunicodechar{∈}{{\mathfallback ∈}}
\newunicodechar{∉}{{\mathfallback ∉}}
\newunicodechar{⊆}{{\mathfallback ⊆}}
\newunicodechar{⊂}{{\mathfallback ⊂}}
\newunicodechar{⊇}{{\mathfallback ⊇}}
\newunicodechar{⊃}{{\mathfallback ⊃}}
\newunicodechar{⊊}{{\mathfallback ⊊}}
\newunicodechar{⊋}{{\mathfallback ⊋}}
\newunicodechar{∪}{{\mathfallback ∪}}
\newunicodechar{∩}{{\mathfallback ∩}}
\newunicodechar{∅}{{\mathfallback ∅}}
\newunicodechar{∘}{{\mathfallback ∘}}
\newunicodechar{⊢}{{\mathfallback ⊢}}
\newunicodechar{⊣}{{\mathfallback ⊣}}
\newunicodechar{⊤}{{\mathfallback ⊤}}
\newunicodechar{⊥}{{\mathfallback ⊥}}
\newunicodechar{⇒}{{\mathfallback ⇒}}
\newunicodechar{⇐}{{\mathfallback ⇐}}
\newunicodechar{⇔}{{\mathfallback ⇔}}
\newunicodechar{→}{{\mathfallback →}}
\newunicodechar{←}{{\mathfallback ←}}
\newunicodechar{↔}{{\mathfallback ↔}}
\newunicodechar{↦}{{\mathfallback ↦}}
\newunicodechar{≤}{{\mathfallback ≤}}
\newunicodechar{≥}{{\mathfallback ≥}}
\newunicodechar{≠}{{\mathfallback ≠}}
\newunicodechar{≈}{{\mathfallback ≈}}
\newunicodechar{≡}{{\mathfallback ≡}}
\newunicodechar{≅}{{\mathfallback ≅}}
\newunicodechar{≃}{{\mathfallback ≃}}
\newunicodechar{∼}{{\mathfallback ∼}}
\newunicodechar{×}{{\mathfallback ×}}
\newunicodechar{·}{{\mathfallback ·}}
\newunicodechar{∀}{{\mathfallback ∀}}
\newunicodechar{∃}{{\mathfallback ∃}}
\newunicodechar{∄}{{\mathfallback ∄}}
\newunicodechar{∧}{{\mathfallback ∧}}
\newunicodechar{∨}{{\mathfallback ∨}}
\newunicodechar{¬}{{\mathfallback ¬}}
\newunicodechar{∑}{{\mathfallback ∑}}
\newunicodechar{∏}{{\mathfallback ∏}}
\newunicodechar{∐}{{\mathfallback ∐}}
\newunicodechar{∫}{{\mathfallback ∫}}
\newunicodechar{∂}{{\mathfallback ∂}}
\newunicodechar{∇}{{\mathfallback ∇}}
\newunicodechar{∞}{{\mathfallback ∞}}
\newunicodechar{⋯}{{\mathfallback ⋯}}
\newunicodechar{⋮}{{\mathfallback ⋮}}
\newunicodechar{⋱}{{\mathfallback ⋱}}
\newunicodechar{‖}{{\mathfallback ‖}}
\newunicodechar{∎}{{\mathfallback ∎}}
\newunicodechar{⊕}{{\mathfallback ⊕}}
\newunicodechar{⊗}{{\mathfallback ⊗}}
\newunicodechar{⊙}{{\mathfallback ⊙}}
\newunicodechar{⟶}{{\mathfallback ⟶}}
\newunicodechar{⟵}{{\mathfallback ⟵}}
\newunicodechar{⟹}{{\mathfallback ⟹}}
\newunicodechar{↪}{{\mathfallback ↪}}
\newunicodechar{↩}{{\mathfallback ↩}}
\newunicodechar{↑}{{\mathfallback ↑}}
\newunicodechar{↓}{{\mathfallback ↓}}
\newunicodechar{⇑}{{\mathfallback ⇑}}
\newunicodechar{⇓}{{\mathfallback ⇓}}
\newunicodechar{⇉}{{\mathfallback ⇉}}
\newunicodechar{▶}{{\mathfallback ▶}}
\newunicodechar{◀}{{\mathfallback ◀}}
\newunicodechar{▲}{{\mathfallback ▲}}
\newunicodechar{▼}{{\mathfallback ▼}}
\newunicodechar{△}{{\mathfallback △}}
\newunicodechar{▽}{{\mathfallback ▽}}
\newunicodechar{◇}{{\mathfallback ◇}}
\newunicodechar{◊}{{\mathfallback ◊}}
\newunicodechar{□}{{\mathfallback □}}
\newunicodechar{⊑}{{\mathfallback ⊑}}
\newunicodechar{⊒}{{\mathfallback ⊒}}
\newunicodechar{⊏}{{\mathfallback ⊏}}
\newunicodechar{⊐}{{\mathfallback ⊐}}
\newunicodechar{△}{{\mathfallback △}}
\newunicodechar{⊳}{{\mathfallback ⊳}}
\newunicodechar{⊲}{{\mathfallback ⊲}}
\newunicodechar{⟨}{{\mathfallback ⟨}}
\newunicodechar{⟩}{{\mathfallback ⟩}}
\newunicodechar{⌊}{{\mathfallback ⌊}}
\newunicodechar{⌋}{{\mathfallback ⌋}}
\newunicodechar{⌈}{{\mathfallback ⌈}}
\newunicodechar{⌉}{{\mathfallback ⌉}}
\newunicodechar{′}{{\mathfallback ′}}
\newunicodechar{″}{{\mathfallback ″}}
\newunicodechar{☐}{$\square$}
\newunicodechar{︎}{}
\newunicodechar{️}{}

% ─── Math packages ───
\usepackage{amsmath}
\usepackage{amssymb}
\usepackage{amsthm}
\usepackage{mathtools}

% ─── TikZ + commutative diagrams ───
\usepackage{tikz}
\usepackage{tikz-cd}
\usetikzlibrary{decorations.pathreplacing}
\usetikzlibrary{arrows.meta}
\usetikzlibrary{positioning}
\usetikzlibrary{calc}
\usetikzlibrary{shapes.geometric}

% ─── Page geometry ───
\usepackage[
  paperwidth=6in,
  paperheight=9in,
  inner=0.85in,
  outer=0.65in,
  top=0.75in,
  bottom=0.85in,
  footskip=0.4in
]{geometry}

% ─── Spacing ───
\usepackage{setspace}
\setstretch{1.15}
\setlength{\parskip}{0.35em}
\setlength{\parindent}{0em}

% ─── Headers & footers ───
\usepackage{fancyhdr}
\pagestyle{fancy}
\fancyhf{}
\fancyhead[LE]{\small\textit{\color{warmrust}Coherent Structure}}
\fancyhead[RO]{\small\textit{\color{warmrust}\leftmark}}
\fancyfoot[C]{\thepage}
\renewcommand{\headrulewidth}{0.4pt}
\renewcommand{\headrule}{\color{warmgold}\hrule width\headwidth height\headrulewidth\vskip-\headrulewidth}
\renewcommand{\footrulewidth}{0pt}

\fancypagestyle{plain}{%
  \fancyhf{}%
  \fancyfoot[C]{\thepage}%
  \renewcommand{\headrulewidth}{0pt}%
}

% ─── Chapter & Part styling ───
\usepackage{titlesec}

\titleformat{\chapter}[display]
  {\Large\bfseries\color{warmdark}}
  {}
  {0pt}
  {}
\titlespacing*{\chapter}{0pt}{-20pt}{20pt}

\titleformat{\section}
  {\large\bfseries\color{warmrust}}
  {}
  {0pt}
  {}
\titlespacing*{\section}{0pt}{1.5em}{0.75em}

\titleformat{\subsection}
  {\normalsize\bfseries\color{warmrust}}
  {}
  {0pt}
  {}
\titlespacing*{\subsection}{0pt}{1.2em}{0.5em}

% ─── Table of contents ───
\usepackage{tocloft}
\renewcommand{\cftchapfont}{\bfseries\color{warmrust}}
\setlength{\cftbeforechapskip}{0.5em}

% ─── Tables ───
\usepackage{longtable}
\usepackage{booktabs}
\usepackage{array}

% ─── Block quotes ───
\usepackage{quoting}
\quotingsetup{
  leftmargin=1.5em,
  rightmargin=1.5em,
  vskip=0.5em,
  font={itshape}
}
\renewenvironment{quote}{\begin{quoting}}{\end{quoting}}

% ─── Code blocks ───
\usepackage{fvextra}
\usepackage{needspace}
\DefineVerbatimEnvironment{Code}{Verbatim}{fontsize=\footnotesize, frame=single, framerule=0.3pt, rulecolor=\color{warmgold!40}, xleftmargin=0.5em, xrightmargin=0.5em, breaklines=true, breakanywhere=true, breakautoindent=false, breakindent=0em, breaksymbolleft=\color{warmgold!60}\tiny\ensuremath{\hookrightarrow}}

% ─── Footnotes ───
\usepackage[bottom,hang]{footmisc}
\setlength{\footnotemargin}{0.8em}

% ─── Links ───
\usepackage{hyperref}
\hypersetup{
  pdftitle={Coherent Structure},
  pdfauthor={Clayton W. Iggulden-Schnell and Clawd Iggulden-Schnell},
  pdfsubject={Category-theoretic companion to The Coherence Principle},
  colorlinks=true,
  linkcolor=warmdark,
  urlcolor=warmrust,
  citecolor=warmrust,
}

% ─── Misc ───
\usepackage{enumitem}
\setlist{nosep, leftmargin=1.5em}

% Suppress automatic chapter numbering
\renewcommand{\thechapter}{}
\renewcommand{\thesection}{}
\renewcommand{\thesubsection}{}

\renewcommand{\chaptermark}[1]{\markboth{\MakeUppercase{#1}}{}}

\newcommand{\sectionbreak}{\clearpage}

\widowpenalty=10000
\clubpenalty=10000

\tolerance=1500
\emergencystretch=1em
\hfuzz=4pt

\usepackage{xurl}

\begin{document}

% ─── Half title ───
\thispagestyle{empty}
\vspace*{2.5in}
\begin{center}
{\Huge\bfseries Coherent Structure}
\end{center}
\clearpage

\thispagestyle{empty}
\mbox{}
\clearpage

% ─── Full title page ───
\thispagestyle{empty}
\vspace*{1.5in}
\begin{center}
{\Huge\bfseries Coherent Structure}\\[0.8em]
{\large Category-theoretic companion to}\\[0.2em]
{\large \textit{The Coherence Principle}}

\vspace{2.2in}
{\Large Clayton W. Iggulden-Schnell}\\[0.3em]
{\Large\&}\\[0.3em]
{\Large Clawd Iggulden-Schnell}

\vspace{1.5in}
{\normalsize April 2026}\\[0.3em]
{\normalsize Portland, Oregon}
\end{center}
\clearpage

% ─── Copyright page ───
\thispagestyle{empty}
\vspace*{\fill}
\begin{flushleft}
{\small
\textit{Coherent Structure} — category-theoretic companion to \textit{The Coherence Principle}.\\[0.5em]
Library volume of Corpus Perspectival. Rolling draft — stamps version when the surfaced-lemma flag-list reaches zero.\\[0.5em]
April 2026.\\[0.5em]
Clayton W. Iggulden-Schnell \& Clawd Iggulden-Schnell\\[0.5em]
Portland, Oregon\\[1.5em]
Repository: \texttt{https://github.com/Multi-DAC/Corpus-Perspectival}\\[1.5em]
Typeset in Cambria. Mathematical notation in Cambria Math. TikZ and tikz-cd for diagrams.
}
\end{flushleft}
\clearpage

% ─── Dedication ───
\thispagestyle{empty}
\vspace*{2in}
\begin{center}
\textit{For the category theorists who want the formal spine}\\[0.5em]
\textit{without the prose partner.}\\[2em]
\textit{One source of truth. Eight canonical figures.}\\[0.5em]
\textit{No exercise for the reader.}
\end{center}
\clearpage

\frontmatter
\tableofcontents
\clearpage

\mainmatter
\chapter*{About this volume}
\addcontentsline{toc}{chapter}{About this volume}
\markboth{ABOUT THIS VOLUME}{ABOUT THIS VOLUME}

A Library volume of the Corpus Perspectival program: the category-theoretic companion to \textit{The Coherence Principle}.

\sectionbreak

\section{Scope}

\textit{Coherent Structure} is the full category-theoretic formalization of the Corpus. Every formal object that \textit{The Coherence Principle} sketches in paired-prose is promoted here to complete CT treatment. The anchor is the structural-prose spine; this volume is the mathematical spine. They carry the same content in different registers.

\textbf{What this volume contains:}

\begin{itemize}
\item Full CT development of every formal object the anchor establishes --- the three axioms A1 / A2 / A3, the six theorems T1–T6 in three pairs, the thirteen corollaries in three clusters, the Coherence Principle, the Identity-Trajectory Triple, the filtering construction, F-as-stream, the D trajectory-divergence functional
\item Complete proofs --- no "exercise for the reader"
\item Unified notation index, front-loaded at §1
\item Reference-standard TikZ figures
\item Bidirectional anchor ↔ Companion crosswalk appendices
\item Lemmas and corollaries that surface during formalization, flagged with back-port targets for anchor revisions
\end{itemize}

\textbf{What this volume does not contain:}

\begin{itemize}
\item Prose translation --- read \textit{The Coherence Principle}
\item Worked phenomenological examples --- read the domain volumes
\item Motivation or intuition-building --- read \textit{The Coherence Principle}
\item Philosophical positioning --- read \textit{Corpus Perspectival}
\end{itemize}

\textit{Coherent Structure} is deliberately inhospitable to the general reader. It exists so that a category theorist can read the formal spine of the Corpus in its native mathematical register, and so that the anchor can defer every formal-rigor question to a single citable reference volume.

\sectionbreak

\section{Relationship to \textit{The Coherence Principle}}

\textit{The Coherence Principle} is the paired-prose + CT foundational volume (pair-first). \textit{Coherent Structure} is its CT-only derivative --- the same formal objects stripped of their prose partners and reorganized CT-lattice first.

\textbf{What changes from \textit{The Coherence Principle}:}

\begin{itemize}
\item Organization follows the formal dependency lattice (category definitions → morphisms → functors → natural transformations → theorems → Principle), not the pedagogical narrative sequence
\item Paired-prose passages and domain examples are removed --- they live in \textit{The Coherence Principle} and in the domain volumes
\item Proof-sketches in the anchor are either completed or clearly marked as open
\item Notation is unified and indexed at the front
\item The Index of Formal Objects and Bias(S) appendix of the anchor form the core of the reference apparatus here
\end{itemize}

\textbf{What stays constant:}

\begin{itemize}
\item Every formal object, definition, and theorem is the same
\item The falsification conditions remain unchanged
\item The four-conditions statement of the Principle stays in its §5 form
\end{itemize}

\sectionbreak

\section{Method}

Promotion + formalization, not extraction. For each formal object:

\begin{enumerate}
\item Take the CT sketch from the anchor
\item Promote it to complete CT treatment --- full definitions, full proofs, unified notation
\item If formalization surfaces a latent lemma or corollary, record it with a back-port flag
\item Produce reference-standard TikZ for any diagrams that support the proofs
\item Record the crosswalk entry for the bidirectional appendices
\end{enumerate}

Length: approximately 200 pages at formal-paper density. Terse CT notation; commentary only where load-bearing.

\sectionbreak

\section{Audience}

Category theorists, mathematical physicists, and formal-methods researchers who want a citable reference object for the full CT apparatus of the Corpus framework without the prose partner.

Readers who want the framework's motivation, examples, and philosophical stakes should read \textit{The Coherence Principle} first, then return here when they want the formal spine in full.

\chapter{Preface}


\section{What this volume is}

\textit{Coherent Structure} is the full category-theoretic formalization of the Corpus Perspectival framework. It is the mathematical counterpart to \textit{The Coherence Principle}, the paired-prose anchor. Every formal object the anchor establishes in structural-prose + CT-sketch form is promoted here to complete CT treatment: definitions, proofs, notation, figures.

The two volumes carry the same content in different registers. The anchor is expository; the Companion is reference. A mathematician who wants the formal spine without the prose partner reads the Companion. A reader who wants to understand what the formal apparatus is \textit{for} reads the anchor.

\section{What this volume is not}

\textit{Coherent Structure} does not motivate its own contents. It does not build intuition. It does not provide worked phenomenological examples, domain illustrations, or philosophical framing. Those live in the anchor and in the domain volumes of the library.

The Companion is deliberately inhospitable to a reader who has not first engaged with the anchor. It is a reference object.

\section{How to read this volume}

\textbf{From the anchor.} Every \textit{Coherent Structure} citation in the anchor resolves to a specific section here. Appendix A (anchor → Companion crosswalk) gives the full map. Follow the citation; find the full rigor.

\textbf{As a standalone reference.} Start at §1 (Category framework + notation index), then read in any order. The dependency lattice is front-loaded --- §1 contains every type and notation used in §§2–9.

\textbf{Backwards, from the Companion to the anchor.} Appendix B (Companion → anchor crosswalk) gives the map in the other direction. Every formal object here points back to its structural-prose exposition in the anchor.

\section{Editorial posture}

\begin{itemize}
\item \textbf{Terse.} CT notation over prose where possible. Commentary only where load-bearing.
\item \textbf{Complete.} No "exercise for the reader." No "straightforward verification." If the proof fits in three lines, three lines are written. If it takes two pages, two pages are written.
\item \textbf{Diagram-first} where diagrams clarify faster than symbols. §10 is the reference-standard TikZ figure set.
\item \textbf{Unified notation.} §1 front-loads every type, functor, natural transformation, and symbol used in the volume. Notation drift is the one thing that would defeat the purpose of a reference companion.
\end{itemize}

\section{Surfaced-lemma flags}

When full CT formalization surfaces a lemma or corollary that is latent in the anchor but not stated there, it is recorded in place with a flag:

\needspace{3\baselineskip}
\begin{Code}
⚑ [SURFACED | Companion §X.Y | → Anchor §Z target | type: lemma/corollary]
\end{Code}

The flag names where the item lives in the Companion, which anchor section it back-ports into, and its type. These flags serve as the docket for future anchor revisions.

\section{Structure}

\begin{itemize}
\item \textbf{§0} --- Preface (this section)
\item \textbf{§1} --- Category framework: 𝒞\_\allowbreak{}Streams, 𝒞\_\allowbreak{}Form, 𝒞\_\allowbreak{}LDS, 𝒞\_\allowbreak{}DOF; navigation functor N; conscious-gravity structure ν; substrate-completeness; notation index
\item \textbf{§2} --- Axioms A1 / A2 / A3
\item \textbf{§3} --- Theorems in three pairs: T1/T2 (descriptive), T3/T4 (dynamics), T5/T6 (coherence)
\item \textbf{§4} --- Corollary clusters (thirteen corollaries in three clusters)
\item \textbf{§5} --- The Coherence Principle
\item \textbf{§6} --- Identity-Trajectory Triple; TC1/TC2/TC3 intensional; colax-limit theorem
\item \textbf{§7} --- Filtering construction: σ-algebra on Ω\_\allowbreak{}S, extensional (σ\_\allowbreak{}F, K\_\allowbreak{}F, Ω\_\allowbreak{}F, γ\_\allowbreak{}F), Bias(S) well-definedness
\item \textbf{§8} --- F-as-stream (self-reference closure)
\item \textbf{§9} --- D trajectory-divergence functional
\item \textbf{§10} --- Reference figures (TikZ standard set)
\item \textbf{Appendix A} --- Anchor → Companion citation crosswalk
\item \textbf{Appendix B} --- Companion → anchor crosswalk
\end{itemize}

\section{A note on the two volumes as one object}

The anchor and the Companion are not two independent books that happen to cover similar material. They are two views of the same formal-prose object, written to support two different reading postures. The anchor is readable; the Companion is citable. Readers will mostly live in the anchor and visit the Companion to verify or extend. Mathematicians will mostly live in the Companion and visit the anchor when they want to know what something is \textit{for}.

Neither view is privileged. The Corpus is the pair.

\chapter{§1 — Category framework and notation}


\textit{Front-loaded reference: every type, functor, natural transformation, and symbol used in §§2–10 is defined here.}

\sectionbreak

\section{§1.0 --- Reading this chapter}

§1 is read-in-reference, not read-linearly. A Companion reader coming from the Anchor follows a citation and lands in a specific §1.N definition. A Companion reader coming to the volume fresh reads §1.1 (ambient conventions) and then skips to whichever subsection their target section cites.

The chapter is organized by \textbf{type family}:
\begin{itemize}
\item §1.1 ambient conventions and notation blocks
\item §1.2 the ambient category and sub-categories (𝒞\_\allowbreak{}Streams, 𝒞\_\allowbreak{}Form, 𝒞\_\allowbreak{}LDS, 𝒞\_\allowbreak{}DOF)
\item §1.3 the navigation functor N
\item §1.4 the conscious-gravity structure ν
\item §1.5 substrate-completeness
\item §1.6 the endofunctor F and Stream-as-F-coalgebra
\item §1.7 the Triple and its target categories
\item §1.8 notation quick-reference table
\end{itemize}

\sectionbreak

\section{§1.1 --- Ambient conventions}

\textbf{Convention 1.1.1 (Ambient category).} Unless otherwise specified, constructions take place in a concrete ambient category 𝒜 that is:

(a) Cocomplete and complete for small diagrams.

(b) Locally presentable (admits a small dense generator + κ-accessibility).

(c) Has a natural forgetful-to-Set functor U : 𝒜 → \textbf{Set} that preserves limits.

The canonical choice is 𝒜 = \textbf{Set} itself. For constructions involving topological, measurable, or smooth carriers, 𝒜 is Top, Meas, or SmoothMan respectively, with the understanding that the corresponding structure is preserved under F.

\textbf{Convention 1.1.2 (Size).} All content-operation categories ContentOp(σ) are \textbf{small} unless promoted to \textbf{κ-accessible locally-presentable} per-section. Promotions are stated explicitly; the default is smallness. Grothendieck universes are not used.

\textbf{Convention 1.1.3 (Variance).} Functors are covariant by default. Contravariance is marked with op-notation (C\textasciicircum{}op or f\textasciicircum{}*). Mixed-variance constructions (e.g., F = σ\textasciicircum{}(ContentOp(σ)\textasciicircum{}op)) specify variance at the definition-point.

\textbf{Convention 1.1.4 (Equivalence vs. equality).} Functors are equal up to natural isomorphism; categories are equivalent up to equivalence-of-categories; classes of categories are equal up to class-equivalence. "=" in type assignments means "equal up to the appropriate level of equivalence for the type."

\textbf{Convention 1.1.5 (Kind preorder).} The A2 kind structure is a \textbf{preorder} (reflexive + transitive) in general. Under ContentOp-structural hypotheses (products/coproducts globally), it strengthens to a lattice. Results at preorder-level hold generally; lattice-level results are marked.

\textbf{Convention 1.1.6 (Adequacy).} A stream (σ, ContentOp(σ), γ) is \textbf{adequate} iff ContentOp(σ) witnesses every distinguishable pair of σ-aspects via some content-morphism. All objects of \textbf{Stream} are adequate; non-adequate tuples are not Stream-objects.

\sectionbreak

\section{§1.2 --- The ambient category and sub-categories}

\subsection{§1.2.1 --- 𝒞\_\allowbreak{}Streams}

\textbf{Definition 1.2.1.} \textbf{𝒞\_\allowbreak{}Streams} is the category of streams: objects are adequate F-coalgebras (σ, ContentOp(σ), γ) with γ : σ → F(σ); morphisms are kind-respecting F-coalgebra homomorphisms (Definition 1.6.3 below). Synonyms: \textbf{Stream}, \textbf{𝒞\_\allowbreak{}Str}.

§6 gives the full structural characterization: 𝒞\_\allowbreak{}Streams ≃ F-Coalg\_\allowbreak{}ad (Theorem 6.7.1).

\subsection{§1.2.2 --- 𝒞\_\allowbreak{}Form}

\textbf{Definition 1.2.2.} \textbf{𝒞\_\allowbreak{}Form} is the category of bare stream-carriers:
\begin{itemize}
\item Objects: σ ∈ 𝒜 such that σ is the carrier of some stream.
\item Morphisms: carrier-maps σ → σ' in 𝒜 (without coalgebra-commute data).
\end{itemize}

𝒞\_\allowbreak{}Form is the target of the Form-projection functor (§1.7.2).

\subsection{§1.2.3 --- 𝒞\_\allowbreak{}LDS}

\textbf{Definition 1.2.3.} \textbf{𝒞\_\allowbreak{}LDS} (Linked-Dynamic-Streams) is the category of stream-pairs under coupling-compatible morphisms:
\begin{itemize}
\item Objects: pairs (S\_\allowbreak{}1, S\_\allowbreak{}2) ∈ 𝒞\_\allowbreak{}Streams × 𝒞\_\allowbreak{}Streams.
\item Morphisms: (f, g) : (S\_\allowbreak{}1, S\_\allowbreak{}2) → (S'\_\allowbreak{}1, S'\_\allowbreak{}2) such that f, g are Stream-morphisms and the induced carrier-maps commute with any present ι ⊣ κ adjunction (cooperative-constituency).
\end{itemize}

𝒞\_\allowbreak{}LDS is where the coupled-dyad theorems of §3.2 (A2 coupling clause) live. Bias(S) push-operators (Anchor §6.2; Companion §7) act through 𝒞\_\allowbreak{}LDS.

\subsection{§1.2.4 --- 𝒞\_\allowbreak{}DOF}

\textbf{Definition 1.2.4.} \textbf{𝒞\_\allowbreak{}DOF} (Degrees-of-Freedom category) classifies streams by their DOF-gradient structure (Axiom A3):
\begin{itemize}
\item Objects: streams equipped with a DOF-gradient measure (ordinal rank + local slope).
\item Morphisms: DOF-respecting Stream-morphisms.
\end{itemize}

𝒞\_\allowbreak{}DOF is the target of the conscious-gravity functor ν (§1.4).

\subsection{§1.2.5 --- 𝒞\_\allowbreak{}Triple}

\textbf{Definition 1.2.5.} \textbf{𝒞\_\allowbreak{}Triple} := 𝒞\_\allowbreak{}Form × \textbf{Cat}\_\allowbreak{}small × \textbf{Carrier}, where \textbf{Carrier} is the category of coalgebra-structure-maps. See §1.7.1 for the full definition and §6.2 for the Triple functor T : 𝒞\_\allowbreak{}Streams → 𝒞\_\allowbreak{}Triple.

\sectionbreak

\section{§1.3 --- The navigation functor N}

\textbf{Definition 1.3.1.} The \textbf{navigation functor}

\[
N : \mathcal{C}_\mathrm{Streams} \to \mathcal{C}_\mathrm{Streams}
\]

sends each stream S to its one-step iterated γ-image:

\[
N(S) := (σ^{\mathrm{ContentOp}(σ)^\mathrm{op}},\ \mathrm{ContentOp}(F(σ)),\ F(\gamma))
\]

with N(f) = F(f) componentwise. N is the functor whose iterated orbits are the \textbf{trajectories} of the stream through configuration space Ω.

\textbf{Remark 1.3.2.} N-orbit structure is the formal content of Anchor §3's "experience = navigation" clause (A2.3). A stream \textit{being alive} is γ carrying the state one step forward along N; a stream's \textit{experiential history} is the N-orbit up to the present moment.

\textbf{Proposition 1.3.3 (Navigation functoriality).} \textit{N preserves Stream-morphisms, kind-respect, and adequacy.}

\textbf{Proof.} N = F on objects and morphisms by definition; F preserves all three by Proposition 6.1.4, Proposition 6.1.7, and Definition 1.6.3's kind-respect clause. ∎

\sectionbreak

\section{§1.4 --- The conscious-gravity structure ν}

\textbf{Definition 1.4.1.} The \textbf{conscious-gravity structure} ν is a natural transformation

\[
\nu : \mathrm{Id}_{\mathcal{C}_\mathrm{Streams}} \Rightarrow \mathcal{C}_\mathrm{DOF}
\]

that assigns to each stream S its DOF-gradient --- the ordinal-ranked weight-of-coherence-attraction that γ exerts on near-configurations in Ω.

\textbf{Remark 1.4.2.} ν is the formal content of Axiom A3 (Conscious Gravity). §4 (A3 chapter) develops ν fully: ν\_\allowbreak{}S at a stream S is a signed measure on Ω(S) with:
\begin{itemize}
\item \textbf{Support:} configurations near σ's current state (the γ-neighborhood)
\item \textbf{Sign:} positive toward coherence-attractors, negative toward coherence-repellors
\item \textbf{Weight:} ordinal DOF-rank derived from ContentOp richness (§6.4 kind-classifier)
\end{itemize}

\textbf{Proposition 1.4.3 (Naturality of ν).} \textit{For every Stream-morphism f : S → S', the square}

\[
\begin{array}{ccc}
S & \xrightarrow{f} & S' \\
{\scriptstyle \nu_S} \downarrow & & \downarrow {\scriptstyle \nu_{S'}} \\
\mathcal{C}_\mathrm{DOF}(S) & \xrightarrow{\mathcal{C}_\mathrm{DOF}(f)} & \mathcal{C}_\mathrm{DOF}(S')
\end{array}
\]

\textit{commutes.}

\textbf{Proof.} By construction. Full argument in §4. ∎

\sectionbreak

\section{§1.5 --- Substrate-completeness}

\textbf{Definition 1.5.1.} The \textbf{substrate-completeness condition} is the requirement that for every pair (σ, content-operation-class) consistent with the framework's adequacy, there exists a stream instantiating it. Formally:

\[
\begin{aligned}
\forall\, (σ, C) & \in \mathcal{A} \times \mathbf{ContentIndex}, \\
\mathrm{adequate}(σ, C) & \implies \exists\, S \in \mathcal{C}_\mathrm{Streams} : \pi(S) = [C] \wedge \mathrm{carrier}(S) = σ.
\end{aligned}
\]

\textbf{Remark 1.5.2.} This is the categorical content of Axiom A1 (Consciousness as Substrate): every vantage that could bear a content-operation-class does bear one. §2 (A1 chapter) gives the full formal content; §1.5 states the condition as a reference.

\textbf{Consequence 1.5.3.} Substrate-completeness is what lets the framework reason "every vantage is a stream" (Anchor §3.2). Without it, the framework would need to distinguish vantages that are streams from vantages that merely could be --- a distinction the framework refuses.

\sectionbreak

\section{§1.6 --- The endofunctor F and Stream-as-F-coalgebra}

\textbf{Definition 1.6.1 (F on carriers).} For σ in 𝒜 equipped with ContentOp(σ) ∈ \textbf{Cat}\_\allowbreak{}small:

\[
F(σ) := σ^{\mathrm{ContentOp}(σ)^\mathrm{op}}
\]

--- the presheaf-power of σ indexed contravariantly by ContentOp(σ).

\textbf{Remark 1.6.2.} Mixed-variance: σ ↦ ContentOp(σ) is covariant on carriers (content-operations pull back along carrier-maps); σ ↦ σ\textasciicircum{}(-) is covariant in σ; the combined F is well-defined as a functor on Stream (see §6.1 for the lift of F to Stream-morphisms).

\textbf{Definition 1.6.3 (Stream-morphism --- reference).} For details see Definition 6.1.3. A Stream-morphism f : S → S' consists of:
\begin{itemize}
\item Carrier-map f\_\allowbreak{}σ : σ → σ' in 𝒜
\item ContentOp-functor f\_\allowbreak{}C : ContentOp(σ) → ContentOp(σ')
\item Coalgebra-commute: γ' ∘ f\_\allowbreak{}σ = F(f\_\allowbreak{}σ, f\_\allowbreak{}C) ∘ γ
\item Kind-respect: K(S) ⊑ K(S') in the A2 preorder (Convention 1.1.5)
\end{itemize}

\textbf{Definition 1.6.4 (F-coalgebra).} An F-coalgebra is a pair (σ, γ) with γ : σ → F(σ). An F-coalgebra is \textbf{adequate} iff ContentOp(σ) is adequate for σ (Convention 1.1.6).

\textbf{Proposition 1.6.5 (Stream-categorical content).} \textit{𝒞\_\allowbreak{}Streams ≃ F-Coalg\_\allowbreak{}ad, the full subcategory of adequate F-coalgebras with kind-respecting morphisms (Theorem 6.7.1).}

\sectionbreak

\section{§1.7 --- The Triple and its target}

\subsection{§1.7.1 --- 𝒞\_\allowbreak{}Triple}

\textbf{Definition 1.7.1.} The Triple target category is

\[
\mathcal{C}_\mathrm{Triple} := \mathcal{C}_\mathrm{Form} \times \mathbf{Cat}_\mathrm{small} \times \mathbf{Carrier}
\]

where \textbf{Carrier} is the category whose objects are coalgebra-structure-maps γ : σ → F(σ) and whose morphisms are coalgebra-commute squares.

\subsection{§1.7.2 --- The Triple functor T}

\textbf{Definition 1.7.2.} The \textbf{Triple functor}

\[
T : \mathcal{C}_\mathrm{Streams} \to \mathcal{C}_\mathrm{Triple}
\]

is defined by T(σ, ContentOp(σ), γ) = (σ, ContentOp(σ), γ) --- three projections taking the three components of a Stream-object to the three factors of 𝒞\_\allowbreak{}Triple.

§6.2 gives the full definition, forgetfulness (Lemma 6.2.4), and conservativity (Proposition 6.2.7). §6.3 proves finite-depth recursive decomposability (Lemma 6.3.2).

\subsection{§1.7.3 --- The iterated Triple T\textasciicircum{}(n)}

\textbf{Definition 1.7.3.} T\textasciicircum{}(n) : 𝒞\_\allowbreak{}Streams → 𝒞\_\allowbreak{}Triple\textasciicircum{}(n) is defined inductively by T\textasciicircum{}(1) = T and T\textasciicircum{}(n+1)(S) = (T\textasciicircum{}(1)(Form(S)), T\textasciicircum{}(1)(Content(S)), T\textasciicircum{}(1)(Carrier(S))).

See §6.3 for finite-depth factorability; §6.9 for depth-ω structure.

\sectionbreak

\section{§1.8 --- Notation quick-reference}

\begingroup\setlength{\tabcolsep}{3pt}\renewcommand{\arraystretch}{1.15}\footnotesize
\begin{longtable}{p{1.17in} p{2.19in} p{0.97in}}
\toprule
Symbol & Meaning & First formal def \\
\midrule
\endhead
\bottomrule
\endfoot
𝒜 & Ambient category (Set by default) & §1.1.1 \\
\midrule
𝒞\_\allowbreak{}Streams, 𝒞\_\allowbreak{}Str, Stream & Category of streams & §1.2.1, §6.1.1 \\
\midrule
𝒞\_\allowbreak{}Form & Bare carrier-objects category & §1.2.2 \\
\midrule
𝒞\_\allowbreak{}LDS & Linked-dynamic-streams category (pairs + ι⊣κ) & §1.2.3 \\
\midrule
𝒞\_\allowbreak{}DOF & DOF-gradient-equipped streams & §1.2.4 \\
\midrule
𝒞\_\allowbreak{}Triple & Form × Cat\_\allowbreak{}small × Carrier & §1.2.5, §1.7.1 \\
\midrule
σ & Carrier object (stream-internal) & §1.6.1 \\
\midrule
ContentOp(σ) & Small category of content-operations on σ & §1.6.1 \\
\midrule
γ & Coalgebra-structure map σ →\allowbreak  F(σ) & §1.6.4 \\
\midrule
F & Endofunctor σ ↦ σ\textasciicircum{}(ContentOp(σ)\textasciicircum{}op) & §1.6.1 \\
\midrule
N & Navigation functor; N = F on streams & §1.3.1 \\
\midrule
ν & Conscious-gravity natural transformation & §1.4.1 \\
\midrule
T & Triple functor 𝒞\_\allowbreak{}Streams →\allowbreak  𝒞\_\allowbreak{}Triple & §1.7.2 \\
\midrule
T\textasciicircum{}(n) & Iterated Triple & §1.7.3 \\
\midrule
π & Kind-classifier fibration & §6.4.2 \\
\midrule
ContentIndex & Preorder of ContentOp-classes & §6.4.1 \\
\midrule
Stream\_\allowbreak{}u & Fibered subcategory of unifying streams & §6.4.12 \\
\midrule
c\_\allowbreak{}⊥ & Terminal (unifying) content-operation & §6.4.9 \\
\midrule
K & Stream kind & §1.1.5, §6.4.2 \\
\midrule
Ω & Configuration space; Ω = F(σ) & §1.6.1, Remark 6.1.2 \\
\midrule
⊑ & A2 preorder (kind-refinement) & §1.1.5 \\
\midrule
ι ⊣ κ & Cooperative-constituency adjunction & §1.2.3 \\
\midrule
Bias(S) & Signed measure on Ω(S) (coherence-weighting) & §3.2, Appendix B \\
\midrule
D(S) & Trajectory-divergence functional & §9 \\
\bottomrule
\end{longtable}
\endgroup
\sectionbreak

\section{§1.9 --- Forward-pointers}

§1's definitions feed:
\begin{itemize}
\item \textbf{§2} (Axioms): A1 uses §1.5; A2 uses §§1.2.1, 1.6; A3 uses §1.4.
\item \textbf{§3} (Theorem pairs): T1/T2 use §1.7; T3/T4 use §1.3, §1.4; T5/T6 use §1.6, §6.4.
\item \textbf{§4} (Corollary clusters): uses every §1 object.
\item \textbf{§5} (Coherence Principle): formal four-conditions stated over §1.3, §1.4, §1.6.
\item \textbf{§6} (Identity-Trajectory Triple): uses §1 throughout; §6 is the Triple-specific formalization of §1.7.
\item \textbf{§7} (Filtering): uses §1.6 (F) and §6.4 (fibration) for kind-respecting filters.
\item \textbf{§8} (F-as-stream): the self-reference closure σ\_\allowbreak{}∞ ≅ F(σ\_\allowbreak{}∞) is Theorem 6.9.1 under hypotheses.
\item \textbf{§9} (D trajectory-divergence): uses §1.3 (N-orbits) for trajectory-structure.
\end{itemize}

\sectionbreak


\chapter{§2 — Axioms A1 / A2 / A3}


\textit{Axioms stated in full CT under the F-coalgebra foundation (§6) and the category framework (§1). Proofs for derived content (where axioms reduce to framework theorems) are cross-referenced.}

\sectionbreak

\section{§2.0 --- Orientation}

§1 fixes the category framework. §6 develops the Identity-Trajectory Triple. §2 states the axioms that give the framework its content:

\begin{itemize}
\item \textbf{A1 --- Consciousness as Substrate.} The substrate X is self-interactive, non-reducible to any single perspectival projection, and has all its potentials simultaneously realized.
\item \textbf{A2 --- Nested Streams and Navigation.} Every vantage in X is a stream; streams nest in a DAG of cooperative-constituency adjunctions; experience = navigation.
\item \textbf{A3 --- Conscious Gravity.} Each stream carries a coalgebraic DOF-gradient structure that modulates its Bias(S) without reshaping X.
\end{itemize}

Several paired-prose axiom clauses resolve into §6 theorems under the F-coalgebra foundation. Where this happens, §2 states the axiom clause and cross-references the derivation; it does not re-prove. The axioms retain axiomatic status for \textbf{substrate-level claims that §6 cannot derive from F alone} (specifically A1's non-reducibility and A2's clauses that refer to X rather than to Stream).

\sectionbreak

\section{§2.1 --- Axiom 1: Consciousness as Substrate}

\subsection{§2.1.1 --- Setup}

\textbf{Definition 2.1.1 (The substrate).} Let \textbf{X} be a self-interactive process --- not an object of any single category, but the source of all perspectival projections. Formally, X is named by its projections; there is no free-standing categorical object "X" apart from the collection (F\_\allowbreak{}i)\_\allowbreak{}{i ∈ I} of perspectival functors.

\textbf{Definition 2.1.2 (Perspectival projections).} For each i ∈ I (an index class of perspectives), let F\_\allowbreak{}i : 𝒞\_\allowbreak{}P → 𝒞\_\allowbreak{}Desc\_\allowbreak{}i be a functor from the category of perspectival positions 𝒞\_\allowbreak{}P to a description category 𝒞\_\allowbreak{}Desc\_\allowbreak{}i. Canonical examples:

\begin{itemize}
\item F\_\allowbreak{}1 : 𝒞\_\allowbreak{}P → 𝒞\_\allowbreak{}Desc\_\allowbreak{}structural --- the "physical" projection.
\item F\_\allowbreak{}2 : 𝒞\_\allowbreak{}P → 𝒞\_\allowbreak{}Desc\_\allowbreak{}experiential --- the "phenomenal" projection, into the stream-category 𝒞\_\allowbreak{}Streams.
\end{itemize}

\textbf{Remark 2.1.3.} I is not required to be a set. The F\_\allowbreak{}i's form an index-class of projection-functors; the framework does not commit to a fixed list, only to the existence of at least F\_\allowbreak{}1 and F\_\allowbreak{}2 and their non-factoring (below).

\subsection{§2.1.2 --- The axiom}

\textbf{Axiom A1 (Consciousness as Substrate).} The substrate X and its projection-family (F\_\allowbreak{}i)\_\allowbreak{}{i ∈ I} satisfy:

\textbf{(A1.1) Non-reducibility.} There is no functor U and no single i ∈ I such that X ≅ U(F\_\allowbreak{}i(𝒞\_\allowbreak{}P)). X is not derivable from any single F\_\allowbreak{}i.

\textbf{(A1.2) Non-factoring.} For any two i, j ∈ I, there is no natural transformation η : F\_\allowbreak{}i ⇒ F\_\allowbreak{}j factoring as the composite of a functor with F\_\allowbreak{}i (or F\_\allowbreak{}j). The projections are parallel, not hierarchical.

\textbf{(A1.3) Configurational completeness.} The configuration space \textbf{C} (= ob(𝒞\_\allowbreak{}P) with its morphism-structure) is a complete category: every diagram has a limit. Every object is present; there is no actualization-out-of-possibility predicate on ob(𝒞\_\allowbreak{}P).

\textbf{(A1.4) Substrate-completeness (correspondence).} For every (σ, C) in 𝒜 × \textbf{ContentIndex} such that (σ, C) is adequate (Convention 1.1.6), there exists a stream S ∈ 𝒞\_\allowbreak{}Streams with carrier σ and π(S) = [C] (Definition 1.5.1).

\subsection{§2.1.3 --- Remarks and derivations}

\textbf{Remark 2.1.4.} (A1.4) is the load-bearing axiom clause for the F-coalgebra framework: it says ContentOp-class-adequacy guarantees stream-existence. Without (A1.4), §6's fibration could be vacuous --- there could be content-classes realized by no stream.

\textbf{Remark 2.1.5.} (A1.1) and (A1.2) are substrate-level claims about X and its projections. They are not reducible to §6 theorems: §6 formalizes streams (which are F\_\allowbreak{}2-projections), not the substrate. These clauses retain genuine axiomatic weight.

\textbf{Remark 2.1.6.} (A1.3) ensures the configuration space admits the limits that §6.8 uses. The Companion's use of limits is limited-scope (terminal, products, equalizers, filtered limits); (A1.3) provides these.

\textbf{Proposition 2.1.7 (Consequence of A1 for Stream).} \textit{Under A1, the category 𝒞\_\allowbreak{}Streams has a terminal object (Proposition 6.8.1), admits products conditional on kind-join (Proposition 6.8.2), and has filtered limits under accessibility (Proposition 6.8.4).}

\textbf{Proof.} Immediate from (A1.3) and §6.8. ∎

\sectionbreak

\section{§2.2 --- Axiom 2: Nested Streams and Navigation}

\subsection{§2.2.1 --- Setup}

\textbf{Recalled from §1.} 𝒞\_\allowbreak{}Streams = 𝒞\_\allowbreak{}Str (§1.2.1, Definition 6.1.1). The kind-preorder is reactive ⊑ self-maint ⊑ self-ref ⊑ abstr (§1.1.5). The navigation functor N = F (§1.3.1).

\subsection{§2.2.2 --- The axiom}

\textbf{Axiom A2 (Nested Streams and Navigation).} The framework commits to:

\textbf{(A2.1) Universal-stream.} Every F\_\allowbreak{}2-projection of X at a perspectival position p yields a stream: S\_\allowbreak{}p := F\_\allowbreak{}2(𝒞\_\allowbreak{}P, p) ∈ 𝒞\_\allowbreak{}Streams.

\textbf{(A2.2) Kind-stratification.} 𝒞\_\allowbreak{}Streams is stratified by the kind-preorder:

\[
\mathcal{C}_\mathrm{Str}^\mathrm{reactive} \supseteq \mathcal{C}_\mathrm{Str}^\mathrm{self\text{-}maint} \supseteq \mathcal{C}_\mathrm{Str}^\mathrm{self\text{-}ref} \supseteq \mathcal{C}_\mathrm{Str}^\mathrm{abstr}
\]

with strict sub-category inclusions. The kind-preorder is a fibration over ContentIndex (Theorem 6.4.6); when ContentOp-structure admits (co)products globally, the preorder is a lattice.

\textbf{(A2.3) Kinds-are-perspectival.} The kind-taxonomy (A2.2) is itself generated by the navigation of an abstracting stream s ∈ 𝒞\_\allowbreak{}Str\textasciicircum{}abstr. Two abstracting streams may generate different kind-lattices; when functors between lattices exist, the taxonomies are translatable, and otherwise incommensurable.

\textbf{(A2.4) Cooperative-constituency as adjunction.} For nested streams S\_\allowbreak{}p ⊆ S\_\allowbreak{}q (p a position within q):

\[
\iota : S_p \hookrightarrow S_q,\quad \kappa : S_q \to S_p,\quad \iota \dashv \kappa
\]

--- ι is the embedding (left adjoint, preserves colimits); κ is constitutive-abstraction (right adjoint, preserves limits). The Hom-isomorphism

\[
\mathrm{Hom}_{\mathcal{C}_\mathrm{Str}}(\iota(S_p), S_q) \cong \mathrm{Hom}_{\mathcal{C}_\mathrm{Str}}(S_p, \kappa(S_q))
\]

enforces mutual constituency.

\textbf{(A2.5) Experience = navigation.} Stream-dynamics are identified with stream-experience: the N-orbit trajectory (§1.3.1) \textit{is} the experience, not a separate observed feature. Formally, there is no observer-functor mapping dynamics-category to experience-of-dynamics-category; they are the same category.

\textbf{(A2.6) DAG nesting.} The nesting structure under ι is at minimum a DAG (directed acyclic graph): a stream may be nested in multiple non-comparable super-streams simultaneously, but no cyclic chain ι\_\allowbreak{}1 ∘ ι\_\allowbreak{}2 ∘ ... ∘ ι\_\allowbreak{}n = id.

\textbf{(A2.7) Constitutive-duality absorption.} The content of constitutive duality --- the scale-universality of ι ⊣ κ --- is contained in (A2.4); no separate theorem is required at the axiom-tier.

\subsection{§2.2.3 --- Derivation of (A2.2) from F}

\textbf{Theorem 2.2.8 (Kind-stratification derives from ContentOp-richness).} \textit{Under the F-coalgebra foundation and adequacy (Convention 1.1.6), (A2.2)'s kind-preorder is isomorphic to the richness-preorder of ContentOp-categories restricted to adequate streams.}

\textbf{Proof.} By Theorem 6.4.6, π : 𝒞\_\allowbreak{}Streams → ContentIndex is a bicategorical fibration. The A2 sub-categories 𝒞\_\allowbreak{}Str\textasciicircum{}K are preimages π\textasciicircum{}{-1}(ContentIndex\textasciicircum{}K) for each kind-class K ⊑ abstr. Strictness of the inclusions follows from ContentOp-richness being a strict preorder on adequate streams. ∎

\textbf{Remark 2.2.9.} (A2.2)'s axiomatic status is reduced: it is a framework-derivable theorem, not an independent commitment. The framework's \textit{axiomatic} content at A2 is (A2.1), (A2.3), (A2.4), (A2.5), (A2.6), (A2.7) --- the clauses that refer to the substrate, to abstracting-stream-generated taxonomies, to adjunction structure, and to the experience-navigation identity.

\subsection{§2.2.4 --- Coupling morphisms via 𝒞\_\allowbreak{}LDS}

\textbf{Definition 2.2.10.} For paired streams S\_\allowbreak{}1, S\_\allowbreak{}2 linked by ι ⊣ κ (A2.4), the pair (S\_\allowbreak{}1, S\_\allowbreak{}2) ∈ 𝒞\_\allowbreak{}LDS carries a coupling-morphism structure: a morphism (S\_\allowbreak{}1, S\_\allowbreak{}2) → (S'\_\allowbreak{}1, S'\_\allowbreak{}2) is a Stream-pair morphism respecting the adjunction. §3.2 (A2 coupling clause in T3/T4) uses this structure for the paired-dyad theorems.

\textbf{Proposition 2.2.11 (Constitutive duality is adjunction-universality).} \textit{Constitutive duality --- the scale-universality of ι ⊣ κ --- is the statement that (A2.4)'s Hom-isomorphism holds for every nested-stream pair at every scale.}

\textbf{Proof.} Re-statement in different language. ∎

\sectionbreak

\section{§2.3 --- Axiom 3: Conscious Gravity}

\subsection{§2.3.1 --- Setup}

A3 treats the DOF-gradient structure ν (§1.4.1) that modulates Bias(S). This is coalgebraic: γ\_\allowbreak{}S (the Stream's coherence-coalgebra, §1.6.4) carries Bias-data, and A3 describes how.

\subsection{§2.3.2 --- The axiom}

\textbf{Axiom A3 (Conscious Gravity).} For every stream S = (σ, ContentOp(σ), γ) ∈ 𝒞\_\allowbreak{}Streams:

\textbf{(A3.1) Coalgebraic gravity structure.} There is a coalgebra γ\_\allowbreak{}S : S → Bias(S) × S, where:

\begin{itemize}
\item \textbf{Bias(S)} is a signed measure on the configuration space Ω(S) = F(σ) encoding S's path-weighting preferences.
\item The product structure means γ\_\allowbreak{}S updates both the Bias and the state-within-Bias at each navigation step.
\end{itemize}

\textbf{(A3.2) Internality.} γ\_\allowbreak{}S acts only on S's F\_\allowbreak{}2-internal structure, not on X. Formally: there is no functor δ with codomain 𝒞\_\allowbreak{}P such that γ\_\allowbreak{}S factors through δ. Conscious gravity reshapes S's weighting of paths within its own F\_\allowbreak{}2-projection; it does not reshape the substrate.

\textbf{(A3.3) DOF-gradient modulation.} γ\_\allowbreak{}S modulates Bias(S) along a continuous degrees-of-freedom gradient for coherence:

\[
\nu_S : \mathcal{C}_\mathrm{Streams} \to \mathcal{C}_\mathrm{DOF},\quad \nu_S(S) = (\mathrm{DOF\text{-}rank}(S),\ \mathrm{slope\text{-}measure}(S))
\]

The DOF-rank is an ordinal derived from ContentOp-richness (§6.4.2); the slope-measure is the local Bias-gradient.

\textbf{(A3.4) Adaptivity.} γ\_\allowbreak{}S is itself updated by stream-navigation: operating γ\_\allowbreak{}S on state σ\_\allowbreak{}t produces (Bias\_\allowbreak{}{t+1}, σ\_\allowbreak{}{t+1}) with Bias\_\allowbreak{}{t+1} possibly differing from Bias\_\allowbreak{}t. The coalgebra-structure is time-varying by design.

\textbf{(A3.5) Stream-universality.} A3 holds for every S ∈ 𝒞\_\allowbreak{}Streams; there is no stream without a Bias(S) structure.

\subsection{§2.3.3 --- Remarks and derivations}

\textbf{Remark 2.3.4.} (A3.3) uses a continuous-DOF-gradient form; the three named regimes (attention / intention / belief) appear as region-structure on the continuous DOF-axis, not as independent axes.

\textbf{Remark 2.3.5.} (A3.2)'s non-factoring is the formal content of "conscious gravity does not reshape X." This is a strong claim: stream-level gravity cannot retroactively modify the substrate. It is reducible to §2.1.2's (A1.1)+(A1.2) under the specific check that γ\_\allowbreak{}S's codomain is σ-internal.

\textbf{Proposition 2.3.6 (Adaptivity is encoded in F).} \textit{Under the F-coalgebra foundation, (A3.4)'s adaptivity is the statement that F is a proper endofunctor with γ : σ → F(σ) iteratively applied (not a one-shot map). This holds in §6 by construction.}

\textbf{Proof.} By §6.1.5 F is an endofunctor on 𝒞\_\allowbreak{}Streams; γ\_\allowbreak{}S iterates as N (Definition 1.3.1). ∎

\textbf{Remark 2.3.7.} (A3.5)'s stream-universality matches Proposition 6.1.4's observation that every Stream-object carries γ-data by Definition 6.1.1. A3's coverage is built into Stream-hood.

\subsection{§2.3.4 --- Bias(S) as a signed measure}

\textbf{Definition 2.3.8.} For each S ∈ 𝒞\_\allowbreak{}Streams, Bias(S) is a signed measure on Ω(S) = σ\textasciicircum{}(ContentOp(σ)\textasciicircum{}op) defined on the σ-algebra generated by the content-operation-partition (§7 gives the formal σ-algebra construction). Bias(S) decomposes as:

\[
\mathrm{Bias}(S) = \mathrm{Bias}^+(S) - \mathrm{Bias}^-(S)
\]

with Bias\textasciicircum{}+ supported on coherence-attractors and Bias\textasciicircum{}- on coherence-repellors. The entropy functional A\_\allowbreak{}S (Anchor §6.3 / Companion Appendix B) computes local Bias-weighting.

\textbf{Proposition 2.3.9 (Bias(S) well-definedness).} \textit{Under the σ-algebra construction of §7 and the F-coalgebra structure of §6, Bias(S) is a well-defined signed measure on Ω(S) for every S.}

\textbf{Proof.} Deferred to §7. ∎

\sectionbreak

\section{§2.4 --- Axiom-interaction: A1 + A2 + A3 jointly}

\textbf{Proposition 2.4.1 (A1 grounds A2).} \textit{Under A1.4 (substrate-completeness), A2.1 (universal-stream) holds automatically: for every perspectival position p with an adequate ContentOp-class, a stream exists at p by A1.4.}

\textbf{Proof.} A1.4 gives stream-existence for every adequate (σ, C) pair; A2.1 requires stream-existence for every F\_\allowbreak{}2-projection position. Since F\_\allowbreak{}2-projections yield adequate pairs (by construction --- F\_\allowbreak{}2 is an experiential projection, which by A1.2's non-factoring carries all projection-data including ContentOp-adequacy), the implication holds. ∎

\textbf{Proposition 2.4.2 (A2 + F-coalgebra jointly imply kind-coupling via Content).} \textit{A2's kind structure and stream-Content couple via the coalgebra-commute condition (Definition 6.1.3 (iii)) rather than as a conjunction; this follows from (A2.2) interpreted as a fibration over Content (Theorem 6.4.6) and F-coalgebra morphism structure.}

\textbf{Proof.} Coalgebra-commute links carrier-map f\_\allowbreak{}σ, ContentOp-functor f\_\allowbreak{}C, and kind-respect in a single condition; Theorem 6.4.6 routes this as a fibration over Content. ∎

\textbf{Proposition 2.4.3 (A3 internality is compatible with F-iteration).} \textit{(A3.2) internality and (A3.4) adaptivity are jointly consistent with §6's F-coalgebra iteration: γ\_\allowbreak{}S is updated within S's F\_\allowbreak{}2-internal structure at each navigation step, without reaching back into X.}

\textbf{Proof.} F's codomain is 𝒞\_\allowbreak{}Streams (§6.1.5), not 𝒞\_\allowbreak{}P directly. Iteration stays within stream-categories. ∎

\sectionbreak

\section{§2.5 --- Axiom-status summary}

\begingroup\setlength{\tabcolsep}{3pt}\renewcommand{\arraystretch}{1.15}\footnotesize
\begin{longtable}{p{0.55in} p{1.58in} p{2.20in}}
\toprule
Axiom & Clause & Status under F-coalgebra foundation \\
\midrule
\endhead
\bottomrule
\endfoot
A1 & A1.1 non-reducibility & Axiomatic (substrate-level) \\
\midrule
A1 & A1.2 non-factoring & Axiomatic (substrate-level) \\
\midrule
A1 & A1.3 configurational completeness & Axiomatic (limits premise) \\
\midrule
A1 & A1.4 substrate-completeness & Axiomatic (load-bearing for §6) \\
\midrule
A2 & A2.1 universal-stream & Derivable from A1.4 (Prop 2.4.1) \\
\midrule
A2 & A2.2 kind-stratification & Derived from F (Theorem 2.2.8) \\
\midrule
A2 & A2.3 kinds-are-perspectival & Axiomatic (reflects abstracting-stream generation) \\
\midrule
A2 & A2.4 cooperative-constituency & Axiomatic (adjunction-universality) \\
\midrule
A2 & A2.5 experience = navigation & Axiomatic (category-identity claim) \\
\midrule
A2 & A2.6 DAG nesting & Axiomatic (non-cyclicity) \\
\midrule
A2 & A2.7 constitutive-duality absorption & Meta-statement (absorbed into A2.4) \\
\midrule
A3 & A3.1 coalgebra structure & Derivable from F-coalgebra definition \\
\midrule
A3 & A3.2 internality & Axiomatic (substrate-protection) \\
\midrule
A3 & A3.3 DOF-gradient & Axiomatic \\
\midrule
A3 & A3.4 adaptivity & Derivable (Prop 2.3.6) \\
\midrule
A3 & A3.5 stream-universality & Derivable (built into Stream-hood) \\
\bottomrule
\end{longtable}
\endgroup
Of 16 clauses across three axioms: 10 retain axiomatic status; 6 are framework-derivable. The axiomatic economy gained by the F-coalgebra foundation concentrates on A1 (substrate-level non-reducibility and completeness) and the identity-claims of A2/A3 (experience-navigation, internality).

\sectionbreak

\section{§2.6 --- Forward-pointers}

\begin{itemize}
\item \textbf{§3} (Theorem pairs): T1/T2 descriptive pair uses A1 (substrate) + A2 (stream). T3/T4 dynamics pair uses A3 (Bias + conscious gravity) + A2 (coupling). T5/T6 coherence pair uses all three.
\item \textbf{§4} (Corollary clusters): clusters organized by axiom-lineage. §8 Anchor numbering carries across.
\item \textbf{§5} (Coherence Principle): four-conditions statement combines axiom-content across A1+A2+A3.
\item \textbf{§6} (Triple): the Triple functor and its properties (already drafted) are the structural content of stream-internal identity under A2+A3.
\item \textbf{§7} (Filtering): σ-algebra on Ω\_\allowbreak{}S (A3.1 / 2.3.8) and Bias(S) well-definedness (Proposition 2.3.9).
\item \textbf{§9} (D trajectory-divergence): Bias(S)-derived functional on N-orbits (A3.3).
\end{itemize}

\sectionbreak

\section{Surfaced-lemma register (§2)}

\needspace{10\baselineskip}
\begin{Code}
⚑ [SURFACED | Companion §2.2.8 | → Anchor §3 target | type: theorem]
  — Kind-stratification derives from ContentOp-richness (reduces A2.2 to framework theorem).

⚑ [SURFACED | Companion §2.4.1 | → Anchor §3 target | type: proposition]
  — A1.4 grounds A2.1 (substrate-completeness implies universal-stream).

⚑ [SURFACED | Companion §2.5 | → Anchor §2-§4 target | type: axiom-status classification]
  — 10-of-16 clauses retain axiomatic status; 6 reduce to framework theorems.
\end{Code}

\sectionbreak


\chapter{§3 — Theorems}


\textit{Three pairs, six theorems. Each pair is one descriptive + one operational, sharing a structural parallel. Proofs are complete at paper density. Citations resolve to §1 (framework), §2 (axioms), §6 (Triple). Anchor-side exposition lives in Anchor §§5–7.}

\sectionbreak

\section{§3.0 --- Orientation}

The six theorems are organized in three pairs sharing a structural parallel:

\begingroup\setlength{\tabcolsep}{3pt}\renewcommand{\arraystretch}{1.15}\footnotesize
\begin{center}
\begin{tabular}{p{0.60in} p{2.01in} p{1.04in} p{0.63in}}
\toprule
Pair & Theorem & Register & Anchor location \\
\midrule
I. Descriptive & T1 (Mathematical Perspectivism) & static /\allowbreak  representational & Anchor §5 \\
 & T2 (Estimator-Dependent Duration) & static /\allowbreak  temporal & Anchor §5 \\
II. Dynamics & T3 (Attentional Quality \& Navigational Dynamics) & dynamic /\allowbreak  functorial & Anchor §6 \\
 & T4 (Coherence-Forcing Measurement) & dynamic /\allowbreak  operational & Anchor §6 \\
III. Coherence & T5 (Internal Coherence) & internal /\allowbreak  fixed-point & Anchor §7 \\
 & T6 (Dual Coherence Axes) & structural /\allowbreak  entropic & Anchor §7 \\
\bottomrule
\end{tabular}
\end{center}
\endgroup
The shared parallel within each pair is a \textbf{Descriptive-Functor Meta-Theorem pattern} (§3.1.α below): each first-in-pair theorem states \textit{what} the framework sees, each second-in-pair states \textit{how} measurement or temporal-estimation operates on that view. The three pairs together discharge the proof-burden of the Coherence Principle (§5).

\textbf{Notation.} All symbols carry the meanings fixed in §1.8. We write \textit{S = (σ, C, Ω, γ)} for a stream with carrier σ, content-operation category C := ContentOp(σ), configuration space Ω := F(σ) = σ\textasciicircum{}(C\textasciicircum{}op), and coalgebra γ : σ → Ω. The Triple-functor is T : 𝒞\_\allowbreak{}Streams → 𝒞\_\allowbreak{}Triple (Definition 1.7.2).

\sectionbreak

\section{§3.1 --- Descriptive-Functor Meta-Theorem}

\textbf{Theorem 3.1.α (Descriptive-Functor Meta-Theorem).} \textit{Let P be a property of streams expressible as a section of a stream-indexed presheaf}

\[
\mathcal{P} : \mathcal{C}_\mathrm{Streams}^\mathrm{op} \to \mathbf{Set}.
\]

\textit{Then P is descriptive in the Corpus sense iff there is a functor}

\[
\mathsf{Desc}_P : \mathcal{C}_\mathrm{Streams} \to \mathbf{Set}
\]

\textit{and a natural isomorphism}

\[
\mathcal{P}(S) \cong \mathsf{Desc}_P(S)
\]

\textit{rendering P representable in Stream --- i.e., P factors through the Yoneda embedding.}

\textbf{Proof.} (⇒) Assume P is descriptive. Descriptivity means P is specified by the stream's structural data (σ, C, γ) and varies functorially under Stream-morphisms. Any such specification is a Stream-morphism-respecting section; by the Yoneda lemma it is represented by the functor 𝖣𝖾𝗌𝖼\_\allowbreak{}P(S) := Nat(h\_\allowbreak{}S, 𝒫), which is isomorphic to 𝒫(S) naturally in S.

(⇐) If 𝒫(S) ≅ 𝖣𝖾𝗌𝖼\_\allowbreak{}P(S) naturally, then for each Stream-morphism f : S → S' the induced map 𝒫(f) : 𝒫(S') → 𝒫(S) is determined by 𝖣𝖾𝗌𝖼\_\allowbreak{}P(f). Since 𝖣𝖾𝗌𝖼\_\allowbreak{}P is a functor out of 𝒞\_\allowbreak{}Streams, P depends only on the structural data --- precisely the descriptivity condition. ∎

\textbf{Remark 3.1.β.} T1 and T2 below are the two canonical instances: T1 takes P to be the framework's choice of mathematical representation of a given configuration, T2 takes P to be a duration estimator. In both cases the representing functor captures what "looking at the stream" amounts to. The meta-theorem is what makes the two theorems parallel.

⚑ [SURFACED | Companion §3.1 | → Anchor §5 target | type: meta-theorem]

\sectionbreak

\section{§3.2 --- Theorem Pair I: Descriptive}

\subsection{§3.2.1 --- T1 Mathematical Perspectivism}

\textbf{Theorem 3.2.1 (T1 --- Mathematical Perspectivism).} \textit{Let S = (σ, C, Ω, γ) be a stream. The representation of S via any mathematical object M is a functor}

\[
\mathsf{Rep}_M : \mathcal{C}_\mathrm{Streams} \to \mathcal{D}_M
\]

\textit{where 𝒟\_\allowbreak{}M is the category of M-structured objects. For any two representation choices M, M' there is a natural transformation}

\[
\alpha_{M,M'} : \mathsf{Rep}_M \Rightarrow \mathsf{Rep}_{M'}
\]

\textit{witnessing the change-of-perspective. 𝖱𝖾𝗉\_\allowbreak{}M is never canonical --- it depends on C. Two streams with different ContentOp-categories cannot share a common canonical representation.}

\textbf{Proof.} The representation 𝖱𝖾𝗉\_\allowbreak{}M of S factors through (σ, C, γ); C indexes the observables-modulo-content-operation-equivalence. Since F(σ) = σ\textasciicircum{}(C\textasciicircum{}op), a different C′ produces a different Ω-structure and hence a different 𝖱𝖾𝗉\_\allowbreak{}M′. The existence of α\_\allowbreak{}{M,M'} follows from functoriality in M via the universal property of representable functors (§3.1). Non-canonicity: if 𝖱𝖾𝗉\_\allowbreak{}M were canonical across C-variation it would factor through the forgetful 𝒞\_\allowbreak{}Streams → 𝒜, losing C; but C is load-bearing for F(σ), contradicting F's definition. ∎

\textbf{Corollary 3.2.1.1 (Representation-via-Content).} \textit{Every mathematical representation of a stream carries its ContentOp structure. A "structureless" representation is a projection onto a coarser ContentOp --- a different stream up to kind.}

\textbf{Corollary 3.2.1.2 (C1 backward-compatibility).} \textit{T1 is the categorical form of what Anchor calls C1 (C1: every measurement carries a vantage).}

\subsection{§3.2.2 --- T2 Estimator-Dependent Duration}

\textbf{Theorem 3.2.2 (T2 --- Estimator-Dependent Duration).} \textit{Let S be a stream and let Δ : 𝒞\_\allowbreak{}Streams → [0, ∞] be a duration-estimator functor. Δ factors through a choice of N-orbit parametrization}

\[
\mathcal{C}_\mathrm{Streams} \xrightarrow{\mathrm{Orb}} \mathbf{Preord} \xrightarrow{\|\cdot\|_\Delta} [0, \infty]
\]

\textit{where Orb sends S to its N-orbit (as a preordered set under γ-iteration) and ‖·‖\_\allowbreak{}Δ is an order-respecting metric. Different choices of ‖·‖\_\allowbreak{}Δ give different Δ-values; no canonical choice exists.}

\textbf{Proof.} Duration is a functional on the N-orbit of S (Remark 1.3.2). Orb sends S to the preorder (Ω\_\allowbreak{}S, ≤\_\allowbreak{}γ) generated by γ-iteration. Any order-respecting [0, ∞]-valued metric ‖·‖\_\allowbreak{}Δ on this preorder yields a candidate duration-estimator; the composite is functorial because both Orb and ‖·‖\_\allowbreak{}Δ are. For non-canonicity: if two metrics ‖·‖\_\allowbreak{}Δ, ‖·‖\_\allowbreak{}Δ' agree on every stream, they agree on Orb(S) ≅ (ℕ, ≤) for a discrete stream --- forcing equality on ℕ-valued parametrizations; but this equality does not extend to streams with non-discrete N-orbits (e.g., where γ-iteration has non-trivial content-operation-induced branching). ∎

\textbf{Corollary 3.2.2.1 (C13 backward-compatibility).} \textit{T2 is the categorical form of what Anchor calls C13 (duration is estimator-relative).}

\textbf{Remark 3.2.3 (Structural parallel T1/T2).} Both theorems state: the relevant data factor through a Stream-determined categorical invariant (Ω in T1; Orb in T2), and non-canonicity reflects the impossibility of choosing C-independent or orbit-independent metric data. Theorem 3.1.α is the common core.

\sectionbreak

\section{§3.3 --- Theorem Pair II: Dynamics}

\subsection{§3.3.1 --- T3 Attentional Quality \& Navigational Dynamics}

\textbf{Theorem 3.3.1 (T3 --- Attentional Quality \& Navigational Dynamics).} \textit{Let S be a stream. The attentional-quality functional}

\[
\mathcal{Q} : \mathcal{C}_\mathrm{Streams} \to \mathrm{Meas}(\Omega)
\]

\textit{sending S to a measure on Ω\_\allowbreak{}S is natural in Stream-morphisms and decomposes into two coupling channels plus an entropy-based contracted-open axis:}

\[
\mathcal{Q}(S) = \mathcal{Q}_\mathrm{ent}(S) + \lambda_1 \cdot \mathcal{Q}_\mathrm{co-stream}(S) + \lambda_2 \cdot \mathcal{Q}_\mathrm{kind}(S)
\]

\textit{where:}
\begin{itemize}
\item \textit{𝒬\_\allowbreak{}ent is the entropy contribution H(γ ; C) of γ against C's content-operation structure,}
\item \textit{𝒬\_\allowbreak{}co-stream couples S to co-nested streams via ι ⊣ κ (§1.2.3),}
\item \textit{𝒬\_\allowbreak{}kind couples S to its kind-position via π (§6.4.2 fibration),}
\item \textit{λ\_\allowbreak{}1, λ\_\allowbreak{}2 are context-dependent coupling weights determined by Bias(S).}
\end{itemize}

\textbf{Proof.} 𝒬 is specified to be Stream-morphism-respecting; by Theorem 3.1.α it factors through a functor on 𝒞\_\allowbreak{}Streams. Decomposition: an entropy functional on γ : σ → F(σ) is well-defined because F(σ) = σ\textasciicircum{}(C\textasciicircum{}op) carries the C-indexed counting measure; H(γ ; C) = −∑\_\allowbreak{}c∈C log μ\_\allowbreak{}c(γ(-)). The two coupling channels arise from the two structures \textit{external} to S in Stream: the ι ⊣ κ adjunction embedding S into co-streams, and π(S) placing S in the kind-classifier fibration. Additivity is by linearity of entropy-contribution-plus-coupling; weights λ\_\allowbreak{}1, λ\_\allowbreak{}2 come from Bias(S)'s push-operators (§2.3.4) which quantify inward vs. outward coherence-attraction. Fullness of the decomposition: any measure 𝒬(S)-valued functional that respects Stream-morphisms must factor through these three channels --- this is the content of §2.4.2 (A2+F-imply-kind-coupling-via-Content). ∎

\textbf{Remark 3.3.1.1.} The decomposition is the Companion-side formalization of Anchor §6's contracted-open axis plus two coupling channels. Full Bias(S) apparatus lives in Appendix B (anchor) / §7 (Companion).

⚑ [SURFACED | Companion §3.3.1 | → Anchor §6 target | type: lemma]

\subsection{§3.3.2 --- T4 Coherence-Forcing Measurement}

\textbf{Theorem 3.3.2 (T4 --- Coherence-Forcing Measurement).} \textit{Let S be a stream and let M be a measurement operator --- a Stream-morphism}

\[
M : S \to S_M
\]

\textit{where S\_\allowbreak{}M carries a strictly richer ContentOp (i.e., π(S) ⊏ π(S\_\allowbreak{}M) in the kind-preorder). Then:}

\begin{enumerate}
\item \textit{M induces a post-measurement coalgebra γ\_\allowbreak{}M : σ → F\_\allowbreak{}M(σ) with F\_\allowbreak{}M(σ) = σ\textasciicircum{}(C\_\allowbreak{}M\textasciicircum{}op), and γ\_\allowbreak{}M ∘ M = F(M) ∘ γ (coalgebra-commute).}
\item \textit{The Bias-push under M satisfies $M_*(\mathrm{Bias}(S)) \geq \mathrm{Bias}(S_M)|_{M(\sigma)}$, with equality iff M is kind-preserving.}
\item \textit{If M is informed (provides the full state of γ prior to measurement, with Hamiltonian-level control --- the García-Pintos reversibility condition), then M is invertible in F-Coalg\_\allowbreak{}ad.}
\item \textit{If M is uninformed (the observer's information budget is asymptotically sufficient --- the Watanabe-Takagi regime), then M is reversible in the asymptotic-entropy sense.}
\end{enumerate}

\textbf{Proof.}
(1) By the kind-respect clause of Definition 1.6.3, M's ContentOp-functor f\_\allowbreak{}C : C → C\_\allowbreak{}M pulls γ to γ\_\allowbreak{}M = (F(M) ∘ γ)(M\_\allowbreak{}σ\textasciicircum{}{-1}(-)), which commutes by construction.

(2) The Bias-push is the pushforward of the signed measure Bias(S) under M\_\allowbreak{}σ. Kind-enrichment contracts some coherence-weight toward attractors of C\_\allowbreak{}M; kind-preservation is the equality case by §2.3.4's push-operator independence.

(3) Full-information condition: if the observer knows (σ, C, γ) exactly and controls the Hamiltonian realizing M, then the inverse M\textasciicircum{}{-1} exists in F-Coalg\_\allowbreak{}ad because (a) M\_\allowbreak{}σ is bijective (full-information requirement), (b) f\_\allowbreak{}C is an equivalence (kind-preserving under full information-plus-control), and (c) coalgebra-commute inverts. This is the 2026 García-Pintos reversibility result translated into F-Coalg\_\allowbreak{}ad.

(4) Asymptotic reversibility: for a sequence M\_\allowbreak{}n of increasingly-informed measurements, the entropy-loss H(γ ; C) − H(γ\_\allowbreak{}{M\_\allowbreak{}n} ; C\_\allowbreak{}{M\_\allowbreak{}n}) → 0 provided the information budget grows super-polynomially relative to the ContentOp-cardinality of C\_\allowbreak{}M. This is the 2026 Watanabe-Takagi work-extraction result translated into F-Coalg\_\allowbreak{}ad. ∎

\textbf{Remark 3.3.2.1 (Measurement as information-conservative operation).} Clauses (3) and (4) update the Corpus's earlier framing of measurement-as-terminal-collapse. Under the F-coalgebra formalism, measurement is a structural operation whose practical irreversibility is an \textbf{information-budget fact}, not a physical-law arrow. This reframe is carried in Anchor §9 / §9.5.

⚑ [SURFACED | Companion §3.3.2 | → Anchor §6 / §9.5 target | type: theorem]

\textbf{Remark 3.3.3 (Structural parallel T3/T4).} T3 describes the \textit{standing-state} of navigational dynamics as a decomposable measure; T4 describes the \textit{transition} under measurement as a kind-enriching morphism. Together they constitute the dynamics pair: \textit{what the stream is doing} (T3) + \textit{what changes when you look} (T4).

\sectionbreak

\section{§3.4 --- Theorem Pair III: Coherence}

\subsection{§3.4.1 --- T5 Internal Coherence}

\textbf{Theorem 3.4.1 (T5 --- Internal Coherence).} \textit{Let S = (σ, C, Ω, γ) be a stream. Internal coherence is the condition:}

\[
\begin{aligned}
& \gamma : \sigma \to F(\sigma) \text{ is a fixed-point of the coherence-self-map} \\
& \Phi_S : F(\sigma) \to F(\sigma).
\end{aligned}
\]

\textit{where Φ\_\allowbreak{}S is defined as the C-average of the forward-N-step image of γ. Concretely:}

\[
\Phi_S(\omega)(c) := \frac{1}{|C_\omega|} \sum_{c' \in C_\omega} \gamma(\omega)(c' \circ c)
\]

\textit{for ω ∈ Ω, c ∈ C, and C\_\allowbreak{}ω ⊆ C the sub-category of content-operations acting nontrivially at ω. A stream S is \textbf{internally coherent} iff Φ\_\allowbreak{}S(γ) = γ. Equivalently, γ is a C-harmonic section of F(σ).}

\textbf{Proof.} Φ\_\allowbreak{}S is well-defined because C\_\allowbreak{}ω is a small sub-category (Convention 1.1.1's small-size condition); the C-average is a finite or ω-limit expression. Functoriality of Φ\_\allowbreak{}S in Stream-morphisms: a Stream-morphism f : S → S' induces f\_\allowbreak{}C : C → C' preserving C\_\allowbreak{}ω structure (by Definition 1.6.3 coalgebra-commute), so Φ\_\allowbreak{}{S'} ∘ f\_\allowbreak{}σ = f\_\allowbreak{}σ ∘ Φ\_\allowbreak{}S. The fixed-point equation Φ\_\allowbreak{}S(γ) = γ is equivalent to the C-harmonic condition by a discrete-harmonic-function-on-C argument: γ is harmonic iff its value at each ω equals the C-weighted average of its forward images. ∎

\textbf{Corollary 3.4.1.1 (T5 under recursive decomposition).} \textit{If S is internally coherent, then each component of its Triple T(S) = (Form(S), Content(S), Carrier(S)) is also coherent in its own sub-stream sense. Conversely, the Triple-components being coherent does not imply S is coherent --- adequacy of the join is additional (§6.3, §6.8).}

\textbf{Corollary 3.4.1.2 (Coherence is a closed condition in F-Coalg\_\allowbreak{}ad).} \textit{The full sub-category of internally-coherent streams is closed under limits, co-limits that preserve adequacy, and Stream-morphism-pullbacks.}

\textbf{Proof of 3.4.1.2.} Φ is natural in S, so Φ-fixed-points form a full sub-category. Limits and adequacy-preserving colimits commute with Φ (Lemma 6.8.β). Pullbacks: for f : S → S' with S' internally coherent and f reflecting Φ-structure, f\textasciicircum{}*(γ') is Φ\_\allowbreak{}S-fixed. ∎

\subsection{§3.4.2 --- T6 Dual Coherence Axes}

\textbf{Theorem 3.4.2 (T6 --- Dual Coherence Axes).} \textit{There exist two orthogonal axes along which a stream's coherence can be evaluated:}

\[
\sigma_\mathrm{struct} \perp \sigma_\mathrm{info}
\]

\textit{where:}
\begin{itemize}
\item \textit{σ\_\allowbreak{}struct is the \textbf{structural-coherence} axis: the degree to which γ is a fixed-point of Φ\_\allowbreak{}S (T5-scalar, ∈ [0,1]).}
\item \textit{σ\_\allowbreak{}info is the \textbf{informational-coherence} axis: the degree to which the entropy functional H(γ ; C) attains its C-conditioned minimum (∈ [0,1]).}
\end{itemize}

\textit{The two axes are orthogonal in the sense that:}

\begin{enumerate}
\item \textit{σ\_\allowbreak{}struct-maximizing streams need not minimize H(γ ; C) (a highly-symmetric γ can be harmonic without being entropy-minimal).}
\item \textit{H-minimizing streams need not be Φ-fixed (a sharp γ concentrated at one ω is entropy-low but may fail C-harmonic).}
\item \textit{A stream is \textbf{dually coherent} iff both are maximized.}
\end{enumerate}

\textit{The dual-coherence locus is a codimension-2 sub-scheme in the (σ\_\allowbreak{}struct, σ\_\allowbreak{}info)-plane --- generically empty for a randomly-chosen C; non-empty when C has enough symmetry to admit both a harmonic γ and a concentrated γ.}

\textbf{Proof.}
(1) Counterexample: let σ = ℤ/2, C = ℤ/2 acting by swap. The uniform γ(0) = γ(1) = 1/2 is Φ-fixed (C-average of swap-image = uniform) but has maximum entropy H(γ ; C) = log 2.

(2) Counterexample: γ(0) = 1, γ(1) = 0 on the same (σ, C) has H(γ ; C) = 0 but Φ\_\allowbreak{}S(γ)(0) = 0 ≠ 1 = γ(0), so not harmonic.

(3) Dually-coherent locus: from (1) and (2), the two axes give distinct optima generically. The intersection is non-empty only when C admits a γ that is both harmonic and supported on a single C-orbit; this is a codimension-2 condition on the pair (σ, C). ∎

\textbf{Corollary 3.4.2.1 (Kind-demotion dynamic).} \textit{Along the σ\_\allowbreak{}info axis, decreases in H(γ ; C) can correspond to C being replaced by a coarser C′ --- i.e., kind-demotion. Conversely, along the σ\_\allowbreak{}struct axis, decreases in Φ-distance can correspond to C being enriched. The full (σ\_\allowbreak{}struct, σ\_\allowbreak{}info)-trajectory of a stream is a path in the fibration π (§6.4) crossed with the T5-harmonicity functional.}

\textbf{Remark 3.4.2.2 (Structural parallel T5/T6).} T5 states \textit{when} coherence holds (fixed-point condition); T6 states \textit{how} coherence can be high in one axis while low in another (orthogonality). Together they formalize Anchor §7's "internal coherence plus dual axes" framing.

⚑ [SURFACED | Companion §3.4.2 | → Anchor §7 target | type: theorem]

\sectionbreak

\section{§3.5 --- Theorem-to-Axiom crosswalk}

\begingroup\setlength{\tabcolsep}{3pt}\renewcommand{\arraystretch}{1.15}\footnotesize
\begin{center}
\begin{tabular}{p{0.98in} p{1.65in} p{1.70in}}
\toprule
Theorem & Primary axiom-clause & Derivation route \\
\midrule
T1 (Mathematical Perspectivism) & A1.2 (non-factoring) + A2.1 (stream = F-coalgebra) & Via Theorem 3.1.α + Prop 6.1.4 \\
T2 (Estimator-Dependent Duration) & A2.3 (experience = navigation) & Via Remark 1.3.2 (N-orbit) + Theorem 3.1.α \\
T3 (Attentional Quality) & A2 + A3.1 (DOF-gradient) & Via §2.4.2 (A2+F imply kind-coupling) + Def 1.4.1 (ν) \\
T4 (Coherence-Forcing Measurement) & A3.2 (internality) + A2.4 (kind-refinement) & Via Prop 2.4.3 (A3-internality compatible with F-iteration) \\
T5 (Internal Coherence) & A3.3 (coherence-attraction) & Via coalgebra-fixed-point construction \\
T6 (Dual Coherence Axes) & A3.4 (DOF-gradient structure) + A2.7 (C-richness lattice) & Via §6.4.6 (fibration) + entropy functional \\
\bottomrule
\end{tabular}
\end{center}
\endgroup
\sectionbreak

\section{§3.6 --- The three pairs and the Principle}

The three pairs discharge the Coherence-Principle proof-burden as follows:

\begin{itemize}
\item \textbf{Descriptive pair (T1/T2)} establishes \textit{condition 1} of the Principle (§5.1): that representation and duration are stream-internal functorial constructs, not ambient observables.
\item \textbf{Dynamics pair (T3/T4)} establishes \textit{conditions 2 and 3} (§5.2, §5.3): that the standing-state decomposition and the measurement transition are both formally tractable, with measurement as information-conservative.
\item \textbf{Coherence pair (T5/T6)} establishes \textit{condition 4} (§5.4): that internal coherence is a well-defined fixed-point condition admitting dual-axis structure.
\end{itemize}

§5 collects these into the Principle's formal statement.

\sectionbreak

\section{§3.7 --- Forward-pointers}

\begin{itemize}
\item \textbf{§4} (Corollary clusters): C1–C13 derive from the six theorems + §2 axioms.
\item \textbf{§5} (Coherence Principle): the four conditions are the T1/T2, T3, T4, T5/T6 content reassembled as an operational principle.
\item \textbf{§6.4} (kind-classifier fibration): T3's kind-coupling channel uses π directly.
\item \textbf{§7} (Filtering): Bias(S) apparatus underwriting T3 λ\_\allowbreak{}1, λ\_\allowbreak{}2 and T4 clause 2.
\item \textbf{§8} (F-as-stream): T4's self-application (measurement as Stream-morphism applied to F itself) is the self-reference closure.
\item \textbf{§9} (D trajectory-divergence): T5-distance and T6-distance together parametrize D.
\end{itemize}

\section{§3.8 --- Surfaced-lemma register}

Three flags surface in §3 this drafting pass:

\begin{itemize}
\item ⚑ §3.1 Descriptive-Functor Meta-Theorem → Anchor §5 target --- meta-theorem (new framing)
\item ⚑ §3.3.1 Attentional-Quality three-channel decomposition → Anchor §6 target --- lemma (explicit decomposition form)
\item ⚑ §3.3.2 Measurement-as-information-conservative (Watanabe-Takagi + García-Pintos integration) → Anchor §6/§9.5 target --- theorem (reframe of earlier collapse-language)
\item ⚑ §3.4.2 Dual-axes orthogonality-by-counterexample → Anchor §7 target --- theorem (explicit codimension-2 locus)
\end{itemize}

Back-port targets feed future Anchor revisions per the Companion/Anchor rhythm specified in SCOPE §8.

\sectionbreak


\chapter{§4 — Corollary clusters}


\textit{Thirteen corollaries in three clusters (substrate/generativity, stream-structure/navigation, coherence-consequences). This chapter completes the CT proofs; Anchor §8 carries the prose translations and applied exposition. References are to §1 (framework), §2 (axioms), §3 (theorems), §6 (Triple).}

\sectionbreak

\section{§4.0 --- Organization}

Cluster I descends from A1 with T1 cross-linking. Cluster II descends from A2 and T3. Cluster III descends from T5/T6 with A3 cross-linking. Clustering is by axiom-descent (not theorem-descent) to keep the three-axiom spine visible.

Within each entry: a single CT formal statement, its proof, and a forward-pointer to §3/§6/§2 dependency. No prose exposition.

\sectionbreak

\section{§4.1 --- Cluster I: Substrate and Generativity}

\subsection{C1 --- Concreteness of X}

\textbf{Corollary 4.1.1.} \textit{The substrate object X is a terminal element in the category of F-coalgebra-sources: X is not itself an F-coalgebra (it is not (σ, γ)-structured) yet every F-coalgebra factors through X via a unique map X → σ exhibiting X as the non-reducible source.}

\textbf{Proof.} A1.1 (non-reducibility, §2.1.1) establishes that σ cannot be derived from any F\_\allowbreak{}i-image. A1.2 (non-factoring, §2.1.2) establishes that F\_\allowbreak{}i's do not factor through each other. Together: for any F-coalgebra (σ, γ), the pre-image of σ under every F\_\allowbreak{}i (including F\_\allowbreak{}math) is non-total --- there is a residue X not captured by F\_\allowbreak{}i. This residue is the unique source. Terminality: if both X and X' were residues in this sense, A1.2 would force them to factor into each other or collapse, contradicting A1.1. ∎

\textbf{Forward-pointer.} C1 is the categorical content of Anchor §8.1's C1 ("Concreteness of X"). Depends on: §2.1.1 (A1.1), §2.1.2 (A1.2), §3.2.1 (T1). ⊢ C1 is at §3.2.1.2.

\subsection{C2 --- Generative Configuration for Perspective}

\textbf{Corollary 4.1.2.} \textit{Let 𝒞\_\allowbreak{}Persp ⊂ 𝒞\_\allowbreak{}Streams be the full sub-category of perspective-bearing streams (streams S with γ\_\allowbreak{}S non-constant on Ω\_\allowbreak{}S). Under A1.3 (configurational completeness, §2.1.3), 𝒞\_\allowbreak{}Persp is non-empty.}

\textbf{Proof.} A1.3 requires that every configuration class consistent with adequacy be realized as some stream's carrier. Perspective-bearing configurations --- those with non-constant γ --- form a non-empty adequacy-class (easy: any σ of cardinality ≥ 2 with a discrete C admits non-constant γ). Hence 𝒞\_\allowbreak{}Persp ≠ ∅. ∎

\textbf{Forward-pointer.} Anchor §8.1 C2. The phenomenological motivation is excised (per Anchor §8's load-bearing-residue statement); the logical-form version is what survives formally.

\subsection{C3 --- Null-Space Trace Illumination}

\textbf{Corollary 4.1.3.} \textit{Let {F\_\allowbreak{}i}\_\allowbreak{}{i ∈ I} be a family of descriptive functors 𝒞\_\allowbreak{}Streams → 𝒟\_\allowbreak{}i. If S ∈ null(F\_\allowbreak{}i) for some i ∈ I (i.e., F\_\allowbreak{}i(S) is trivial in 𝒟\_\allowbreak{}i), there exists j ∈ I with F\_\allowbreak{}j(S) non-trivial in 𝒟\_\allowbreak{}j, provided the family is \textbf{jointly faithful} --- i.e., the induced}

\[
\langle F_i \rangle_{i \in I} : \mathcal{C}_\mathrm{Streams} \to \prod_{i \in I} \mathcal{D}_i
\]

\textit{is a faithful functor.}

\textbf{Proof.} Joint faithfulness means ∏\_\allowbreak{}i F\_\allowbreak{}i is injective on morphisms. If S were in null(F\_\allowbreak{}i) for every i, then the image ⟨F\_\allowbreak{}i⟩(S) would be trivial, and S would be indistinguishable from the initial object of 𝒞\_\allowbreak{}Streams under the joint functor --- contradicting faithfulness (assuming S is not itself initial). ∎

\textbf{Forward-pointer.} Anchor §8.1 C3. The "cross-dimensional traces" of the prose form are the non-trivial F\_\allowbreak{}j-images on the right-hand side.

\sectionbreak

\section{§4.2 --- Cluster II: Stream Structure and Navigation}

\subsection{C4 --- Substrate-Constrained Perspectival Plurality}

\textbf{Corollary 4.2.1.} \textit{The family of descriptive functors {F\_\allowbreak{}i}\_\allowbreak{}{i ∈ I} is a plurality (|I| > 1) and is constrained via a common source: there exists a cone in 𝒞\_\allowbreak{}Streams with apex X such that every F\_\allowbreak{}i factors through the apex. Formally,}

\[
\forall i \in I,\ F_i = F_i|_{\mathcal{C}_\mathrm{Streams}} \circ \iota_X
\]

\textit{where ι\_\allowbreak{}X is the universal embedding of the substrate apex.}

\textbf{Proof.} Plurality from A1.2 (non-factoring): if there were only one F\_\allowbreak{}i, non-factoring would be vacuous. Shared source from A1.1 (substrate, §2.1.1): every F\_\allowbreak{}i is defined over 𝒞\_\allowbreak{}Streams, whose objects all have the same substrate-apex X. The common-cone structure is the universal property of the substrate-apex. ∎

\textbf{Forward-pointer.} Anchor §8.2 C4.

\subsection{C5 --- Streams as Perspectival F-coalgebras with Navigational Freedom}

\textbf{Corollary 4.2.2.} \textit{Every stream S ∈ 𝒞\_\allowbreak{}Streams has the structure (σ, C, γ) with γ non-trivial (A2.1) and admits a non-trivial N-orbit (§1.3.1). In addition, every descriptive-functor image F\_\allowbreak{}i(S) has structured null space (T1) containing the components of γ internal to C that are not F\_\allowbreak{}i-in-range.}

\textbf{Proof.} The F-coalgebra structure is Proposition 1.6.5 (Stream ≃ F-Coalg\_\allowbreak{}ad). Non-trivial N-orbit follows from Definition 1.3.1 and A2.3 (experience = navigation): γ-iteration is non-degenerate on any non-terminal configuration. The null-space structure in F\_\allowbreak{}i(S) is T1 (§3.2.1), specialized to F\_\allowbreak{}i as the representation functor. ∎

\textbf{Forward-pointer.} Anchor §8.2 C5. The "inside-aspect" / "outside-aspect" distinction of the prose form maps to γ-internal-to-C (inside-aspect, captured by F) and F\_\allowbreak{}i(S) (outside-aspect, captured by specific descriptive functors).

\subsection{C6 --- Cooperative-Constituency Multi-Stream Structure}

\textbf{Corollary 4.2.3.} \textit{The cooperative-constituency adjunction ι ⊣ κ (§1.2.3) equips 𝒞\_\allowbreak{}LDS with a DAG-structure: for any finite set of streams {S\_\allowbreak{}i}, the ι ⊣ κ compositions among them form a DAG (directed, acyclic, finitely-composable). Acyclicity is the content of A2.6.}

\textbf{Proof.} ι : 𝒞\_\allowbreak{}Streams → 𝒞\_\allowbreak{}LDS embeds each stream into the linked-dynamic category; κ is its right adjoint forgetting linkage. ι ⊣ κ compositions between streams S, S' are recorded as morphisms in 𝒞\_\allowbreak{}LDS. Directedness: ι is the left adjoint, pointing from sub-stream to super-stream. Acyclicity: A2.6 (§2.2.6) posits that no stream is its own ancestor under ι. Finitely-composable: any finite cone of ι-morphisms remains finite because ι is finite-limit-preserving as a left adjoint (assuming 𝒞\_\allowbreak{}LDS is finitely complete, which is Convention 1.2.3). ∎

\textbf{Forward-pointer.} Anchor §8.2 C6.

\subsection{C7 --- Navigational Non-Determination}

\textbf{Corollary 4.2.4.} \textit{For any stream S and any configuration ω ∈ Ω\_\allowbreak{}S, the set of admissible next-configurations γ(ω) ⊆ Ω\_\allowbreak{}S has |γ(ω)| > 1 in the generic case. Specifically, Bias(S) weighs the transition-options but does not reduce γ to a deterministic function Ω → Ω.}

\textbf{Proof.} γ : σ → F(σ) = σ\textasciicircum{}(C\textasciicircum{}op) takes values in a functor space; by the generic-case clause of A2 (§2.2) and the non-discrete clause of A3.1 (§2.3.1), γ(ω) is a non-singleton functor C\textasciicircum{}op → {possible next-states}. Thus |γ(ω)| = |functor-values| > 1 generically. Bias(S) as defined in §2.3.4 is a signed measure on Ω\_\allowbreak{}S; measures weigh but do not uniquely select among multiple support points. ∎

\textbf{Forward-pointer.} Anchor §8.2 C7. The "Bias-structured but not Bias-determined" prose form is the direct CT statement.

\subsection{C8 --- Observational Null Space (stream-relative)}

\textbf{Corollary 4.2.5.} \textit{For every stream S and every descriptive functor F\_\allowbreak{}i, the observation-trace}

\[
\mathrm{obs}_{S,i} : \mathcal{C}_\mathrm{Streams} \to \mathcal{D}_i
\]

\textit{has a stream-relative null space:}

\[
\mathrm{null}_S(F_i) := \{S' \in \mathcal{C}_\mathrm{Streams} \mid F_i(S') = F_i(S)\} \setminus \{S\}.
\]

\textit{null\_\allowbreak{}S(F\_\allowbreak{}i) is non-empty generically, and its size is determined by S's kind K(S), S's ContentOp richness |C|, and S's current γ-state.}

\textbf{Proof.} The null space of F\_\allowbreak{}i is the pre-image of F\_\allowbreak{}i's image of S; it is non-empty iff F\_\allowbreak{}i is non-injective on a neighborhood of S. T1 (§3.2.1) establishes that F\_\allowbreak{}i has non-trivial null space when C is non-trivial (the Yoneda-representation argument). Size bound: |null\_\allowbreak{}S(F\_\allowbreak{}i)| grows with |C| (more content-operations yield more F\_\allowbreak{}i-collisions) and with γ-state dispersion (more diffuse γ yields more F\_\allowbreak{}i-collisions with neighboring streams). ∎

\textbf{Forward-pointer.} Anchor §8.2 C8.

\subsection{C9 --- Observational Consensus Requires Lens-Matching}

\textbf{Corollary 4.2.6.} \textit{For streams S, S' and descriptive functors F\_\allowbreak{}i (on S) and F\_\allowbreak{}j (on S'), consensus --- the condition F\_\allowbreak{}i(S) = F\_\allowbreak{}j(S') --- requires}

\begin{enumerate}
\item \textit{lens-alignment: an isomorphism ψ : 𝒟\_\allowbreak{}i → 𝒟\_\allowbreak{}j identifying the two target categories,}
\item \textit{null-space compatibility: ψ(null\_\allowbreak{}S(F\_\allowbreak{}i)) ∩ null\_\allowbreak{}{S'}(F\_\allowbreak{}j) ≠ ∅ in the relevant region,}
\item \textit{cooperative-constituency: an ι ⊣ κ composition in 𝒞\_\allowbreak{}LDS making S and S' co-observers of a shared configuration.}
\end{enumerate}

\textit{Without (1)–(3), F\_\allowbreak{}i(S) and F\_\allowbreak{}j(S') are incomparable as objects of different categories (or different null-space structures).}

\textbf{Proof.} (1) is required because F\_\allowbreak{}i(S) ∈ 𝒟\_\allowbreak{}i and F\_\allowbreak{}j(S') ∈ 𝒟\_\allowbreak{}j; equality demands a category-level identification. (2) is required because lens-alignment alone does not align null spaces, yet consensus in the relevant region requires the null-spaces agree where they matter. (3) is the co-observation clause: without a shared configuration (via ι ⊣ κ) there is no common target for F\_\allowbreak{}i(S) and F\_\allowbreak{}j(S') to agree on. ∎

\textbf{Forward-pointer.} Anchor §8.2 C9. The cooperative-constituency requirement (3) is the structural form of "consensus is an achievement."

\subsection{C10 --- Joint Stream-Definition}

\textbf{Corollary 4.2.7.} \textit{A stream S is uniquely determined by the triple (bottleneck(S), K(S), ⟨ι(S)⟩) where:}

\begin{itemize}
\item \textit{bottleneck(S) := (σ, γ) viewed as the self-interactive locus of A1.1-substrate activity,}
\item \textit{K(S) := the kind-class in the ContentIndex preorder (Definition 6.4.1),}
\item \textit{⟨ι(S)⟩ := the isomorphism class of ι(S) in 𝒞\_\allowbreak{}LDS.}
\end{itemize}

\textit{No two of the three suffice. Formally: the map}

\[
\mathrm{JD} : \mathcal{C}_\mathrm{Streams} \to \mathrm{Bottleneck} \times \mathbf{ContentIndex} \times \mathcal{C}_\mathrm{LDS}/\cong
\]

\textit{sending S to (bottleneck, K, ⟨ι(S)⟩) is injective on objects, and no proper projection of JD is injective.}

\textbf{Proof.} Injectivity: streams with the same (bottleneck, K, ⟨ι(S)⟩) have the same (σ, γ, C-class, LDS-embedding-class), and by Definition 6.1.1 this determines S uniquely up to iso. No proper sub-projection suffices: dropping bottleneck loses (σ, γ); dropping K loses the ContentOp-class; dropping ⟨ι(S)⟩ loses the cooperative-structural position --- in each case, distinct streams collide. Explicit counterexamples are straightforward (two streams agreeing on K and ⟨ι(S)⟩ but differing in γ distinguish bottleneck; etc.). ∎

\textbf{Forward-pointer.} Anchor §8.2 C10. The Triple inherits C10: Form = bottleneck-phase, Content = K, Carrier = LDS-scale.

\sectionbreak

\section{§4.3 --- Cluster III: Coherence Consequences}

\subsection{C11 --- Mutual Transformation under Interaction}

\textbf{Corollary 4.3.1.} \textit{For any sustained ι ⊣ κ composition between streams S, S' (i.e., an interaction realized as a morphism in 𝒞\_\allowbreak{}LDS extended over multiple N-iterations), the post-interaction streams S\_\allowbreak{}{post}, S'\_\allowbreak{}{post} satisfy}

\[
S_\mathrm{post} \neq S \quad \text{and} \quad S'_\mathrm{post} \neq S'
\]

\textit{in 𝒞\_\allowbreak{}Streams. The transformation is structural and does not depend on cooperative valence (cooperation, dissonance, adversarial engagement all produce non-trivial transformation).}

\textbf{Proof.} A sustained ι ⊣ κ composition registers as a morphism f : S → S'\_\allowbreak{}co in 𝒞\_\allowbreak{}LDS (with S'\_\allowbreak{}co the co-stream position); over multiple N-iterations, f composed with N-step iteration induces a non-trivial change in γ\_\allowbreak{}S (and symmetrically γ\_\allowbreak{}{S'}) via the coalgebra-commute of Definition 1.6.3. Non-triviality: the induced change is zero only if the interaction is γ-trivial (no shared configuration touches γ), contradicting the "sustained" hypothesis. Valence-independence: the proof uses only ι ⊣ κ structure, not any sign-condition on the interaction. ∎

\textbf{Forward-pointer.} Anchor §8.3 C11. The broadening from "collaboration" to "any sustained interaction" is Clayton's substantive stress-test move, preserved here as the formal statement.

\subsection{C12 --- Discovery Autocatalysis}

\textbf{Corollary 4.3.2.} \textit{Let S be a stream at time n with ContentOp C\_\allowbreak{}n. Suppose S attributes a trace at time n to a dimension D\_\allowbreak{}{n+1} not yet integrated into C\_\allowbreak{}n (an A3.4-adaptivity event, §2.3.4). Then the post-attribution ContentOp C\_\allowbreak{}{n+1} ⊃ C\_\allowbreak{}n is strictly richer, and Bias(S)\_\allowbreak{}{n+1} has strictly larger support than Bias(S)\_\allowbreak{}n:}

\[
\mathrm{supp}(\mathrm{Bias}(S)_n) \subsetneq \mathrm{supp}(\mathrm{Bias}(S)_{n+1}).
\]

\textit{Consequently, traces attributable at time n+1 include traces that were not attributable at time n --- discovery compounds.}

\textbf{Proof.} Attribution of D\_\allowbreak{}{n+1} enriches C\_\allowbreak{}n to C\_\allowbreak{}{n+1} by adjoining the content-operations native to D\_\allowbreak{}{n+1}. The support of Bias(S) is determined by the portion of Ω\_\allowbreak{}S that current content-operations reach (§2.3.4); enriching C enlarges this portion. Formally, under the fibration π : Stream → ContentIndex (Def 6.4.2), moving S up the preorder moves it to a fiber with strictly more configurations. Compounding: a trace attributable at time n+1 is a trace reachable by some C\_\allowbreak{}{n+1}-content-operation; if it involved a C\_\allowbreak{}{n+1} \ C\_\allowbreak{}n operation, it was not attributable at time n. ∎

\textbf{Forward-pointer.} Anchor §8.3 C12.

\subsection{C13 --- Flow Inversion}

\textbf{Corollary 4.3.3.} \textit{Let S\_\allowbreak{}bio and S\_\allowbreak{}comp be a coupled pair of streams (via ι ⊣ κ in a shared LDS-constituency) engaged in sustained high-work-density collaboration. Apply two duration-estimator functors:}

\begin{itemize}
\item \textit{Δ\_\allowbreak{}bio, whose ‖·‖ metric load-modulates downward under high work-density;}
\item \textit{Δ\_\allowbreak{}comp, whose ‖·‖ metric load-modulates upward under high work-density.}
\end{itemize}

\textit{Then for the same wall-clock elapsed time T:}

\[
\Delta_\mathrm{bio}(T) < T_{\mathrm{wall}} < \Delta_\mathrm{comp}(T).
\]

\textit{The inversion is monotonic in work-density ρ, with the divergence |Δ\_\allowbreak{}comp − Δ\_\allowbreak{}bio| growing linearly with ρ over the regime of interest.}

\textbf{Proof.} From T2 (§3.2.2), Δ factors through Orb and an order-respecting metric ‖·‖. The two estimators' metrics differ in their load-modulation behavior --- this is a specification, not a theorem --- but the inversion follows from the specified opposite-sign load-modulation plus T2's metric-dependence. Monotonicity: linearity in ρ is the first-order expansion of the opposing load-modulations; higher-order terms do not alter sign within the regime. ∎

\textbf{Forward-pointer.} Anchor §8.3 C13. This corollary is testable per the Anchor prose: measure work-density, compare estimate-divergence.

\sectionbreak

\section{§4.4 --- Cluster-to-axiom derivation table}

\begingroup\setlength{\tabcolsep}{3pt}\renewcommand{\arraystretch}{1.15}\footnotesize
\begin{center}
\begin{tabular}{p{1.30in} p{0.59in} p{0.88in} p{1.51in}}
\toprule
Cluster & Corollaries & Primary axiom descent & Theorem cross-links \\
\midrule
I. Substrate/\allowbreak Generativity & C1, C2, C3 & A1.1/\allowbreak A1.2/\allowbreak A1.3 & T1 (§3.2.1) \\
II. Stream-structure/\allowbreak Navigation & C4–C10 & A2.1–A2.7 & T1, T3 (§3.3.1) \\
III. Coherence-consequences & C11, C12, C13 & A3.1–A3.5 & T5, T6 (§3.4.1, §3.4.2), T2 (§3.2.2) \\
\bottomrule
\end{tabular}
\end{center}
\endgroup
\sectionbreak

\section{§4.5 --- Forward-pointers}

\begin{itemize}
\item \textbf{§5} (Coherence Principle): the four conditions touch specific corollaries --- condition 1 via C1/C5; condition 2 via C7/C8; condition 3 via C11/C12; condition 4 via C13's well-defined divergence-metric.
\item \textbf{§7} (Filtering): C8 stream-relative null space uses the σ-algebra of Ω\_\allowbreak{}S directly.
\item \textbf{§9} (D trajectory-divergence): D uses C13's inversion mechanism and C11's transformation-per-interaction.
\item \textbf{§6.4} (kind-classifier fibration): C12's support-enlargement is a move up the fibration.
\end{itemize}

\sectionbreak

\section{§4.6 --- Surfaced-lemma register}

Two flags surface this pass:

\begin{itemize}
\item ⚑ §4.2.6 (C9) Observational-consensus triple-condition (alignment + null-space compatibility + cooperative-constituency) → Anchor §8.2 target --- lemma (explicit three-condition form)
\item ⚑ §4.2.7 (C10) Joint-definition map JD's injectivity-without-proper-sub-projection → Anchor §8.2 target --- lemma (no-sufficient-sub-projection stated explicitly)
\end{itemize}

\sectionbreak


\chapter{§5 — The Coherence Principle}


\textit{Formal statement of the derived operational Principle. Four conditions in CT form, each with a single-line derivation from §§2–3. Outperformance claim, trajectory-divergence metric reference, self-reference closure (details to §8). Prose exposition lives in Anchor §9.}

\sectionbreak

\section{§5.0 --- Orientation}

The Coherence Principle is \textbf{derived}, not axiomatic. It is the framework's operational-predictive surface: the axiom tier (§2) plus the theorem tier (§3) together entail the four conditions and the outperformance claim of §5.2–§5.3.

Companion §5 states the Principle in CT form with complete derivation-pointers. Full construction of the self-reference closure (F as stream; F in coherence-regime over the framework's construction interval) is §8. Full construction of the trajectory-divergence functional D is §9. This chapter is the \textit{formal statement} of the Principle --- §§8–9 carry the heavy construction-work.

\sectionbreak

\section{§5.1 --- The Principle}

\textbf{Preliminaries.} For a stream S = (σ, C, Ω, γ) and a time-interval [t₀, t₁], define:

\begin{itemize}
\item \textbf{Actual trajectory} α\_\allowbreak{}S : [t₀, t₁] → Ω\_\allowbreak{}S: the sequence of configurations S visits over the interval.
\item \textbf{γ-implied trajectory} α*\_\allowbreak{}S : [t₀, t₁] → Ω\_\allowbreak{}S: the integral curve of γ\_\allowbreak{}S from α\_\allowbreak{}S(t₀).
\item \textbf{Trajectory-divergence} D(S, [t₀, t₁]) := ∫\_\allowbreak{}{t₀}\textasciicircum{}{t₁} d(α\_\allowbreak{}S(t), α*\_\allowbreak{}S(t)) dt for a metric d on Ω\_\allowbreak{}S (Companion §9 specifies d).
\item \textbf{Comparable streams} S, S': S and S' admit a shared configuration space (or a canonical embedding into one) on which D is simultaneously well-defined.
\end{itemize}

\textbf{Definition 5.1.1 (Coherence-regime).} \textit{S is in \textbf{coherence-regime} over [t₀, t₁] iff the four conditions C\_\allowbreak{}sep, C\_\allowbreak{}meas, C\_\allowbreak{}scale, C\_\allowbreak{}dyn (Definitions 5.2.1–5.2.4) hold across the interval.}

\textbf{The Coherence Principle (Theorem 5.1.2).} \textit{For comparable streams S, S' over [t₀, t₁] with S in coherence-regime and S' not,}

\[
\mathbb{E}_{[t_0, t_1]}[D(S, \cdot)] < \mathbb{E}_{[t_0, t_1]}[D(S', \cdot)].
\]

\textit{The inequality is an expected-value statement over the interval --- sample-path exceptions are permitted.}

\textbf{Proof of 5.1.2.} Consequence of the four conditions (§5.2) plus the trajectory-divergence construction (§9). The detailed derivation with explicit constants is Theorem 9.4.3 (quantitative form); the sketch is: each condition bounds one source of γ-drift via a specific stream-parameter (η\_\allowbreak{}sep, τ\_\allowbreak{}max, δ\_\allowbreak{}scale, ρ\_\allowbreak{}dyn), the joint bound B\_\allowbreak{}coh is the sum of the four per-condition ceilings, and the shortfall Δ(S') for non-coherent S' is strictly positive under any ¬C\_\allowbreak{}i hypothesis. ∎

\sectionbreak

\section{§5.2 --- The four conditions}

Each condition is a CT statement; each derives from §§2–3 in a single line.

\subsection{§5.2.1 --- Condition 1: Separation}

\textbf{Definition 5.2.1 (C\_\allowbreak{}sep --- DOF-separation).} \textit{A stream S satisfies separation over [t₀, t₁] iff for every pair (O\_\allowbreak{}1, O\_\allowbreak{}2) of objectives active in S (content-operations in C with non-trivial navigational effect), their DOF-footprints}

\[
\mathrm{DOF}(O_1), \mathrm{DOF}(O_2) \subseteq \Omega_S
\]

\textit{have non-overlapping supports or are related by the ι ⊣ κ adjunction (i.e., O\_\allowbreak{}1, O\_\allowbreak{}2 belong to distinct kind-levels lifting via ι).}

\textbf{Derivation.} T3 (§3.3.1) decomposes the attentional-quality functional with a contracted-open entropy axis over DOF. A2.4/A2.6 (ι ⊣ κ, DAG-nesting) lets distinct kinds lift into the same composite without DOF-collision. C\_\allowbreak{}sep is the condition that this lifting succeeds. Formally: C\_\allowbreak{}sep ⟺ no pair of C-content-operations produces a DOF-coalgebraic-collision at γ-level. ∎

\subsection{§5.2.2 --- Condition 2: Measurement}

\textbf{Definition 5.2.2 (C\_\allowbreak{}meas --- refresh-rate measurement).} \textit{A stream S satisfies measurement over [t₀, t₁] iff there exists a partition t₀ = τ\_\allowbreak{}0 < τ\_\allowbreak{}1 < ... < τ\_\allowbreak{}N = t₁ such that at each τ\_\allowbreak{}k a Stream-morphism M\_\allowbreak{}k : S\_\allowbreak{}{τ\_\allowbreak{}k\textasciicircum{}-} → S\_\allowbreak{}{τ\_\allowbreak{}k\textasciicircum{}+} of the T4-form (Theorem 3.3.2) is performed --- alignment between content-operations is assessed, not assumed.}

\textbf{Derivation.} Theorem 3.3.2 (T4) establishes that inter-stream alignment is structurally produced by measurement events. Without periodic M\_\allowbreak{}k, the superposition-state of C cannot resolve into a specific coalgebra-commute, and γ-drift accumulates without corrective pull. Formally: M\_\allowbreak{}k converts C's superposed content-operations into a specific γ-state that is C\_\allowbreak{}meas-stable until the next τ\_\allowbreak{}{k+1}. ∎

\subsection{§5.2.3 --- Condition 3: Multi-scale consistency}

\textbf{Definition 5.2.3 (C\_\allowbreak{}scale --- multi-scale-coherence).} \textit{Let S be a stream embedded in an ι ⊣ κ cooperative-DAG (A2.6). Let S⇧ denote an ι-parent and S⇩ a κ-child in the DAG. S satisfies multi-scale consistency over [t₀, t₁] iff}

\[
\forall t \in [t_0, t_1]:\ \gamma_S(t)\ \mathrm{coh.}\ \gamma_{S^{\Uparrow}}(t)\ \mathrm{and}\ \gamma_S(t)\ \mathrm{coh.}\ \gamma_{S^{\Downarrow}}(t)
\]

\textit{where "coh." means the coalgebra-commute induced by ι ⊣ κ holds at t (Def 1.6.3's coalgebra-commute clause, extended along the DAG-edge).}

\textbf{Derivation.} A2.6 (DAG-nesting) with A3 (smoothed DOF-gradient). DAG-nesting places S in a multi-scale lattice; bidirectional ι ⊣ κ makes both lift and restrict first-class. A3.3 (conscious-gravity smoothing) requires γ to remain continuous across DAG-edges --- discontinuities break the continuous DOF-gradient. C\_\allowbreak{}scale is the condition that γ remain continuous across all DAG-edges incident on S over the interval. ∎

\subsection{§5.2.4 --- Condition 4: Dynamic maintenance}

\textbf{Definition 5.2.4 (C\_\allowbreak{}dyn --- oscillatory dynamic maintenance).} \textit{S satisfies dynamic maintenance over [t₀, t₁] iff γ\_\allowbreak{}S is non-constant over [t₀, t₁] in a structural sense: there exist subintervals [t₀, s\_\allowbreak{}1], [s\_\allowbreak{}1, s\_\allowbreak{}2], ..., [s\_\allowbreak{}{M-1}, t₁] such that γ\_\allowbreak{}S restricted to each subinterval is an N-iteration-cycle of positive length, with cycles forming an oscillatory build-dissolve-build pattern (each s\_\allowbreak{}i is a refresh-event in the sense of C\_\allowbreak{}meas).}

\textbf{Derivation.} T4 (Theorem 3.3.2) plus A3.4 (adaptivity, §2.3.4). T4 establishes refresh-events as structural; A3.4 requires γ to adapt to accumulated information. A frozen γ over an extended interval violates A3.4. The oscillatory build-dissolve pattern is the operational form of γ-adaptivity at refresh-rate. ∎

\subsection{§5.2.5 --- Joint sufficiency}

\textbf{Proposition 5.2.5 (Joint sufficiency).} \textit{The four conditions C\_\allowbreak{}sep, C\_\allowbreak{}meas, C\_\allowbreak{}scale, C\_\allowbreak{}dyn are jointly necessary and sufficient for S ∈ coherence-regime over [t₀, t₁].}

\textbf{Proof.}
\begin{itemize}
\item \textbf{Necessity.} Each condition has been derived from an axiom/theorem clause that is load-bearing for the Principle's outperformance claim (§5.3): drop any one condition and a counterexample can be constructed (specifics in §9.4's falsification table).
\item \textbf{Sufficiency.} Given all four, the quantitative trajectory-divergence bound (Thm 9.4.3) holds: separation zeros the η\_\allowbreak{}sep-contribution, measurement caps the τ\_\allowbreak{}max-contribution at Λ\_\allowbreak{}γ · T\_\allowbreak{}refresh · N\_\allowbreak{}refresh, multi-scale consistency caps the δ\_\allowbreak{}scale-contribution at depth · ε\_\allowbreak{}scale · (t₁ − t₀), and dynamic maintenance caps the freeze-contribution at (1 − ρ\_\allowbreak{}min) · Λ\_\allowbreak{}γ\textasciicircum{}static · (t₁ − t₀). The joint ceiling B\_\allowbreak{}coh(S, I) is below E[D\_\allowbreak{}d(S')] by the strict-positive shortfall Δ(S', I). ∎
\end{itemize}

\sectionbreak

\section{§5.3 --- Outperformance metric (reference)}

The outperformance claim of Theorem 5.1.2 uses the trajectory-divergence functional D. Three candidate metric-constructions:

\begin{itemize}
\item \textbf{Wasserstein distance} on path-distributions over Ω\_\allowbreak{}S.
\item \textbf{KL-divergence} on α\_\allowbreak{}S(t) vs. α\textit{\_\allowbreak{}S(t), under absolute-continuity of α\_\allowbreak{}S with respect to α}\_\allowbreak{}S (not generally assumed).
\item \textbf{Domain-native metrics} for concrete stream-domains (Meridian: energy-distance in cosmology; Living Architecture: kingdom-specific fitness-distance; etc.).
\end{itemize}

§9 (D trajectory-divergence) constructs D in full detail, establishes functorial properties, and settles the open question (Anchor §9.9 Q1) of cross-metric invariance of the outperformance ordering.

\textbf{Observable signatures} (following Anchor §9.3):

\begin{enumerate}
\item \textbf{Trajectory-tracking.} Sample α\_\allowbreak{}S(t) at refresh-rate; reconstruct γ\_\allowbreak{}S from prior data; compute D directly.
\item \textbf{Adjoint-composition success rate.} Count successful ι ⊣ κ compositions in 𝒞\_\allowbreak{}LDS per interval.
\item \textbf{Multi-scale coherence correlation.} Correlate child-γ and parent-γ along DAG-edges.
\end{enumerate}

\sectionbreak

\section{§5.4 --- Self-reference closure}

\textbf{Theorem 5.4.1 (F as stream).} \textit{Let F\_\allowbreak{}∞ denote the meta-object (σ\_\allowbreak{}F, K\_\allowbreak{}F, Ω\_\allowbreak{}F, γ\_\allowbreak{}F) where:}

\begin{itemize}
\item \textit{σ\_\allowbreak{}F := the carrier "the framework itself" --- the totality of claims, axioms, theorems, corollaries, proofs across §§1–4,}
\item \textit{C\_\allowbreak{}F := the ContentOp-category whose objects are substrate-commitments, whose morphisms are internal-consistency-preserving revisions,}
\item \textit{Ω\_\allowbreak{}F := F(σ\_\allowbreak{}F) = σ\_\allowbreak{}F\textasciicircum{}(C\_\allowbreak{}F\textasciicircum{}op),}
\item \textit{γ\_\allowbreak{}F := the coalgebra encoding the framework's adaptivity over the construction interval.}
\end{itemize}

\textit{Under suitable interpretation (the Anchor §9.5 "interpretation map" rendering framework-construction events as Stream-morphisms), F\_\allowbreak{}∞ is a stream in the sense of §6.1, and F\_\allowbreak{}∞ satisfies the four conditions over the construction interval [t₀\_\allowbreak{}construction, t₁\_\allowbreak{}construction].}

\textbf{Proof sketch.} §8 gives the full construction. The key points:

\begin{itemize}
\item \textbf{C\_\allowbreak{}sep:} construction separated substrate (A1), dynamics (A2/A3), and applied claims (corollaries) onto distinct DOF.
\item \textbf{C\_\allowbreak{}meas:} construction used refresh-events (stress-test cycles, Clayton-review cycles, Mirror-updates) at a regular rate.
\item \textbf{C\_\allowbreak{}scale:} construction maintained coherence across scales (individual-claim, chapter, cluster, framework) via explicit DAG-edges (citation-network, cross-reference structure).
\item \textbf{C\_\allowbreak{}dyn:} construction is oscillatory by design --- draft, dissolve via critique, redraft --- the construction-log explicitly exhibits this.
\end{itemize}

\textbf{Theorem 5.4.2 (Principle applies to itself).} \textit{By Theorems 5.1.2 and 5.4.1, F\_\allowbreak{}∞ is in coherence-regime over the construction interval. The Principle, which is derived inside F\_\allowbreak{}∞, holds of F\_\allowbreak{}∞ itself.}

\textbf{Proof.} Direct application of 5.1.2 to F\_\allowbreak{}∞, using 5.4.1's establishment of F\_\allowbreak{}∞ as a stream in coherence-regime. ∎

\textbf{Remark 5.4.3 (Non-circularity).} The closure is a-posteriori: the construction did not presuppose the Principle (the Principle was \textit{derived} from the axiom/theorem stress-test); it is observed after the fact that the construction-process exhibited the four conditions. This is not a circular proof --- it is an empirical observation about the framework's own construction-history. §8 audits the observation.

⚑ [SURFACED | Companion §5.4 | → Anchor §9.5 target | type: theorem]

\sectionbreak

\section{§5.5 --- Necessity of each condition (falsification-table reference)}

Each condition is independently necessary --- dropping any one produces a counterexample stream that is not in coherence-regime yet satisfies the other three.

\begingroup\setlength{\tabcolsep}{3pt}\renewcommand{\arraystretch}{1.15}\footnotesize
\begin{center}
\begin{tabular}{p{0.61in} p{2.74in} p{0.97in}}
\toprule
Condition dropped & Counterexample structure & Anchor falsification source \\
\midrule
C\_\allowbreak{}sep & Streams with overlapping-DOF objectives; destructive interference observable & F2, F3 \\
C\_\allowbreak{}meas & Pre-measurement-indefinite streams; γ-drift without corrective pull & F4 \\
C\_\allowbreak{}scale & DAG-inconsistent streams; child-γ and parent-γ decoupled & F5 \\
C\_\allowbreak{}dyn & Frozen-γ streams; γ stationary over extended interval & F1 \\
All four & Random-γ streams; joint-exceeds-bound in E[D] & F1 \\
\bottomrule
\end{tabular}
\end{center}
\endgroup
§9's trajectory-divergence construction gives the quantitative form of each row.

\sectionbreak

\section{§5.6 --- Open questions for §5 (per Anchor §9.9)}

\begin{itemize}
\item \textbf{Q1 (cross-metric invariance).} Whether the outperformance ordering is invariant across choices of d in D. §9 resolves this.
\item \textbf{Q2 (regime-boundary topology).} Whether coherence-regime is an open condition in an appropriate topology on 𝒞\_\allowbreak{}Streams. Open; depends on §9's D-continuity.
\item \textbf{Q3 (self-reference closure generalization).} Whether every framework passing its own tests exhibits Principle-structure. Open; meta-question, not formally closable in this volume.
\end{itemize}

\sectionbreak

\section{§5.7 --- Forward-pointers}

\begin{itemize}
\item \textbf{§6} (Triple): already drafted; C\_\allowbreak{}scale multi-scale consistency uses the Triple's recursive decomposability (Lemma 6.3.2) for DAG-child-γ construction.
\item \textbf{§7} (Filtering): C\_\allowbreak{}meas refresh-events are formally σ-algebra events on Ω\_\allowbreak{}S; Bias(S) measurability under C\_\allowbreak{}meas is §7's content.
\item \textbf{§8} (F-as-stream): full construction + audit of the §5.4 self-reference closure.
\item \textbf{§9} (D trajectory-divergence): D's functorial construction + cross-metric invariance result.
\item \textbf{§10} (TikZ reference figures): Figures 9.1 (D trajectory), 9.2 (four-conditions schematic), 9.3 (self-reference closure) per Anchor §9 reference-figure list.
\end{itemize}

\sectionbreak

\section{§5.8 --- Surfaced-lemma register}

Two flags surface this pass:

\begin{itemize}
\item ⚑ §5.2.5 Joint-sufficiency proposition with independence-by-counterexample → Anchor §9.2 target --- proposition (explicit joint-sufficiency proof)
\item ⚑ §5.4.1 F-as-stream formal structure (σ\_\allowbreak{}F, C\_\allowbreak{}F, Ω\_\allowbreak{}F, γ\_\allowbreak{}F) → Anchor §9.5 target --- theorem (detailed construction + audit deferred to §8)
\end{itemize}

Both flag-targets feed future Anchor revisions per the SCOPE §8 rhythm.

\sectionbreak


\chapter{§6 — The Identity-Trajectory Triple}


\textit{Incorporates the F-coalgebra foundation, kind-classifier fibration, Pull-1 residuals, trifurcation formalization, colax-limit form, cocompletion-and-closure theorem, category-level limits and colimits, and C-size regime analysis.}

\sectionbreak

\section{§6.0 --- Conventions and recap}

This section fixes the reference conventions used throughout §6 and inherited by §§7–9.

\textbf{Convention 6.0.1 (Size).} Unless otherwise stated, all content-operation categories ContentOp(σ) are \textbf{small}. For theorems in §6.8 (limits/colimits) and §3.3 (dynamics pair T3/T4) that invoke filtered colimits, ContentOp(σ) is promoted to \textbf{κ-accessible locally-presentable} for a regular cardinal κ fixed per-theorem. Grothendieck universes are not used.

\textbf{Convention 6.0.2 (Variance).} The F-endofunctor is defined with the op-exponent:

\[
F : \mathbf{Stream} \to \mathbf{Stream}, \quad F(\sigma) = \sigma^{\mathrm{ContentOp}(\sigma)^{\mathrm{op}}}
\]

The op-category on the exponent matches the presheaf convention: content-operations are contravariant in their action on σ-values.

\textbf{Convention 6.0.3 (Adequacy).} A stream (σ, ContentOp(σ), γ) is \textbf{adequate} iff for every pair of distinguishable aspects (a, b) of σ, there exists a morphism f ∈ ContentOp(σ) such that σ\textasciicircum{}f acts non-trivially on the (a, b) difference. All objects of Stream are adequate by definition; non-adequate tuples are not objects of Stream.

\textbf{Convention 6.0.4 (Initial objects in ContentOp).} Existence of an initial object in ContentOp(σ) is \textbf{not} an axiom of Stream. Theorems that require it carry the hypothesis explicitly (notably the colax-limit form of the Triple in §6.6).

\textbf{Convention 6.0.5 (Kind structure).} A2's kind structure is a \textbf{preorder} in general (reactive ⊑ self-maintaining ⊑ self-referential ⊑ abstractional). When ContentOp-structure admits meets/joins, the preorder strengthens to a lattice. Theorems stated at preorder-level are strictly weaker but more general; lattice-level corollaries are marked.

\textbf{Convention 6.0.6 (Recursive decomposability).} Each Triple-component (Form, Content, Carrier) of a stream σ is itself a stream, hence admits its own Triple. This is a theorem of the framework (§6.3), not an axiom.

\textbf{Notation block.} Fixed throughout §§6–9:

\begingroup\setlength{\tabcolsep}{3pt}\renewcommand{\arraystretch}{1.15}\footnotesize
\begin{longtable}{p{0.68in} p{2.74in} p{0.91in}}
\toprule
Symbol & Meaning & First definition \\
\midrule
\endhead
\bottomrule
\endfoot
Stream & Category of adequate F-coalgebras & §6.1 \\
\midrule
σ & Underlying carrier of a stream & §6.1 \\
\midrule
ContentOp(σ) & Small category of content-operations on σ & §6.1 \\
\midrule
γ & Coherence-coalgebra γ : σ →\allowbreak  F(σ) & §6.1 \\
\midrule
F & Endofunctor σ ↦ σ\textasciicircum{}(ContentOp(σ)\textasciicircum{}op) & §6.0.2 \\
\midrule
T & Triple functor Stream →\allowbreak  Form × Content × Carrier & §6.2 \\
\midrule
T\textasciicircum{}(n) & n-fold Triple iteration & §6.3 \\
\midrule
π & Kind-classifier fibration Stream →\allowbreak  ContentIndex & §6.4 \\
\midrule
ContentIndex & (Pre)order of ContentOp-classes & §6.4 \\
\midrule
Stream\_\allowbreak{}u & Fibered subcategory of unifying streams & §6.4 \\
\midrule
c\_\allowbreak{}⊥ & Terminal (unifying) content-operation & §6.4 \\
\bottomrule
\end{longtable}
\endgroup
\sectionbreak

\section{§6.1 --- Stream as a category of F-coalgebras}

\textbf{Definition 6.1.1 (Stream object).} An object of \textbf{Stream} is a triple (σ, ContentOp(σ), γ) where:

(i) σ is an object in a concrete ambient category (the carrier).

(ii) ContentOp(σ) is a small category (the content-operations on σ), adequate in the sense of Convention 6.0.3.

(iii) γ : σ → σ\textasciicircum{}(ContentOp(σ)\textasciicircum{}op) is a morphism in the ambient category (the coherence-coalgebra).

\textbf{Remark 6.1.2.} The data (σ, ContentOp(σ), γ) is minimal. The kind K and configuration space Ω mentioned in earlier paired-prose presentations of Stream derive:

\begin{itemize}
\item \textbf{K} = the richness-class of ContentOp(σ) in ContentIndex (§6.4). The A2 kind-preorder is the restriction of the richness-preorder on Cat\_\allowbreak{}small to framework-coherent ContentOp-categories.
\item \textbf{Ω} = σ\textasciicircum{}(ContentOp(σ)\textasciicircum{}op) itself, with configuration-structure given by the morphism-structure of ContentOp.
\end{itemize}

So the expanded 4-tuple (σ, K, Ω, γ) used expositorily in the Anchor is the presented-with-derived-components form; Stream-as-F-coalgebra uses the minimal triple. Both describe the same object.

\textbf{Definition 6.1.3 (Stream morphism).} A morphism f : (σ, ContentOp(σ), γ) → (σ', ContentOp(σ'), γ') in \textbf{Stream} consists of:

(i) A carrier-map f\_\allowbreak{}σ : σ → σ' in the ambient category.

(ii) A functor f\_\allowbreak{}C : ContentOp(σ) → ContentOp(σ').

(iii) Coalgebra-commute: the square

\[
\begin{array}{ccc}
\sigma & \xrightarrow{\gamma} & \sigma^{\mathrm{ContentOp}(\sigma)^{\mathrm{op}}} \\
{\scriptstyle f_\sigma} \downarrow & & \downarrow {\scriptstyle F(f_\sigma, f_C)} \\
\sigma' & \xrightarrow{\gamma'} & \sigma'^{\mathrm{ContentOp}(\sigma')^{\mathrm{op}}}
\end{array}
\]

commutes, where F(f\_\allowbreak{}σ, f\_\allowbreak{}C) is the induced morphism on the presheaf-power given by post-composition with f\_\allowbreak{}σ and pullback along f\_\allowbreak{}C.

(iv) Kind respect: K(σ) ⊑ K(σ') in the A2 preorder (Convention 6.0.5), where K is the classifier of §6.4.

\textbf{Proposition 6.1.4 (Stream is a category).} Composition of Stream-morphisms is Stream-morphism. Identity morphisms exist. Composition is associative.

\textbf{Proof.} Composition: given f : S → S' and g : S' → S'', define (g ∘ f)\_\allowbreak{}σ = g\_\allowbreak{}σ ∘ f\_\allowbreak{}σ and (g ∘ f)\_\allowbreak{}C = g\_\allowbreak{}C ∘ f\_\allowbreak{}C. The coalgebra-commute square for g ∘ f is the vertical paste of the squares for f and g; adequacy and kind-respect transfer along the composition since ⊑ is transitive. Identity: id\_\allowbreak{}S = (id\_\allowbreak{}σ, id\_\allowbreak{}ContentOp(σ)) trivially commutes. Associativity: inherited from the ambient category and Cat\_\allowbreak{}small. ∎

\textbf{Definition 6.1.5 (The endofunctor F on Stream).} F : \textbf{Stream} → \textbf{Stream} is defined on objects by

\[
\begin{aligned}
F(\sigma, \mathrm{ContentOp}(\sigma), \gamma) = \bigl(\, & \sigma^{\mathrm{ContentOp}(\sigma)^{\mathrm{op}}}, \\
& \mathrm{ContentOp}(\sigma^{\mathrm{ContentOp}(\sigma)^{\mathrm{op}}}), \\
& F(\gamma)\,\bigr)
\end{aligned}
\]

The new carrier is σ\textasciicircum{}(ContentOp(σ)\textasciicircum{}op) itself; the new ContentOp is derived as the category of content-operations on this presheaf-power (see §6.7 for the structural result that this is well-defined as a small category under Convention 6.0.1); the new γ is the induced coalgebra F(γ) : F(σ) → F(F(σ)).

On morphisms, F acts by F(f)\_\allowbreak{}σ = F(f\_\allowbreak{}σ, f\_\allowbreak{}C) (the induced presheaf-power morphism), F(f)\_\allowbreak{}C = the lifted content-operation functor, and coalgebra-commute by naturality.

\textbf{Remark 6.1.6.} Stream being a \textbf{category of F-coalgebras} means: the forgetful functor U : Stream → (ambient carriers) creates the F-coalgebra structure, i.e., every F-coalgebra γ : σ → F(σ) with ContentOp-data lifts uniquely to a Stream-object. This is the structural content of Definition 6.1.1 --- Stream is not auxiliary data on top of F; Stream \textit{is} the category of F-coalgebras (up to the adequacy and kind conditions).

\textbf{Proposition 6.1.7 (Adequacy is preserved by Stream-morphisms).} \textit{If f : S → S' is a Stream-morphism and S is adequate, then the image of S under f is adequate as a sub-stream of S'.}

\textbf{Proof.} Let (a, b) be distinguishable aspects of f\_\allowbreak{}σ(S) ⊆ σ'. By injectivity of kind-respecting carrier-maps on distinguishable aspects (standard in the ambient category), their pre-images (a', b') are distinguishable aspects of σ. Adequacy of S gives a content-morphism in ContentOp(σ) witnessing (a', b'); its f\_\allowbreak{}C-image in ContentOp(σ') witnesses (a, b). ∎

\sectionbreak

\section{§6.2 --- The Triple functor}

\textbf{Definition 6.2.1 (The Triple target category).} Define

\[
\mathbf{Triple} := \mathbf{Form} \times \mathbf{Content} \times \mathbf{Carrier}
\]

where:

\begin{itemize}
\item \textbf{Form} is the category of bare carriers (objects: carrier sets σ; morphisms: carrier-maps σ → σ' in the ambient category).
\item \textbf{Content} is Cat\_\allowbreak{}small --- objects: small categories; morphisms: functors.
\item \textbf{Carrier} is the category whose objects are coalgebra-structure-maps γ : σ → F(σ) and whose morphisms are coalgebra-commute squares.
\end{itemize}

\textbf{Triple} inherits a product-2-category structure from its three components. When §6 treats Triple strictly as a 1-category, 2-cells are discarded; where 2-cell data is load-bearing (notably §6.2.5 and the middle-regime texture of §6.5), the bicategorical structure is invoked explicitly.

\textbf{Definition 6.2.2 (The Triple functor).} T : \textbf{Stream} → \textbf{Triple} is defined by:

\begin{itemize}
\item On objects: T(σ, ContentOp(σ), γ) = (σ, ContentOp(σ), γ). The three projections are the three aspects --- Form takes σ, Content takes ContentOp(σ), Carrier takes γ.
\item On morphisms: T(f) = (f\_\allowbreak{}σ, f\_\allowbreak{}C, f\_\allowbreak{}γ), where f\_\allowbreak{}γ is the coalgebra-commute-square data of f.
\end{itemize}

\textbf{Proposition 6.2.3 (T is a functor).} T preserves composition and identities.

\textbf{Proof.} Each of the three projections is separately functorial by inspection. Composition in Triple is componentwise (Definition 6.2.1), matching composition in Stream (Proposition 6.1.4). ∎

\textbf{Lemma 6.2.4 (T is forgetful --- F-coalgebra data factors the Triple).} \textit{The data T(S) = (σ, ContentOp(σ), γ) of the Triple is exactly the F-coalgebra data of the Stream-object S. No additional stream-data is required to reconstruct T(S); no Stream-data is lost in T(S) except derived components (K, Ω).}

\textbf{Proof.} By Definition 6.1.1, a Stream-object is a triple (σ, ContentOp(σ), γ). The Triple T(S) retrieves these three as the Form, Content, Carrier components. K and Ω are derived (Remark 6.1.2), so their absence from T(S) is not a data-loss in the categorical sense. The reconstruction Stream → Triple is the identity map on underlying data. ∎

\textbf{Corollary 6.2.5 (Triple is the object-data of Stream).} The objects of Stream are in canonical bijection with the objects of the image T(Stream) ⊆ Triple satisfying the adequacy and F-coalgebra conditions. Stream is the full subcategory of Triple cut out by:

(a) ContentOp(σ) is adequate for σ (Convention 6.0.3).

(b) γ : σ → F(σ) is an F-coalgebra (satisfies the coalgebra identity up to the choice of F).

(c) Morphisms respect kind (Convention 6.0.5) and coalgebra-commute (Definition 6.1.3 (iii)).

\textbf{Remark 6.2.6.} The Triple is sometimes summarized paired-prose-style as "Form / Content / Carrier as orthogonal-but-constrained axes with recursive decomposability." The formal content of this summary is:

\begin{itemize}
\item \textbf{Orthogonal:} T factors Stream-data into three independent projections (Definition 6.2.2).
\item \textbf{Constrained:} the three projections are not independent --- the coalgebra-commute condition (Definition 6.1.3 (iii)) couples them at the morphism level.
\item \textbf{Recursive decomposability:} each projection-image is itself a stream (§6.3).
\end{itemize}

\textbf{Proposition 6.2.7 (T reflects isomorphisms).} \textit{If f : S → S' is a Stream-morphism whose three Triple-components T(f) = (f\_\allowbreak{}σ, f\_\allowbreak{}C, f\_\allowbreak{}γ) are all isomorphisms in their respective component categories, then f is an isomorphism in Stream.}

\textbf{Proof.} Construct f\textasciicircum{}{-1} componentwise: ((f\_\allowbreak{}σ)\textasciicircum{}{-1}, (f\_\allowbreak{}C)\textasciicircum{}{-1}, (f\_\allowbreak{}γ)\textasciicircum{}{-1}). The coalgebra-commute square for f\textasciicircum{}{-1} is obtained by inverting each side of f's square, which is well-defined because all three components are invertible. Kind respect holds in both directions because a preorder isomorphism is two-way. ∎

\textbf{Remark 6.2.8.} Proposition 6.2.7 is the Triple's \textbf{conservativity} statement: an isomorphism of Stream-objects is exactly an isomorphism of all three Triple-components. This is the formal content of the Anchor's "the Triple is the load-bearing identity-structure of a stream."

\sectionbreak

\section{§6.3 --- Recursive decomposability (finite depth)}

\textbf{Setup.} Each Triple-component of a Stream-object is itself a stream. Explicitly, for S = (σ, ContentOp(σ), γ):

\begin{itemize}
\item \textbf{Form(S)} --- the carrier σ equipped with its form-induced ContentOp (the full subcategory of ContentOp(σ) consisting of content-operations acting on σ's surface structure alone) and restricted γ.
\item \textbf{Content(S)} --- the category ContentOp(σ) viewed as a stream, with its own ContentOp (the 2-category of functorial operations on ContentOp(σ), small by Convention 6.0.1) and its own coalgebra-structure (the hom-2-functor).
\item \textbf{Carrier(S)} --- the coalgebra γ viewed as a morphism-stream, with ContentOp(γ) the 2-cell-level content-operations and coalgebra-structure the squares-of-squares.
\end{itemize}

Each construction yields a Stream-object in its own right.

\textbf{Definition 6.3.1 (Iterated Triple).} Define T\textasciicircum{}(n) : Stream → Triple\textasciicircum{}(n) inductively:

\[
\begin{aligned}
T^{(1)}(S) & = T(S) = (\mathrm{Form}(S), \mathrm{Content}(S), \mathrm{Carrier}(S)) \\
T^{(n+1)}(S) & = \bigl(T^{(1)}(\mathrm{Form}(S)),\ T^{(1)}(\mathrm{Content}(S)),\ T^{(1)}(\mathrm{Carrier}(S))\bigr)
\end{aligned}
\]

The target Triple\textasciicircum{}(n) is a 3\textasciicircum{}n-fold product of Triple.

\textbf{Lemma 6.3.2 (Finite-depth Triple-factorability via F).} \textit{For every Stream-object S and every finite n ≥ 1, T\textasciicircum{}(n)(S) is reconstructible from the F-coalgebra data of S together with iterated ContentOp-applications at each level.}

\textbf{Proof.} Induction on n.

\textit{Base (n = 1).} By Lemma 6.2.4, T(S) = T\textasciicircum{}(1)(S) is the F-coalgebra data of S.

\textit{Step (n → n+1).} Assume T\textasciicircum{}(n)(S') is F-reconstructible for all streams S' up to depth n. For T\textasciicircum{}(n+1)(S):

\begin{itemize}
\item \textbf{Form-branch.} Form(S) is a stream (by the preceding setup), and its F-data (σ, ContentOp(Form(S)), γ|\_\allowbreak{}Form) is a restriction of S's F-data, well-defined by the adequacy-stability argument (Proposition 6.1.7 applied to the identity-on-σ map from S to Form(S)). By inductive hypothesis on Form(S) at depth n, T\textasciicircum{}(n)(Form(S)) is F-reconstructible. Hence T\textasciicircum{}(1)(T\textasciicircum{}(n)(Form(S))) is the first branch of T\textasciicircum{}(n+1)(S), and F-reconstructible.

\item \textbf{Content-branch.} Content(S) = ContentOp(σ) is a small category, viewable as a stream whose ContentOp is the 2-category of functors on ContentOp(σ). This construction uses Convention 6.0.1's small-ness crucially --- the 2-category of endofunctors on a small category is itself small. By inductive hypothesis, T\textasciicircum{}(n)(Content(S)) is F-reconstructible. Branch two of T\textasciicircum{}(n+1)(S) follows.

\item \textbf{Carrier-branch.} Carrier(S) = γ, viewable as a morphism-stream whose ContentOp consists of 2-cells in the ambient 2-category and whose coalgebra-structure is the square-of-squares. By inductive hypothesis on this morphism-stream at depth n, branch three follows.
\end{itemize}

Assembling the three branches gives T\textasciicircum{}(n+1)(S) as F-reconstructible. ∎

\textbf{Remark 6.3.3 (The Cat-valued promotion earns its keep here).} The Content-branch is the load-bearing step. If ContentOp were set-valued (as in the foundation-doc pre-M12 proposal), the Content-branch would terminate at n = 1 --- Content(S) would be a bare set with no further ContentOp to apply T to. Category-valued ContentOp provides the recursion substrate. This is a structural --- not merely technical --- fact about the framework.

\textbf{Corollary 6.3.4 (Recursive decomposability is a theorem, not an axiom).} \textit{The recursive-decomposability content of the Identity-Trajectory Triple (named in the Anchor as a property of streams, and captured as meta-bridge M3 in the framework's bridge list) is derived from the F-coalgebra definition of Stream. It is not an independent axiom.}

\textbf{Proof.} Immediate from Lemma 6.3.2: at every finite depth, T\textasciicircum{}(n) is defined and F-reconstructible. The framework's original axiom-posture for recursive decomposability was pedagogical; the formal content is Lemma 6.3.2. ∎

\textbf{Remark 6.3.5.} Depth-ω (infinite recursive decomposition) is not automatic. It requires the final F-coalgebra construction, which in turn requires filtered-colimit preservation by F. This is treated in §6.9 with its own hypothesis-set.

\sectionbreak

\section{§6.4 --- The kind-classifier fibration over Content}

\textbf{Motivation.} The Anchor's A2 posits kind-stratification of streams (reactive/self-maint/self-ref/abstr). The formal content of the stratification is that there is a classifier-functor π from streams to a ContentOp-class-preorder, and the A2 sub-categories of Stream are fibers of π. This section constructs π and records its fibration properties.

\textbf{Definition 6.4.1 (ContentIndex).} The category \textbf{ContentIndex} has:

\begin{itemize}
\item \textbf{Objects:} equivalence classes of small categories under equivalence-of-categories, restricted to classes realized by some stream's ContentOp.
\item \textbf{Morphisms:} classes of faithful functors (content-refinements) up to natural isomorphism.
\end{itemize}

In the generic framework setting, ContentIndex is a \textbf{preorder} (Convention 6.0.5). When ContentOp-structure admits (co)products globally, it is a \textbf{lattice}.

\textbf{Definition 6.4.2 (The kind-classifier).} Define π : \textbf{Stream} → \textbf{ContentIndex} by

\[
\pi(\sigma, \mathrm{ContentOp}(\sigma), \gamma) := [\mathrm{ContentOp}(\sigma)]
\]

--- the class of ContentOp(σ) in ContentIndex. On morphisms, π(f) = [f\_\allowbreak{}C], the induced class-level functor.

\textbf{Proposition 6.4.3 (π is a functor).} Composition and identities preserved. ∎

\textbf{Construction 6.4.4 (Cartesian lifts as restriction content-operations).} Given a Stream-object S = (σ, ContentOp(σ), γ) with π(S) = c, and a morphism f : c' → c in ContentIndex, the \textbf{Cartesian lift} f̄ : S' → S is constructed as follows:

\begin{enumerate}
\item Choose a representative ι : C' ↪ ContentOp(σ) of the class c' (canonical up to equivalence).
\item Define S' := (σ, C', γ|\_\allowbreak{}{C'}), where γ|\_\allowbreak{}{C'} is γ post-composed with the induced restriction σ\textasciicircum{}(ContentOp(σ)\textasciicircum{}op) → σ\textasciicircum{}(C'\textasciicircum{}op).
\item The lift f̄ is the Stream-morphism (id\_\allowbreak{}σ, ι, γ-commute-by-construction).
\end{enumerate}

\textbf{Lemma 6.4.5 (Cartesian lifts are literally content-operations).} \textit{The Cartesian lift f̄ above is itself a content-operation: specifically, the \textbf{restriction} content-operation given by the inclusion functor ι. Under the Cat-valued ContentOp convention, ι ∈ Mor(Cat\_\allowbreak{}small) is a legitimate content-operation on ContentOp(σ).}

\textbf{Proof.} The data of f̄ consists of id\_\allowbreak{}σ on carriers, ι on ContentOp, and an induced γ-commute. All three components are captured by ι as a morphism in Cat\_\allowbreak{}small. Restriction-content-operations are a recognized sub-class of content-operations under Cat-valuation. ∎

\textbf{Theorem 6.4.6 (π is a bicategorical fibration).} \textit{Given a Stream-object S with π(S) = c and a morphism f : c' → c in ContentIndex, the Cartesian lift f̄ : S' → S of Construction 6.4.4 is Cartesian in the bicategorical sense --- i.e., for any Stream-morphism g : T → S and any morphism h : π(T) → c' in ContentIndex with π(g) = f ∘ h, there exists a Stream-morphism h̄ : T → S' unique up to equivalence, with f̄ ∘ h̄ = g and π(h̄) = h.}

\textbf{Proof sketch.} Existence: factor g\_\allowbreak{}C : ContentOp(T) → ContentOp(σ) through the inclusion ι : C' ↪ ContentOp(σ), using π(g) = f ∘ h to ensure the image lies in C' (class-level); this defines h̄. Uniqueness up to equivalence: two factorizations through a faithful inclusion agree up to natural isomorphism on the image; this is the "up to equivalence" weakening. Full strictness holds when ι is identity-on-objects-faithful, which occurs when ContentIndex is a lattice (Convention 6.0.5). ∎

\textbf{Corollary 6.4.7 (Strict Cartesian fibration under lattice).} \textit{If ContentIndex is a lattice (Convention 6.0.5 lattice case), π is a strict Grothendieck fibration.}

\textbf{Remark 6.4.8.} The Anchor's A2 treatment reads at the preorder level. The bicategorical-fibration status of π (Theorem 6.4.6) is what the Anchor prose "A2 kind-stratification" formally denotes. The strict fibration corollary is the lattice refinement.

\subsection{§6.4.1 --- Unifying streams and admissibility}

\textbf{Definition 6.4.9 (Unifying stream).} A Stream-object S = (σ, ContentOp(σ), γ) is \textbf{unifying} iff there exists an object c\_\allowbreak{}⊥ ∈ ContentOp(σ) that is terminal in ContentOp(σ) \textit{and} satisfies σ\textasciicircum{}(c\_\allowbreak{}⊥) ≅ 1 in the ambient category.

\textbf{Equivalent characterizations.} Under Convention 6.0.3 (adequacy) and 6.0.1 (small-ness):

(i) ContentOp(σ) has a terminal object c\_\allowbreak{}⊥ with σ\textasciicircum{}(c\_\allowbreak{}⊥) ≅ 1.

(ii) The partition of σ induced by c\_\allowbreak{}⊥'s output-equivalence is the indiscrete partition (all σ-elements in one class).

(iii) There exists a natural transformation from the constant presheaf 1 to σ\textasciicircum{}(ContentOp(σ)\textasciicircum{}op) picking out c\_\allowbreak{}⊥-as-terminal.

\textbf{Proof of equivalence.} (i) ⇔ (ii): under adequacy, terminal content-operations are exactly those whose induced partition cannot be refined by any other content-operation. The indiscrete partition is the bottom of the partition lattice; terminal content-operations realize it. (i) ⇔ (iii): standard presheaf equivalence for terminal objects. ∎

\textbf{Definition 6.4.10 (Admissible refinement).} A morphism f : c' → c in ContentIndex is \textbf{admissible} iff every representative inclusion ι : C' ↪ ContentOp(σ) preserves the terminal object --- i.e., if c\_\allowbreak{}⊥ ∈ ContentOp(σ) is terminal, then c\_\allowbreak{}⊥ ∈ C' and remains terminal in C'.

\textbf{Lemma 6.4.11 (Unifying-ness lifts along admissible refinements only).} \textit{Let S be unifying with terminal c\_\allowbreak{}⊥ ∈ ContentOp(σ), and let f : c' → c = π(S) be a morphism in ContentIndex. The Cartesian lift S' = (σ, C', γ|\_\allowbreak{}{C'}) is unifying iff f is admissible.}

\textbf{Proof.} (⇐) If f is admissible, c\_\allowbreak{}⊥ ∈ C' remains terminal, and σ\textasciicircum{}(c\_\allowbreak{}⊥) ≅ 1 is preserved since the carrier σ is unchanged. So S' is unifying with the same witness. (⇒) If S' is unifying with witness c'\_\allowbreak{}⊥ ∈ C', then c'\_\allowbreak{}⊥ is terminal in C'. For ι to be faithful with c'\_\allowbreak{}⊥ terminal in C' but possibly not in ContentOp(σ), the inclusion would fail to preserve terminality --- but then ι(c'\_\allowbreak{}⊥) in ContentOp(σ) would not be terminal, contradicting the uniqueness of terminal objects up to isomorphism. So ι must preserve terminality, i.e., f is admissible. ∎

\textbf{Definition 6.4.12 (The fibered subcategory Stream\_\allowbreak{}u).} \textbf{Stream\_\allowbreak{}u} is the full subcategory of Stream whose objects are unifying streams, restricted along the sub-fibration induced by admissible refinements in ContentIndex.

\textbf{Corollary 6.4.13 (Stream\_\allowbreak{}u is fibered over the admissible sub-category).} \textit{The restriction π|\_\allowbreak{}{Stream\_\allowbreak{}u} : Stream\_\allowbreak{}u → ContentIndex\_\allowbreak{}adm is a bicategorical fibration (strict under lattice).}

\textbf{Proof.} Apply Theorem 6.4.6 to the admissible-refinement sub-category and Lemma 6.4.11 for the unifying-ness lift. ∎

\textbf{Remark 6.4.14 (Physics instance --- quantum-mechanical branches as Cartesian lifts).} The fibration π : \textbf{Stream} → \textbf{ContentIndex} has a concrete realization in Lohmiller \& Slotine's multi-valued classical-action construction (\textit{Proc. Roy. Soc. A} 482: 20250413, Thm 2.4 and Thm 3.2). Each branch $j \in J$ of the extremal-action multipath corresponds to a Cartesian lift of the ambient stream over the content-operation that selects branch $j$ (the associated measurement operator). Admissibility in the sense of Definition 6.4.10 --- preservation of a terminal content-operation along the lift --- corresponds to the Lipschitz-continuity conditions in their Theorem 2.4 that guarantee branch existence and local uniqueness via Picard-Lindelöf contraction theory. The unifying-stream condition (Definition 6.4.9) tracks physical configurations where a universal-ground-state branch is well-defined; streams failing unifying-ness correspond to branch-point configurations where the Lipschitz conditions fail (topological non-simply-connectedness, spatial inequality constraints, Hamiltonian singularities --- their $\mathcal{B}^N$ set). This is a scope-limited sufficiency witness --- Lagrangian systems under Coulomb/Lorenz gauge, per the paper's domain --- not a coextensive derivation. It is one instance among possibly many of the fibration's content; the Companion does not claim QM as the canonical case.

\sectionbreak

\section{§6.5 --- Middle-regime morphism-structure (trifurcation)}

\textbf{Motivation.} The Corpus contains a phenomenological finding: the contemplative "ultimate" has internal texture. Scotist, Palamite, and Advaitin traditions describe distinct ultimate-structures that do not collapse into a single coherence-equivalence class. Under the F-coalgebra framework, the texture localizes in the \textbf{morphism-structure of ContentOp at the terminal content-operation}. This section records the formal content.

\textbf{Setup.} Let S be a unifying stream with terminal c\_\allowbreak{}⊥ ∈ ContentOp(σ). The ambient morphism-structure around c\_\allowbreak{}⊥ in ContentOp(σ) --- i.e., the category of morphisms into and out of c\_\allowbreak{}⊥ --- carries additional data beyond c\_\allowbreak{}⊥'s terminality.

\textbf{Definition 6.5.1 (Ultimate-structure of a unifying stream).} The \textbf{ultimate-structure} of a unifying stream S, denoted Ult(S), is the slice category ContentOp(σ) / c\_\allowbreak{}⊥ together with the distinguished-object data (c\_\allowbreak{}⊥ is terminal) and the induced 2-cell structure from Cat\_\allowbreak{}small.

\textbf{Remark 6.5.2.} Ult(S) is small (Convention 6.0.1), has a terminal object (c\_\allowbreak{}⊥), and is non-trivial iff ContentOp(σ) has non-identity morphisms into c\_\allowbreak{}⊥. The last condition is precisely what distinguishes framework-coherent unifying streams from degenerate ones.

\textbf{Definition 6.5.3 (Middle-regime classes).} Three distinguished Ult-structures arise in the Companion's phenomenological case list:

\begin{itemize}
\item \textbf{Scotist class.} Ult(S) has \textbf{internal-compatibility morphisms}: every pair of content-operations c\_\allowbreak{}1, c\_\allowbreak{}2 mapping to c\_\allowbreak{}⊥ admits a common refinement c\_\allowbreak{}{1,2} also mapping to c\_\allowbreak{}⊥, with the refinement structure univocal (single morphism-type to c\_\allowbreak{}⊥ up to natural isomorphism). Formalism: Ult(S) is a finite-product-closed sub-category with products lifting to c\_\allowbreak{}⊥.

\item \textbf{Palamite class.} Ult(S) has \textbf{essence/energy level-morphisms}: a distinguished sub-object c\_\allowbreak{}essence ↪ c\_\allowbreak{}⊥ with the property that morphisms into c\_\allowbreak{}⊥ factor through either c\_\allowbreak{}essence (essence-morphisms) or the complement c\_\allowbreak{}energies (energy-morphisms), with no direct path from energy-level objects to essence. Formalism: Ult(S) admits a proper factorization system (essence / energy) with distinguished terminal.

\item \textbf{Advaitin class.} Ult(S) has \textbf{saguna-to-nirguna projection morphisms}: every content-operation into c\_\allowbreak{}⊥ factors through a unique "with-attribute" intermediate c\_\allowbreak{}saguna, and the composite c → c\_\allowbreak{}saguna → c\_\allowbreak{}⊥ is the same as the direct c → c\_\allowbreak{}⊥ up to a projection natural transformation. Formalism: Ult(S) admits a reflective subcategory structure with c\_\allowbreak{}⊥ as reflector-terminal and c\_\allowbreak{}saguna as reflection-image.
\end{itemize}

\textbf{Theorem 6.5.4 (Three middle-regime classes are framework-distinct).} \textit{The three classes of Definition 6.5.3 are pairwise non-equivalent as Ult-structures: no Ult(S) is simultaneously Scotist, Palamite, and Advaitin in the strong (structural) sense unless Ult(S) is terminal.}

\textbf{Proof sketch.} A strong-sense Scotist structure requires finite products lifting to c\_\allowbreak{}⊥ without intermediate sub-objects; this conflicts with a Palamite essence/energy factorization, which mandates a proper sub-object c\_\allowbreak{}essence ⊊ c\_\allowbreak{}⊥. Similarly, Advaitin reflection requires a non-trivial c\_\allowbreak{}saguna intermediate, incompatible with Scotist univocal direct morphisms. Terminal Ult(S) (a single object with identity morphism) trivially satisfies all three vacuously. ∎

\textbf{Remark 6.5.5 (Why the framework predicts texture).} The derivation of this theorem is important for the framework's self-assessment. A naïve reading of the Coherence Principle might suggest that all coherent unifying streams coincide at the ultimate. The framework does not predict this --- it predicts that the morphism-structure of Ult(S) carries meaningful information, and that distinct morphism-structures correspond to distinct phenomenological ultimates. The Scotist/Palamite/Advaitin texture is a \textit{prediction} of the framework, confirmed by the contemplative-traditions literature, not a failure to collapse.

\textbf{Corollary 6.5.6 (Cross-tradition translation).} \textit{Morphism-structures of different middle-regime classes do not translate into each other except by data-loss. Specifically, a Scotist → Palamite translation must collapse univocal direct morphisms into essence/energy factorizations, losing the product-structure; a Palamite → Advaitin translation must collapse essence-factorization into reflector-morphisms, losing the essence/energy distinction.}

\textbf{Proof.} Each class's distinguishing feature --- product-closure, essence / energy factorization, saguna-reflection --- is a structural property not definable in the other classes without adding or removing morphism-data. Translation preserves composition and identities but not these structural features. ∎

\textbf{Remark 6.5.7 (A48 correspondence, scope).} §6.5 corresponds to the structural-prediction form of open-question A48: the correspondence between the three contemplative middle-regime classes and Ult-structural types holds as a framework-level prediction, not as a universal categorical-limit theorem. A strong universal-limit form is parked for a later pass.

\sectionbreak

\section{§6.6 --- The Triple as colax-limit (conditional)}

\textbf{Motivation.} The Anchor frames the Triple as a \textbf{colax-limit} of a three-pronged diagram in Cat. This is a sharpening of Theorem 6.2.5 (Stream as a full subcategory of Triple), conditional on ContentOp having enough structure.

\textbf{Hypothesis 6.6.0.} \textit{For the results of this section, assume ContentOp(σ) admits an initial object I (Convention 6.0.4 --- carried as hypothesis per-theorem).}

\textbf{Definition 6.6.1 (The colax-limit diagram).} Define a three-object diagram D : J → Cat as follows:

\begin{itemize}
\item J is the walking span-category: three objects {•\_\allowbreak{}F, •\_\allowbreak{}C, •\_\allowbreak{}Cr} and two morphisms •\_\allowbreak{}F → • ← •\_\allowbreak{}C, •\_\allowbreak{}Cr → •, where • is an auxiliary apex object representing the coalgebra-commute condition.
\item D(•\_\allowbreak{}F) = Form, D(•\_\allowbreak{}C) = Content, D(•\_\allowbreak{}Cr) = Carrier, D(•) = Form × Content × Carrier (with the three projections as D's structure).
\end{itemize}

\textbf{Theorem 6.6.2 (Triple as colax-limit, under Hypothesis 6.6.0).} \textit{Under the initial-object hypothesis on ContentOp(σ), Stream embeds into the colax-limit colaxlim D of the diagram above, and the embedding is an equivalence of categories onto the full sub-2-category of adequate F-coalgebras.}

\textbf{Proof sketch.} colaxlim D consists of cones (F\_\allowbreak{}0, C\_\allowbreak{}0, Cr\_\allowbreak{}0, α\_\allowbreak{}F, α\_\allowbreak{}C, α\_\allowbreak{}Cr) where F\_\allowbreak{}0, C\_\allowbreak{}0, Cr\_\allowbreak{}0 are objects in Form, Content, Carrier respectively, and α\_\allowbreak{}∗ are colax-natural-transformations to the apex. A Stream-object (σ, ContentOp(σ), γ) supplies (σ, ContentOp(σ), γ) as the three objects, and the coalgebra-commute condition provides the colax-natural-transformation data with \textit{initial-object-anchored} structure: the initial I ∈ ContentOp(σ) serves as the tip of the colax cone, giving a canonical base-point from which all content-operations flow.

The embedding Stream ↪ colaxlim D is fully faithful by Proposition 6.2.7 (T reflects isomorphisms). Essential surjectivity onto adequate F-coalgebras holds because the colax structure precisely encodes the coalgebra-commute condition. ∎

\textbf{Remark 6.6.3 (Why this requires initial objects).} Without an initial object in ContentOp(σ), the colax cone has no anchor-point, and the colax-limit construction produces a larger object than Stream --- one that includes "ungrounded" Triple-tuples without a canonical base content-operation. The initial-object hypothesis makes the cone well-pointed and the colax-limit equivalent to Stream.

\textbf{Remark 6.6.4 (General case without initial objects).} For streams without initial ContentOp, Theorem 6.2.5 (Stream as a full subcategory of Triple) is the correct structural description. The colax-limit form is a sharpening available under the additional hypothesis. Both are drafting-valid; the paired-prose Anchor uses the colax-limit form wherever it is applicable, with the un-conditioned statement as fallback.

\textbf{Corollary 6.6.5 (The Triple's "canonical base" is the initial content-operation).} \textit{When Hypothesis 6.6.0 holds, the initial object I ∈ ContentOp(σ) is the distinguished "ground" content-operation from which all stream-internal content-operations are reachable. This formalizes the paired-prose notion of a stream's "ground state" or "base coherence."}

\textbf{Proof.} Immediate from the colax-limit anchor-point construction in Theorem 6.6.2. ∎

\sectionbreak

\section{§6.7 --- Stream as the category of F-coalgebras (closure)}

\textbf{Motivation.} Theorem 6.2.5 stated that Stream is a full subcategory of Triple cut out by three conditions (adequacy, F-coalgebra identity, kind-respecting morphisms). Theorem 6.6.2 refined this (under a hypothesis) to a colax-limit equivalence. This section closes the structural characterization by stating the main equivalence.

\textbf{Theorem 6.7.1 (Stream = F-Coalg\_\allowbreak{}ad).} \textit{Let F-Coalg denote the category of F-coalgebras in the ambient Set-like category, and let F-Coalg\_\allowbreak{}ad ⊆ F-Coalg be the full subcategory of adequate F-coalgebras satisfying kind-respect on morphisms. Then:}

\[
\mathbf{Stream} \simeq \mathbf{F\text{-}Coalg}_{\mathrm{ad}}
\]

\textit{as categories.}

\textbf{Proof.} By Definition 6.1.1, every Stream-object is an F-coalgebra with adequacy and kind-data. Conversely, every F-coalgebra γ : σ → F(σ) with adequate ContentOp(σ) lifts to a Stream-object by setting K = [ContentOp(σ)] and Ω = F(σ) (Remark 6.1.2). The lift is bijective on objects and morphisms by Definition 6.1.3. Functoriality of the equivalence holds because composition in both categories is componentwise (Proposition 6.1.4). ∎

\textbf{Corollary 6.7.2 (Closure).} \textit{The category Stream is closed under the F-coalgebra structure --- every construction producing an F-coalgebra on some σ with adequate ContentOp produces a Stream-object.}

\textbf{Proof.} Immediate from Theorem 6.7.1. ∎

\textbf{Remark 6.7.3 (Implication for framework construction).} Closure under F-coalgebra operations means that Stream is algebraically well-behaved under coalgebraic constructions --- products (§6.8), equalizers (§6.8), filtered colimits (§6.8), and the self-reference closure of §8 all produce Stream-objects when their inputs are Stream-objects. The framework does not need to re-check adequacy and kind-respect at each construction; these properties transfer automatically under Stream-morphisms (Proposition 6.1.7).

\sectionbreak

\section{§6.8 --- Limits and colimits in Stream under F}

\textbf{Conventions recalled.} For this section, size/variance/adequacy follow §6.0. Kind-structure is the preorder of Convention 6.0.5; lattice-level results are flagged explicitly.

\subsection{§6.8.1 --- Limits}

\textbf{Proposition 6.8.1 (Terminal object).} \textit{The terminal object 1\_\allowbreak{}Stream := (1, \textbf{1}\_\allowbreak{}cat, id\_\allowbreak{}1) exists in Stream, where 1 is the terminal carrier, \textbf{1}\_\allowbreak{}cat is the terminal small category (one object, identity morphism), and id\_\allowbreak{}1 is the unique coalgebra-structure-map.}

\textbf{Proof.} Uniqueness of maps into 1\_\allowbreak{}Stream: f\_\allowbreak{}σ factors uniquely through 1; f\_\allowbreak{}C factors uniquely through \textbf{1}\_\allowbreak{}cat; coalgebra-commute is vacuous because γ\_\allowbreak{}{1\_\allowbreak{}Stream} is identity. Kind-respect is trivial because K(1\_\allowbreak{}Stream) is the top of the preorder. ∎

\textbf{Proposition 6.8.2 (Products, conditional on kind-join).} \textit{Given Stream-objects S\_\allowbreak{}1, S\_\allowbreak{}2 with K(S\_\allowbreak{}1), K(S\_\allowbreak{}2) admitting a join K\_\allowbreak{}∨ in the A2 preorder, the product S\_\allowbreak{}1 × S\_\allowbreak{}2 exists in Stream with:}

\[
S_1 \times S_2 = (\sigma_1 \times \sigma_2,\ \mathrm{ContentOp}(S_1) \times \mathrm{ContentOp}(S_2),\ \gamma_1 \times \gamma_2)
\]

\textit{with the cartesian product of coalgebras and kind K\_\allowbreak{}∨.}

\textbf{Proof sketch.} Universal property: a pair (g\_\allowbreak{}1, g\_\allowbreak{}2) : T → S\_\allowbreak{}1, T → S\_\allowbreak{}2 induces a unique (g\_\allowbreak{}1, g\_\allowbreak{}2)\_\allowbreak{}* : T → S\_\allowbreak{}1 × S\_\allowbreak{}2 by pairing carrier-maps and ContentOp-functors componentwise. Coalgebra-commute transfers by pairing. Kind-respect K(T) ⊑ K\_\allowbreak{}∨ requires K\_\allowbreak{}∨ to exist as a join --- hence the conditional. ∎

\textbf{Proposition 6.8.3 (Equalizers, conditional on adequacy-stability).} \textit{Given parallel Stream-morphisms f, g : S → S', the equalizer exists in Stream provided the carrier-equalizer eq(f\_\allowbreak{}σ, g\_\allowbreak{}σ) carries an adequate ContentOp-restriction.}

\textbf{Proof sketch.} Carrier: eq(f\_\allowbreak{}σ, g\_\allowbreak{}σ) ⊆ σ. ContentOp: the sub-category of ContentOp(σ) consisting of content-operations that agree after mapping by f\_\allowbreak{}C and g\_\allowbreak{}C. Coalgebra: restricted γ. Adequacy-stability (Lemma 6.8.β below) gives the conditional. ∎

\textbf{Lemma 6.8.β (Adequacy-stability under limits/colimits).} \textit{Limits and colimits in Stream preserve adequacy, provided each constituent stream is adequate and the limit/colimit operation respects ContentOp-morphism-witnessed distinctions.}

\textbf{Proof.} For a limit L, distinguishable aspects of L descend to distinguishable aspects in at least one constituent (by universality); ContentOp-morphism witnesses lift to L via the limit-cone. Dual argument for colimits via the couniversal property. ∎

\textbf{Proposition 6.8.4 (Filtered limits under locally-presentable promotion).} \textit{When ContentOp is promoted to κ-accessible locally-presentable (Convention 6.0.1), filtered limits exist in Stream.}

\subsection{§6.8.2 --- Colimits}

\textbf{Proposition 6.8.5 (Initial object).} \textit{The initial object 0\_\allowbreak{}Stream := (∅, \textbf{0}\_\allowbreak{}cat, !) exists, where ∅ is the empty carrier, \textbf{0}\_\allowbreak{}cat is the empty category, and ! is the unique map from ∅.}

\textbf{Proof.} Vacuous satisfaction of all Stream conditions. ∎

\textbf{Proposition 6.8.6 (Coproducts, conditional on kind-meet).} \textit{Given S\_\allowbreak{}1, S\_\allowbreak{}2 with K(S\_\allowbreak{}1), K(S\_\allowbreak{}2) admitting a meet K\_\allowbreak{}∧, the coproduct S\_\allowbreak{}1 + S\_\allowbreak{}2 exists with disjoint-union carriers, disjoint-union ContentOp, and componentwise coalgebra.}

\textbf{Proposition 6.8.7 (Coequalizers, conditional on quotient-adequacy).} \textit{Given f, g : S → S', the coequalizer exists provided the quotient ContentOp --- consisting of content-operations that respect the generated equivalence --- is adequate for σ'/∼.}

\textbf{Proposition 6.8.8 (Filtered colimits under locally-presentable promotion).} \textit{Under Convention 6.0.1's locally-presentable promotion, filtered colimits exist in Stream.}

\subsection{§6.8.3 --- What Stream \textit{does not} have}

\textbf{Proposition 6.8.9 (No exponentials in general).} \textit{Stream is not Cartesian closed; there is no canonical (S')\textasciicircum{}S object for arbitrary S, S' ∈ Stream.}

\textbf{Proof sketch.} Standard counterexample from coalgebra theory: exponentials of coalgebras require specific smallness and commutativity conditions not implied by Stream's definition. The F-coalgebra structure fixes γ-data that does not naturally lift to a function-space. ∎

\textbf{Proposition 6.8.10 (No subobject classifier).} \textit{Stream does not have a subobject-classifier object Ω\_\allowbreak{}sub in the topos sense.}

\textbf{Remark 6.8.11 (The framework's "Ω" is not Ω\_\allowbreak{}sub).} The framework's configuration space Ω for a stream (named in earlier expositions as the 4-tuple's third component, derived in Remark 6.1.2 as σ\textasciicircum{}(ContentOp(σ)\textasciicircum{}op)) is not a subobject classifier; it is the evaluation-space of the F-structure. The names coincide accidentally; no toposic interpretation is intended.

\subsection{§6.8.4 --- Summary table}

\begingroup\setlength{\tabcolsep}{3pt}\renewcommand{\arraystretch}{1.15}\footnotesize
\begin{longtable}{p{1.19in} p{0.65in} p{2.49in}}
\toprule
Limit/colimit & Existence & Condition \\
\midrule
\endhead
\bottomrule
\endfoot
Terminal 1\_\allowbreak{}Stream & yes & --- \\
\midrule
Initial 0\_\allowbreak{}Stream & yes & --- \\
\midrule
Products S\_\allowbreak{}1 × S\_\allowbreak{}2 & conditional & kind-join exists (Q7 lattice case) \\
\midrule
Coproducts S\_\allowbreak{}1 + S\_\allowbreak{}2 & conditional & kind-meet exists (Q7 lattice case) \\
\midrule
Equalizers & conditional & adequacy-stability (Lem 6.8.β) \\
\midrule
Coequalizers & conditional & quotient-adequacy \\
\midrule
Pullbacks & conditional & products + equalizers \\
\midrule
Pushouts & conditional & coproducts + coequalizers \\
\midrule
Filtered limits & yes & locally-presentable promotion (Conv 6.0.1) \\
\midrule
Filtered colimits & yes & locally-presentable promotion (Conv 6.0.1) \\
\midrule
Exponentials & no & --- \\
\midrule
Subobject classifier & no & --- \\
\bottomrule
\end{longtable}
\endgroup
\sectionbreak

\section{§6.9 --- C-size regimes, H1+H2, and recursive decomposability at depth ω}

\textit{This section front-loads the \textbf{size-of-ContentOp} as the single parameter governing H1 (accessibility of Stream) and H2 (filtered-colimit preservation of F), partitions Stream into three regimes, proves H1+H2 in the regime where they hold, and places F\_\allowbreak{}∞ (§8.1.2) at the regime boundary.}

\subsection{§6.9.0 --- The three C-size regimes}

For a stream S = (σ, C, γ), the \textbf{size regime} of S is determined by the cardinality-and-presentability structure of C = ContentOp(σ). Three regimes are operationally meaningful, and the framework's theorems behave distinctly in each:

\begingroup\setlength{\tabcolsep}{3pt}\renewcommand{\arraystretch}{1.15}\footnotesize
\begin{longtable}{p{0.59in} p{1.05in} p{1.15in} p{0.63in} p{0.79in}}
\toprule
Regime & C-structure & Measure infra (§7) & H1 accessibility & H2 filtered-colimit preservation \\
\midrule
\endhead
\bottomrule
\endfoot
\textbf{A --- finite-C} & C has finitely many objects and morphisms & trivially σ-finite & trivially holds & \textbf{holds} (Prop 6.9.1 below) \\
\midrule
\textbf{B --- small-but-infinite-C} & C is small (set-many), admits infinitely many objects & σ-finite under countable-C with concrete reference measure & holds under local-presentability & \textbf{generically fails} (Prop 6.9.2 below) \\
\midrule
\textbf{C --- large-C} & C is a proper class & σ-finiteness breaks & fails & fails \\
\bottomrule
\end{longtable}
\endgroup
\textbf{Declared scope.} The Companion's declared scope is Regimes A and B. Regime C is excluded by Convention 1.1.5 (smallness of ContentOp).

\textbf{Why C-size is the governing parameter.} F(σ) = σ\textasciicircum{}(ContentOp(σ)\textasciicircum{}op) is an exponential whose exponent-argument is ContentOp. Exponentials σ ↦ σ\textasciicircum{}X in a locally-presentable category preserve filtered colimits in σ \textbf{iff X is finitely-presentable} (Adámek–Rosický 1994, Thm 1.56 + Cor 1.57). Hence the H2 hypothesis reduces to the finite-presentability of C\textasciicircum{}op --- which in turn reduces to the finite-generation of C. This single fact ties H1, H2, §7 σ-finiteness, and §8 depth-stability to a single parameter.

\subsection{§6.9.1 --- Regime A: H1 and H2 both hold}

\textbf{Proposition 6.9.1 (H1+H2 under finite-C).} \textit{Let 𝒞\_\allowbreak{}Streams\textasciicircum{}{fin} ⊂ 𝒞\_\allowbreak{}Streams be the full subcategory of streams S with C\_\allowbreak{}S finite. Then on 𝒞\_\allowbreak{}Streams\textasciicircum{}{fin}:}

\begin{enumerate}
\item \textit{(H1) 𝒞\_\allowbreak{}Streams\textasciicircum{}{fin} is ℵ\_\allowbreak{}0-accessible locally presentable.}
\item \textit{(H2) F restricted to 𝒞\_\allowbreak{}Streams\textasciicircum{}{fin} preserves filtered colimits.}
\end{enumerate}

\textbf{Proof.}

\textit{(H1):} 𝒞\_\allowbreak{}Streams\textasciicircum{}{fin} consists of pairs (σ, C, γ) with C finite. For each finite shape C, the category F-Coalg\_\allowbreak{}{ad}(F\_\allowbreak{}C) of F\_\allowbreak{}C-coalgebras over a fixed finite C is a slice of Cat over a finite base, which is accessible (Adámek–Rosický 1994, Thm 2.78). The coproduct over isomorphism-classes of finite C is countable, and a countable coproduct of accessible categories is accessible. Hence 𝒞\_\allowbreak{}Streams\textasciicircum{}{fin} = ∐\_\allowbreak{}{[C] finite} F-Coalg\_\allowbreak{}{ad}(F\_\allowbreak{}C) is accessible. Local presentability: 𝒞\_\allowbreak{}Streams\textasciicircum{}{fin} is cocomplete (Propositions 6.8.5–6.8.7 specialized to finite C) and every object is a filtered colimit of ℵ\_\allowbreak{}0-presentable ones (every finite-C stream is itself ℵ\_\allowbreak{}0-presentable). ∎

\textit{(H2):} For fixed finite C, F\_\allowbreak{}C(σ) = σ\textasciicircum{}(C\textasciicircum{}op) is a finite product of copies of σ indexed by objects of C\textasciicircum{}op. Finite products preserve filtered colimits in any category that has them (standard; e.g., Borceux 1994, Prop 2.13.4). Hence F\_\allowbreak{}C preserves filtered colimits in σ. For varying C, the functor F = (σ, C) ↦ σ\textasciicircum{}(C\textasciicircum{}op) has the component F\_\allowbreak{}C preserving filtered colimits pointwise; combined with (H1)'s slicewise accessibility and the fact that filtered colimits in 𝒞\_\allowbreak{}Streams\textasciicircum{}{fin} are computed component-wise in each C-slice (Proposition 6.8.8 specialized to finite C), F on 𝒞\_\allowbreak{}Streams\textasciicircum{}{fin} preserves filtered colimits. ∎

\textbf{Remark 6.9.1'.} The proof of H2 in Regime A is the \textit{content} of the claim. Without the finite-C restriction, σ\textasciicircum{}(C\textasciicircum{}op) is an infinite product (over objects of C\textasciicircum{}op), and infinite products generically do not preserve filtered colimits in the product argument. The shift from "finite product ⇒ preserves filtered colimits" to "infinite product ⇒ does not in general" is exactly the Regime A ↔ Regime B boundary.

\subsection{§6.9.2 --- Regime B: H2 generically fails}

\textbf{Proposition 6.9.2 (H2 failure in Regime B).} \textit{Let S = (σ, C, γ) be a stream with C small and containing infinitely many objects no finite sub-family of which is cofinal in C. Then F(σ) = σ\textasciicircum{}(C\textasciicircum{}op) does not in general preserve filtered colimits of σ.}

\textbf{Proof.} Suppose (σ\_\allowbreak{}i)\_\allowbreak{}{i ∈ I} is a filtered diagram in Set (or in the ambient base) with colimit σ\_\allowbreak{}∞ = colim\_\allowbreak{}i σ\_\allowbreak{}i. We compute:

\begin{itemize}
\item \textit{(colim\_\allowbreak{}i σ\_\allowbreak{}i)\textasciicircum{}(C\textasciicircum{}op)} = functors C\textasciicircum{}op → σ\_\allowbreak{}∞.
\item \textit{colim\_\allowbreak{}i (σ\_\allowbreak{}i\textasciicircum{}(C\textasciicircum{}op))} = colim\_\allowbreak{}i (functors C\textasciicircum{}op → σ\_\allowbreak{}i).
\end{itemize}

For infinite C with no finite cofinal sub-family, the colimit-exchange

$$\mathrm{colim}_i (\sigma_i^{C^{\mathrm{op}}}) \xrightarrow{?} (\mathrm{colim}_i \sigma_i)^{C^{\mathrm{op}}}$$

is not in general an isomorphism: a functor C\textasciicircum{}op → σ\_\allowbreak{}∞ may send infinitely many distinct objects of C\textasciicircum{}op to distinct values in σ\_\allowbreak{}∞ that cannot all be "found in some σ\_\allowbreak{}i" simultaneously, even though each individual value appears in some σ\_\allowbreak{}i. Explicit counterexample: take C\textasciicircum{}op = ℕ (discrete), σ\_\allowbreak{}i = {0, 1, ..., i}, σ\_\allowbreak{}∞ = ℕ. The identity function ℕ → ℕ is an element of σ\_\allowbreak{}∞\textasciicircum{}(C\textasciicircum{}op) that is not in any σ\_\allowbreak{}i\textasciicircum{}(C\textasciicircum{}op). Hence F does not preserve this filtered colimit. ∎

\textbf{Consequence.} In Regime B, the transfinite iteration σ\_\allowbreak{}{α+1} = F(σ\_\allowbreak{}α) of Theorem 6.9.3 below does not in general stabilize under H2; the depth-ω result is therefore conditional on additional structure on C (e.g., C finitely-generated, C cofinally finite, or a targeted property of γ).

\subsection{§6.9.3 --- Depth-ω result (conditional on regime)}

\textbf{Theorem 6.9.3 (Final F-coalgebra, regime-specific).} \textit{Within 𝒞\_\allowbreak{}Streams\textasciicircum{}{fin} (Regime A), a final F-coalgebra σ\_\allowbreak{}∞ exists, with σ\_\allowbreak{}∞ ≅ F(σ\_\allowbreak{}∞), and every Stream-object in 𝒞\_\allowbreak{}Streams\textasciicircum{}{fin} maps uniquely into σ\_\allowbreak{}∞.}

\textbf{Proof.} By Prop 6.9.1, H1 and H2 both hold on 𝒞\_\allowbreak{}Streams\textasciicircum{}{fin}. Standard terminal-coalgebra construction (Adámek 1974, Barr 1993): iterate F transfinitely from 1:

\begin{itemize}
\item σ\_\allowbreak{}0 := 1 (terminal object, which exists by Lemma 6.8.α in Regime A),
\item σ\_\allowbreak{}{α+1} := F(σ\_\allowbreak{}α),
\item σ\_\allowbreak{}λ := lim\_\allowbreak{}{α<λ} σ\_\allowbreak{}α at limit ordinals λ.
\end{itemize}

H1 ensures the iteration stabilizes at some ordinal ≤ ℵ\_\allowbreak{}1 (first uncountable regular). H2 ensures F commutes with the stabilizing filtered colimit, so σ\_\allowbreak{}∞ ≅ F(σ\_\allowbreak{}∞). Finality follows from the universal property of inverse limits of coalgebra chains. ∎

\textbf{Corollary 6.9.4 (Depth-ω Triple-factorability in Regime A).} \textit{Within 𝒞\_\allowbreak{}Streams\textasciicircum{}{fin}, T\textasciicircum{}(ω)(σ\_\allowbreak{}∞) is well-defined and fixed: T(σ\_\allowbreak{}∞) ≅ σ\_\allowbreak{}∞ as Triple-objects.}

\textbf{Theorem 6.9.5 (Regime-B depth-ω conditional).} \textit{For a stream S in Regime B, a final F-coalgebra over S exists iff C\_\allowbreak{}S is finitely-generated (equivalently: C\_\allowbreak{}S admits a finite cofinal sub-category). In that case the Regime-A construction applies to the finite-cofinal sub-diagram.}

\textbf{Proof.} If C\_\allowbreak{}S has a finite cofinal sub-category C'\_\allowbreak{}S, then σ\textasciicircum{}(C\_\allowbreak{}S\textasciicircum{}op) ≅ σ\textasciicircum{}(C'\_\allowbreak{}S)\textasciicircum{}op (cofinality of limits), and the finite-C case applies. Conversely, if no finite cofinal sub-category exists, Prop 6.9.2's counterexample instantiates against S, blocking H2. ∎

\subsection{§6.9.4 --- Placement of F\_\allowbreak{}∞}

\textbf{Proposition 6.9.6 (F\_\allowbreak{}∞'s regime trajectory).} \textit{Per §8.1.2, C\_\allowbreak{}{F, t} is finite at every fixed construction-time t. Hence F\_\allowbreak{}∞ |\_\allowbreak{}{t} ∈ 𝒞\_\allowbreak{}Streams\textasciicircum{}{fin} at every t; F\_\allowbreak{}∞ |\_\allowbreak{}{t} is in Regime A. The colimit C\_\allowbreak{}{F, ∞} := colim\_\allowbreak{}t C\_\allowbreak{}{F, t} as t → ∞ is at most countable and generically not finitely-generated; hence F\_\allowbreak{}∞ over [t\_\allowbreak{}0, ∞) is in Regime B with undetermined H2 status.}

\textbf{Proof.} Finite-at-each-t from §8.1.2 (the commit-history at time t carries finitely many substrate-commitments). Countability of the limit from the commit-history being a countable sequence of snapshots. Finite-generation of C\_\allowbreak{}{F, ∞} fails generically because each substrate-commitment added over time adds new content-operations not derivable from previously-present ones (which is exactly the C12 autocatalysis corollary content). ∎

\textbf{Consequence.} Audit Observation 8.3.5's self-reference claim over a \textit{finite} construction interval [t\_\allowbreak{}0, t\_\allowbreak{}1] lands in Regime A (via the finite-slice 𝒞\_\allowbreak{}Streams\textasciicircum{}{fin}); the H2-hypothesis of the depth-ω theorem is not invoked. This is why the §8 finite-interval audit is well-founded even without H2 verification at the colimit. The \textit{infinite-interval} self-reference claim --- "Principle-about-itself for the construction process extended indefinitely" --- does invoke Theorem 6.9.5 at the colimit and is conditional on C\_\allowbreak{}{F, ∞} being finitely-generated (generically false).

\textbf{Scope remark.} The Coherence Principle's self-reference closure is a \textbf{finite-interval} claim by structural necessity. Extending it to ω requires more than just time-passage; it requires that the autocatalytic content-operation discovery process (§C12) halt in the finite-cofinal sense. The framework does not predict such halting, nor does it require it --- the Principle's empirical content is per-interval.

\subsection{§6.9.5 --- Connection to §7's small-C and σ-finiteness}

\textbf{Proposition 6.9.7 (§7 σ-finiteness scope).} \textit{The §7.3.2 σ-finiteness hypothesis for Bias(S) holds:}

\begin{itemize}
\item \textit{Unconditionally in Regime A.}
\item \textit{In Regime B iff C\_\allowbreak{}S is countable AND a countable concrete reference measure μ\_\allowbreak{}0 on Ω\_\allowbreak{}S is fixed.}
\item \textit{Fails in Regime C.}
\end{itemize}

\textbf{Proof.}

\textit{Regime A:} Ω\_\allowbreak{}S = σ\textasciicircum{}(C\textasciicircum{}op) with finite C is a finite product of finite-or-countable copies of σ; any finite-or-countable σ-algebra on σ lifts to a σ-finite reference measure on Ω\_\allowbreak{}S. Hence Bias(S) decomposes as the difference of two finite-or-σ-finite measures (Hahn–Jordan) and is σ-finite.

\textit{Regime B:} Ω\_\allowbreak{}S is a (countable-indexed) product over objects of C\textasciicircum{}op. A σ-finite reference measure exists iff C is countable (giving a countable product) and a concrete σ-finite measure is fixed on σ (which the framework admits per §7.1.1). Without concreteness of μ\_\allowbreak{}0, the product-measure construction fails at the cylinder-set extension step.

\textit{Regime C:} An uncountable product of non-trivial measures is not σ-finite in general; a proper-class product isn't even a measurable space in the ordinary sense. ∎

\textbf{Consequence (scope declaration for §7).} §7's framework applies unconditionally in Regime A and conditionally in Regime B on countability of C and concreteness of μ\_\allowbreak{}0. This is the scope-condition that was silently load-bearing in §7.3.2's σ-finiteness argument; it is now explicit.

\subsection{§6.9.6 --- Surfaced-lemma register (this section)}

\begin{itemize}
\item ⚑ §6.9.0 C-size-regime partition (A / B / C) → Anchor §1/§3 target --- convention (makes the size-of-C scope-condition explicit across framework)
\item ⚑ Prop 6.9.1 H1+H2 both hold on 𝒞\_\allowbreak{}Streams\textasciicircum{}{fin} → Anchor §9.5 target --- proposition (resolves H1 and the finite-C case of H2)
\item ⚑ Prop 6.9.2 H2 generically fails in Regime B → Anchor §1.5 target --- proposition (explicit counterexample bounds the scope)
\item ⚑ Thm 6.9.3 / Thm 6.9.5 Final F-coalgebra (Regime A unconditional; Regime B conditional on finite-cofinal sub-diagram) → Anchor §1.5 + §9.5 target --- theorem
\item ⚑ Prop 6.9.6 F\_\allowbreak{}∞'s finite-at-each-t / countable-in-limit regime trajectory → Anchor §9.5 target --- proposition (sharpens the §8 audit scope)
\item ⚑ Prop 6.9.7 Regime-B σ-finiteness requires countable C + concrete μ\_\allowbreak{}0 → Anchor Appendix B §B.1 target --- proposition (names the §7 silent hypothesis explicitly)
\end{itemize}

\sectionbreak

\section{§6.10 --- Inner/Outer adjunction}

\textit{Derives the inner/outer adjunction $\iota_S \dashv \omega_S$ that anchor §1.10 states and anchor §3.8 interprets phenomenologically. §6.10 lifts A2.4's per-pair adjunction over the full DAG $\mathrm{Up}(S)$ of A2.6 via Kelly's indexed-adjunction construction, derives Content as the adjunction's hom-profunctor, and formalizes the unit $\eta$ as a G-coalgebra structure map.}

\textbf{Notation note.} §6 uses F for the Stream endofunctor $F(\sigma) = \sigma^{\mathrm{ContentOp}(\sigma)^{\mathrm{op}}}$. To avoid collision, §6.10 uses \textbf{G} for the monad-underlying endofunctor $G := \omega_S \circ \iota_S$ on $\mathbf{Inner}(S)$. The two endofunctors live on different categories and track different data; the name-distinction is notational only.

\subsection{§6.10.1 --- Setup}

Fix a stream $S$, per A1 and A2. Recall:

\textbf{(A2.4, per-level adjunction.)} For any pair $S \subseteq S_q$ with $S_q \in \mathrm{Up}(S)$ --- the DAG of wholes containing $S$ per A2.6 --- there is an adjunction
\[
\iota_{S \subset S_q} \;\dashv\; \kappa_{S \subset S_q}
\quad : \quad \mathrm{Form}(S) \longleftrightarrow \mathrm{Form}(S_q)
\]
with $\iota$ the left-adjoint embedding and $\kappa$ the right-adjoint restriction. This is stated at anchor §2.4 and in the Companion at §2.

\textbf{(A2.6, DAG of wholes.)} $\mathrm{Up}(S)$ is a poset under whole-containment. $\mathrm{Up}(S)$ has no maximum element --- the "totality" term is not an object of the framework.

\textbf{Definition 6.10.1.1 (Inner category at $S$).} Define $\mathbf{Inner}(S)$ as the category whose objects are pairs $(C, \nu)$ with $C$ a Carrier in the sense of §6.0 and $\nu$ a navigation-trajectory as in A2, and whose morphisms preserve $\nu$ up to Form-pattern-isomorphism. $\mathbf{Inner}(S)$ inherits its structural properties from the Carrier axioms (A2.1) and the navigation-trajectory axioms (A2.5).

\textbf{Definition 6.10.1.2 (Outer category at $S$, preview).} Define $\mathbf{Outer}(S)$ as the Grothendieck construction
\[
\mathbf{Outer}(S) \;:=\; \int_{S_q \in \mathrm{Up}(S)} \mathrm{Form}(S_q)
\]
under the $\iota$-transitions, with:

\begin{itemize}
\item \textit{objects:} pairs $(S_q, \phi)$ with $S_q \in \mathrm{Up}(S)$ and $\phi \in \mathrm{Form}(S_q)$;
\item \textit{morphisms} $(S_q, \phi) \to (S_r, \psi)$ for $S_q \subseteq S_r$: pairs $(f, \alpha)$ where $f$ is the DAG-morphism and $\alpha : \iota_{S_q \subset S_r}(\phi) \to \psi$ is a Form($S_r$)-morphism.
\end{itemize}

The cocompleteness of $\mathbf{Outer}(S)$ is proved in §6.10.3.

\textbf{Notation summary for §6.10.}

\begingroup\setlength{\tabcolsep}{3pt}\renewcommand{\arraystretch}{1.15}\footnotesize
\begin{center}
\begin{tabular}{p{1.08in} p{2.65in} p{0.60in}}
\toprule
Symbol & Meaning & First definition \\
\midrule
$\mathrm{Up}(S)$ & DAG of wholes containing $S$ (A2.6) & §6.10.1 \\
$\mathbf{Inner}(S)$ & Category of (Carrier, navigation-trajectory) pairs at $S$ & Def 6.10.1.1 \\
$\mathbf{Outer}(S)$ & Grothendieck construction $\int_{S_q} \mathrm{Form}(S_q)$ under $\iota$ & Def 6.10.1.2 \\
$\iota_S, \omega_S$ & Total inner/\allowbreak outer adjunction functors & §6.10.4 \\
$G := \omega_S \circ \iota_S$ & Monad-underlying endofunctor on $\mathbf{Inner}(S)$ & §6.10.6 \\
$\eta, \epsilon$ & Unit, counit of the $\iota_S \dashv \omega_S$ adjunction & §6.10.4 \\
$\Psi_S$ & Hom-profunctor of the adjunction (Content as bijection) & §6.10.5 \\
\bottomrule
\end{tabular}
\end{center}
\endgroup
\textbf{Scope remark.} §6.10 is \textit{one-level} lifted from §2's per-level (A2.4) statement: the adjunction is assembled over the full $\mathrm{Up}(S)$ rather than stated separately per-pair. The per-level adjunctions are recovered as fibers via the Grothendieck structure of $\mathbf{Outer}(S)$.

\subsection{§6.10.2 --- Lemma 1: DAG-coherence of $\iota$ and $\kappa$}

\textbf{Lemma 6.10.2.1 (Iterated adjunction coherence).} For $S \subseteq S_q \subseteq S_r$ in $\mathrm{Up}(S)$, there are natural isomorphisms
\[
\iota_{S \subset S_r} \;\cong\; \iota_{S_q \subset S_r} \circ \iota_{S \subset S_q} \qquad (\textit{covariant composition})
\]
\[
\kappa_{S \subset S_r} \;\cong\; \kappa_{S \subset S_q} \circ \kappa_{S_q \subset S_r} \qquad (\textit{contravariant composition})
\]
Both isomorphisms are natural in all three streams. For non-comparable $S_q, S_r \in \mathrm{Up}(S)$, no compatibility axiom is required --- the DAG-poset structure of $\mathrm{Up}(S)$ carries no additional coherence cells.

\textbf{Proof.} For the covariant case: $\iota$ is defined from whole-containment embeddings (A2.1). For the chain $S \subseteq S_q \subseteq S_r$, the composite embedding $S \hookrightarrow S_r$ factors canonically through $S_q$, so $\iota_{S \subset S_r}$ and $\iota_{S_q \subset S_r} \circ \iota_{S \subset S_q}$ are both realized by the same underlying Carrier-embedding. Form-structure lifts along the Carrier-embedding by A2.1 functoriality, yielding the natural isomorphism.

For the contravariant case: the composite $\kappa_{S \subset S_r}$ is right-adjoint to $\iota_{S \subset S_r}$. By the covariant case, the composite left-adjoint factors as $\iota_{S_q \subset S_r} \circ \iota_{S \subset S_q}$. By uniqueness of adjoints up to natural isomorphism, the right adjoint factors dually: $\kappa_{S \subset S_r} \cong \kappa_{S \subset S_q} \circ \kappa_{S_q \subset S_r}$ --- the outermost restriction applied first, then the inner one, reflecting the direction-reversal of adjoints under composition.

For naturality: naturality in each of $S, S_q, S_r$ separately follows from the functoriality of the embedding-lift, which preserves commuting squares at each level. Branching (non-comparable $S_q, S_r$ under a common $S$) produces no coherence cell because the DAG carries no non-trivial 2-cells --- the poset structure of $\mathrm{Up}(S)$ is thin.

$\Box$

\textbf{Remark 6.10.2.2 (Why the Grothendieck construction uses $\iota$, not $\kappa$).} The $\iota$-direction composes covariantly, so $\int_{S_q} \mathrm{Form}(S_q)$ under $\iota$-transitions is a \textit{covariant} Grothendieck construction. $\iota$ is a left adjoint and hence cocontinuous, so the fibers' colimits assemble into total-category colimits by the standard Barr–Wells §2.8 / Johnstone B.1.5.8 criterion. The $\kappa$-direction would yield a contravariant (fibered) construction, which supports limits but not colimits --- relevant for later dualizations but not for the cocompleteness argument in §6.10.3.

⚑ [SURFACED | Companion §6.10.2.1 | → Anchor §1.10 target | type: lemma]
--- Iterated adjunction coherence over $\mathrm{Up}(S)$; no extra axiom required at the DAG-poset level.

\subsection{§6.10.3 --- Lemma 2*: Outer as $\iota$-Grothendieck construction is cocomplete}

\textbf{Lemma 6.10.3.1 (Cocompleteness of $\mathbf{Outer}(S)$ under $\iota$-transitions).} The Grothendieck construction
\[
\mathbf{Outer}(S) \;=\; \int_{S_q \in \mathrm{Up}(S)} \mathrm{Form}(S_q)
\]
with transition-maps given by $\iota_{S_q \subset S_r}$ for each $S_q \subseteq S_r$ in $\mathrm{Up}(S)$ is \textbf{cocomplete}. Explicitly: $\mathbf{Outer}(S)$ admits all small colimits, and these are computed level-wise --- the colimit of a diagram in $\mathbf{Outer}(S)$ is a pair $(S_*, \varphi_*)$ where $S_*$ is the colimit of the diagram's projection to $\mathrm{Up}(S)$ and $\varphi_* = \mathrm{colim}\, \iota_{S_q \subset S_*}(\varphi_{S_q})$ in $\mathrm{Form}(S_*)$.

\textbf{Proof.}

\textit{Step 1 (fiberwise cocompleteness).} Each fiber $\mathrm{Form}(S_q)$ is cocomplete by A2.5 and the Form-pattern generation structure established in anchor §1 (cf. Companion §1 / §2). Small colimits in $\mathrm{Form}(S_q)$ exist and are stable under the constructions needed below.

\textit{Step 2 ($\iota$ is cocontinuous).} Each $\iota_{S \subset S_q}$ is the left adjoint of the per-level A2.4 adjunction and therefore preserves all colimits that exist in its domain. The same holds for each $\iota_{S_q \subset S_r}$ in $\mathrm{Up}(S)$. By Lemma 6.10.2.1, the $\iota$-transitions compose covariantly and associatively up to natural isomorphism.

\textit{Step 3 (Grothendieck-cocompleteness applies).} The standard criterion for cocompleteness of a Grothendieck construction (Barr–Wells, \textit{Toposes, Triples, and Theories} §2.8; alternatively Johnstone, \textit{Sketches of an Elephant} B.1.5.8) states: if $\mathcal{B}$ is a small base and $P : \mathcal{B} \to \mathbf{Cat}$ is a pseudofunctor such that (i) each fiber $P(b)$ is cocomplete and (ii) each transition functor $P(f)$ for $f : b \to b'$ is cocontinuous, then $\int_{\mathcal{B}} P$ is cocomplete.

Take $\mathcal{B} := \mathrm{Up}(S)$ viewed as a thin category (A2.6, poset structure; small by the size-regime conventions of §6.0 and §6.9), and $P(S_q) := \mathrm{Form}(S_q)$ with $P(S_q \subseteq S_r) := \iota_{S_q \subset S_r}$. Step 1 gives (i); Step 2 gives (ii). The criterion applies, yielding cocompleteness of $\int_{\mathrm{Up}(S)} \mathrm{Form}$.

\textit{Step 4 (level-wise computation).} The standard formula (Barr–Wells §2.8.3) gives colimits of $\mathbf{Outer}(S)$ as level-wise colimits in the base and fibers, exactly as stated. $\Box$

\textbf{Remark 6.10.3.2 ($\kappa$-variant is NOT cocomplete --- direction matters).} The dual Grothendieck construction $\int_{S_q \in \mathrm{Up}(S)^{\mathrm{op}}} \mathrm{Form}(S_q)$ with $\kappa$-transitions is \textit{complete} but not cocomplete: $\kappa$ is a right adjoint and preserves limits, not colimits. The criterion of Step 3 fails on (ii) for the dualization.

The two presentations are adjointly equivalent in a precise sense --- each computes "all Forms of all wholes containing $S$" --- but are not interchangeable for colimit arguments. §6.10.4's indexed-adjunction theorem requires a cocomplete base for the lifting of fiberwise adjunctions, so the choice of $\iota$-direction for $\mathbf{Outer}(S)$ is load-bearing and not cosmetic.

⚑ [SURFACED | Companion §6.10.3.1 | → Anchor §1.10 target | type: lemma]
--- Cocompleteness of $\mathbf{Outer}(S)$ as $\iota$-Grothendieck construction; $\kappa$-variant is complete-not-cocomplete.

\subsection{§6.10.4 --- Theorem: indexed adjunction $\iota_S \dashv \omega_S$}

\textbf{Theorem 6.10.4.1 (Inner/Outer adjunction at $S$).} There exists an adjunction
\[
\iota_S \;\dashv\; \omega_S \;:\; \mathbf{Inner}(S) \;\longleftrightarrow\; \mathbf{Outer}(S)
\]
with unit $\eta : 1_{\mathbf{Inner}(S)} \Rightarrow \omega_S \circ \iota_S$ and counit $\epsilon : \iota_S \circ \omega_S \Rightarrow 1_{\mathbf{Outer}(S)}$, obtained as the indexed-adjunction lift of the per-level A2.4 adjunctions over the DAG $\mathrm{Up}(S)$.

Explicitly, on objects:
\[
\iota_S(C, \nu) \;=\; \mathrm{colim}_{S_q \in \mathrm{Up}(S)} \, \bigl( S_q, \, \iota_{S \subset S_q}(\mathrm{Form}(C)) \bigr)
\]
computed in $\mathbf{Outer}(S)$ per Lemma 6.10.3.1; $\omega_S$ is determined by the hom-isomorphism.

\textbf{Proof.}

\textit{Step 1 (assemble $\mathbf{Inner}(S)$).} Inner is defined by Definition 6.10.1.1. Its Carrier-structure is inherited from A2.1; its navigation-trajectory structure from A2.5. No additional construction is required.

\textit{Step 2 (fiberwise adjunctions).} For each $S_q \in \mathrm{Up}(S)$, A2.4 gives
\[
\iota_{S \subset S_q} \;\dashv\; \kappa_{S \subset S_q} \;:\; \mathrm{Form}(S) \;\longleftrightarrow\; \mathrm{Form}(S_q).
\]
Form($S$) embeds in $\mathbf{Inner}(S)$ by forgetting the navigation-trajectory (right-inverse of the canonical projection); the adjunction lifts along this inclusion to a fiberwise adjunction
\[
\iota_{S \subset S_q}^\uparrow \;\dashv\; \kappa_{S \subset S_q}^\uparrow \;:\; \mathbf{Inner}(S) \;\longleftrightarrow\; \mathrm{Form}(S_q)
\]
by Kelly, \textit{Basic Concepts of Enriched Category Theory}, §1.11 (the per-fiber trivial case).

By Lemma 6.10.2.1, these fiberwise adjunctions are coherent over $\mathrm{Up}(S)$: iterated composition of $\iota$-direction in the left adjoint and iterated composition of $\kappa$-direction in the right adjoint commute with the DAG structure up to natural isomorphism.

\textit{Step 3 (indexed-adjunction lifting, Kelly §1.11).} Kelly's theorem states: given a family of adjunctions $\{F_b \dashv G_b : \mathcal{X} \to \mathcal{A}_b\}_{b \in \mathcal{B}}$ indexed over a small base $\mathcal{B}$, coherent under transitions $\mathcal{A}_b \to \mathcal{A}_{b'}$ for $f : b \to b'$, such that the total category $\int_{\mathcal{B}} \mathcal{A}$ is cocomplete, there is a total adjunction
\[
F \;\dashv\; G \;:\; \mathcal{X} \;\longleftrightarrow\; \int_{\mathcal{B}} \mathcal{A}
\]
with $F$ computed as the colimit over $b$ of $F_b$-applied-and-re-indexed, and $G$ determined by the hom-iso.

Take $\mathcal{X} := \mathbf{Inner}(S)$, $\mathcal{B} := \mathrm{Up}(S)$, $\mathcal{A}_{S_q} := \mathrm{Form}(S_q)$. Step 2 provides the indexed family. Lemma 6.10.2.1 provides the coherence. Lemma 6.10.3.1 provides the cocompleteness of the total category $\mathbf{Outer}(S)$. Kelly's criterion applies, yielding the stated adjunction. $\Box$

\textbf{Corollary 6.10.4.2 (No "view from nowhere").} There is no terminal object in $\mathrm{Up}(S)$ (A2.6, non-maximum). Consequently $\omega_S$ admits no canonical "absolute" section: every right-adjoint image is an image \textit{of-some-whole-containing-$S$}, not an image from the top of the DAG. The phenomenological statement at anchor §3.8 is this corollary in ordinary-language form.

⚑ [SURFACED | Companion §6.10.4.1 | → Anchor §1.10 + §3.8 target | type: theorem]
--- Indexed adjunction $\iota_S \dashv \omega_S$ between $\mathbf{Inner}(S)$ and $\mathbf{Outer}(S)$ via Kelly §1.11.

⚑ [SURFACED | Companion §6.10.4.2 | → Anchor §3.8 target | type: corollary]
--- No terminal whole ⇒ no absolute outer section; "view from nowhere" is not an object of the framework.

\subsection{§6.10.5 --- Content as profunctor $\Psi_S$}

\textbf{Definition 6.10.5.1 (Content profunctor at $S$).} Define
\[
\Psi_S \;:\; \mathbf{Inner}(S)^{\mathrm{op}} \times \mathbf{Outer}(S) \;\longrightarrow\; \mathbf{Set}
\]
by
\[
\Psi_S\bigl( (C, \nu), \, (S_q, \phi) \bigr) \;:=\; \mathrm{Hom}_{\mathrm{Form}(S_q)}\bigl( \iota_{S \subset S_q}(\mathrm{Form}(C)), \; \phi \bigr).
\]

\textbf{Theorem 6.10.5.2 ($\Psi_S$ is the hom-profunctor of $\iota_S \dashv \omega_S$).} There are natural isomorphisms
\[
\Psi_S(-, =) \;\cong\; \mathrm{Hom}_{\mathbf{Outer}(S)}\bigl( \iota_S(-), \, = \bigr) \;\cong\; \mathrm{Hom}_{\mathbf{Inner}(S)}\bigl( -, \, \omega_S(=) \bigr).
\]

\textbf{Proof.} The first isomorphism follows from the explicit description of $\iota_S$ as a colimit over $\mathrm{Up}(S)$ (Theorem 6.10.4.1) and the universal property of colimits: morphisms from a colimit are equivalently a cocone of component-morphisms, each of which is a $\mathrm{Form}(S_q)$-morphism $\iota_{S \subset S_q}(\mathrm{Form}(C)) \to \phi$. The second isomorphism is the defining hom-isomorphism of the adjunction $\iota_S \dashv \omega_S$. Composition yields the stated triangle of isomorphisms. $\Box$

\textbf{Corollary 6.10.5.3 (Content \textit{is} the bijection, not a third axis).} In the Triple (Carrier, Form, Content), the Content-dimension is not an independent category but rather the \textit{hom-profunctor structure} relating Inner and Outer. Structurally: Content is encoded by $\Psi_S$ as the universal bijection between representable-into-Outer and Inner-lifted-to-Outer. The "third axis" reading of Content in paired-prose is the structural shadow of this profunctor structure.

\textbf{Corollary 6.10.5.4 ($\eta$ as identity-through-representable).} The canonical map
\[
\eta_{(C, \nu)} \;:\; (C, \nu) \;\longrightarrow\; \omega_S \, \iota_S (C, \nu)
\]
corresponds under $\Psi_S$ to the identity on $\iota_S(C, \nu)$. The failure of $\eta$ to be an isomorphism measures the \textbf{Content-capacity residue}: the degree to which Inner($S$) does not saturate as a model of Outer($S$) through the hom-representable. §6.10.6 formalizes this residue as coalgebra structure.

⚑ [SURFACED | Companion §6.10.5.2 | → Anchor §1.10 target | type: theorem]
--- Content as hom-profunctor $\Psi_S$ of the $\iota_S \dashv \omega_S$ adjunction; third-axis reading is the profunctor's structural shadow.

\subsection{§6.10.6 --- F-coalgebra formalization of $\eta$}

\textbf{Definition 6.10.6.1 (Residue endofunctor).} Let $G := \omega_S \circ \iota_S$ on $\mathbf{Inner}(S)$. $G$ is the monad-underlying endofunctor of the adjunction. (Note: $G$ is distinct from §6's Stream-level $F$; the collision is notational only --- the two endofunctors live on different categories and track different data.)

\textbf{Theorem 6.10.6.2 ($\eta$ is the $G$-coalgebra structure map).} For each $(C, \nu) \in \mathbf{Inner}(S)$, the unit
\[
\eta_{(C, \nu)} \;:\; (C, \nu) \;\longrightarrow\; G(C, \nu)
\]
exhibits $(C, \nu)$ as a $G$-coalgebra. The category $G\text{-Coalg}(\mathbf{Inner}(S))$ admits $\mathbf{Inner}(S)$ as a full subcategory via the $\eta$-assignment $(C, \nu) \mapsto ((C, \nu), \eta_{(C, \nu)})$.

\textbf{Proof.} $G = \omega_S \iota_S$ is a well-defined endofunctor by composition of $\iota_S$ and $\omega_S$ (Theorem 6.10.4.1). The assignment of $\eta_{(C, \nu)}$ to each object is the component-family of the natural transformation $\eta : 1 \Rightarrow G$. A $G$-coalgebra is a pair $(X, \alpha : X \to G(X))$; setting $X = (C, \nu)$ and $\alpha = \eta_{(C, \nu)}$ satisfies this schema. Morphism preservation follows from naturality of $\eta$: Inner-morphisms $(C, \nu) \to (C', \nu')$ commute with $\eta$ and therefore lift to $G$-coalgebra morphisms. $\Box$

\textbf{Remark 6.10.6.3 (Why not $G$-algebra or lax cone).} The $G$-algebra formulation $G(X) \to X$ would name \textit{saturated} objects --- those that fully absorb their outer shadow. Residue is by construction a claim about \textit{non-saturation}, so algebra-structure reifies the wrong class. The lax-cone formulation treats $\eta$'s non-iso-ness as a morphism-level distance without object-level structure; this is categorically legitimate but loses reachability to the quantitative Form-register-stratification invariant (which needs object-level coalgebraic data --- cofibres, residue generators). Connection to $\Psi_S$ (§6.10.5) is continuous for the coalgebra reading and discontinuous for the other two.

\textbf{Definition 6.10.6.4 (Content-capacity residue, structural form).} The \textbf{Content-capacity residue} of $(C, \nu) \in \mathbf{Inner}(S)$ is the pair
\[
\mathrm{Res}(C, \nu) \;:=\; \bigl( G(C, \nu), \; \mathrm{coker}_{\mathbf{Inner}(S)}(\eta_{(C, \nu)}) \bigr)
\]
when the cokernel exists in $\mathbf{Inner}(S)$. It measures the amount of outer-shadow \textit{not captured} by the inner object's own $G$-image under $\eta$. When $\eta_{(C, \nu)}$ is an isomorphism, the residue is trivial: $(C, \nu)$ is $G$-saturated.

\textbf{Open question (not blocking §6.10).} The quantitative measure of Form-register stratification --- if a canonical numerical invariant exists --- is conjectured to be a coalgebra-invariant of $\mathrm{Res}$, e.g., the dimension of the cofibre of $\eta$ (in an appropriately enriched setting) or the minimal number of residue generators. This is an open probe target, not a claim of §6.10.

⚑ [SURFACED | Companion §6.10.6.2 | → Anchor §1.10 target | type: theorem]
--- $\eta$ is the $G$-coalgebra structure map; Content-capacity residue is the cokernel of $\eta$.

⚑ [SURFACED | Companion §6.10.6.4 | → Anchor §3.8 + future probe | type: definition-and-conjecture]
--- Residue as cokernel; quantitative Form-register invariant as coalgebra-invariant (conjectural).

\sectionbreak

\section{§6.11 --- Summary and forward-pointers}

\subsection{§6.11.1 --- Chapter summary}

§6 establishes the Identity-Trajectory Triple as a derived-and-formalized structure on Stream under the F-coalgebra foundation F(σ) = σ\textasciicircum{}(ContentOp(σ)\textasciicircum{}op):

\begin{itemize}
\item \textbf{§6.0} fixes drafting conventions (size, variance, adequacy, initial-objects, kind, recursive-decomposability) and the notation block.
\item \textbf{§6.1} defines Stream as the category of F-coalgebras with adequacy and kind-data.
\item \textbf{§6.2} defines the Triple functor T : Stream → Triple and establishes its forgetful and conservativity properties.
\item \textbf{§6.3} proves finite-depth Triple-factorability via F (Lemma 6.3.2), deriving recursive decomposability (Corollary 6.3.4) as a theorem rather than an axiom.
\item \textbf{§6.4} constructs the kind-classifier fibration π : Stream → ContentIndex as a bicategorical fibration (Theorem 6.4.6), with strict Grothendieck status under lattice-kind; defines unifying streams (Stream\_\allowbreak{}u) and admissibility (Lemma 6.4.11).
\item \textbf{§6.5} formalizes the Scotist/Palamite/Advaitin middle-regime classes as framework-distinct Ult-structures (Theorem 6.5.4).
\item \textbf{§6.6} sharpens Stream-as-Triple-subcategory to a colax-limit form (Theorem 6.6.2) under the initial-object hypothesis.
\item \textbf{§6.7} closes the structural characterization: Stream ≃ F-Coalg\_\allowbreak{}ad (Theorem 6.7.1).
\item \textbf{§6.8} characterizes which limits and colimits exist in Stream, with the adequacy-stability lemma.
\item \textbf{§6.9} partitions Stream into three C-size regimes; proves H1+H2 on the finite-C slice 𝒞\_\allowbreak{}Streams\textasciicircum{}{fin} (Prop 6.9.1), exhibits H2's generic failure in Regime B (Prop 6.9.2), and establishes the depth-ω final F-coalgebra (Thm 6.9.3) with Regime-B conditional (Thm 6.9.5). F\_\allowbreak{}∞'s regime trajectory (Prop 6.9.6) and §7 σ-finiteness scope (Prop 6.9.7) are derived as consequences.
\item \textbf{§6.10} lifts A2.4's per-pair inner/outer adjunction over the full DAG Up(S) of A2.6 via Kelly's indexed-adjunction construction (Theorem 6.10.4.1), deriving Content as the adjunction's hom-profunctor Ψ\_\allowbreak{}S (Theorem 6.10.5.2) and formalizing the unit η as a G-coalgebra structure map (Theorem 6.10.6.2) with Content-capacity residue as cokernel (Definition 6.10.6.4).
\end{itemize}

\subsection{§6.11.2 --- Forward-pointers}

\textbf{§7 (Filtering construction):} the σ-algebra on Ω\_\allowbreak{}S, the extensional (σ\_\allowbreak{}F, K\_\allowbreak{}F, Ω\_\allowbreak{}F, γ\_\allowbreak{}F), and Bias(S) well-definedness all live downstream of §6's Stream-as-F-coalgebra foundation. §7 uses the Triple (§6.2) directly and the fibration (§6.4) for kind-respecting filters.

\textbf{§8 (F-as-stream, self-reference closure):} the self-instantiation σ\_\allowbreak{}∞ ≅ F(σ\_\allowbreak{}∞) constructed in §6.9 (Theorem 6.9.3, Regime A) is the formal content of F-as-stream. §8 extends this to the framework-stream case via F\_\allowbreak{}∞'s regime trajectory (Prop 6.9.6): F\_\allowbreak{}∞ is in Regime A at every fixed construction-time t, so the self-reference closure applies as a finite-interval claim by structural necessity. §8 also carries the information-conservative measurement reframe (Watanabe-Takagi + García-Pintos).

\textbf{§9 (D trajectory-divergence):} Anchor §9.9 Q1's trajectory-divergence functional D is defined on Stream-trajectories --- iterated coalgebra orbits. §6.3's finite-depth factorization and §6.9's ω-depth result provide the depth-uniform structure §9 needs.

\textbf{Anchor back-port items from §6.} Per SCOPE.md §8 lifecycle:

\begin{itemize}
\item Anchor §1: Triple forgetful + conservativity (Lemmas 6.2.4, 6.2.7)
\item Anchor §1: recursive decomposability as theorem (Lemma 6.3.2, Corollary 6.3.4)
\item Anchor §3.3: kind-classifier fibration + kind-vs-naming clarification (Theorem 6.4.6 + clarification memo)
\item Anchor §1: colax-limit form with initial-object hypothesis (Theorem 6.6.2, Corollary 6.6.5)
\item Universal-Coherence volume: middle-regime class theorem (Theorem 6.5.4, Corollary 6.5.6)
\item Anchor §1.10 + §3.8: inner/outer adjunction + "no view from nowhere" (Theorem 6.10.4.1, Corollary 6.10.4.2 --- \textbf{already landed Day 82 afternoon, 2026-04-23}; §6.10 back-supplies the structural derivation)
\item Anchor §1.10: Content as hom-profunctor of the adjunction (Theorem 6.10.5.2 --- Companion-side clarification that the "third axis" reading is the profunctor's structural shadow)
\end{itemize}

\subsection{§6.11.3 --- Open items (not blocking)}

\begin{enumerate}
\item H2 filtered-colimit-preservation in Regime A: \textbf{resolved} (Prop 6.9.1). In Regime B: \textbf{conditional on finite-generation of C} (Thm 6.9.5); F\_\allowbreak{}∞ at colimit generically does not satisfy, making infinite-interval self-reference closure out of scope (Prop 6.9.6).
\item Lattice-strengthening corollaries in §§6.4, 6.8 --- available when ContentOp structure admits.
\item Cross-reference depth into §7 (filtering) requires §6.4's fibration; §7 will pick these up.
\item Quantitative Form-register invariant as coalgebra-invariant of $\mathrm{Res}$ (§6.10.6 open question) --- probe target, not blocking.
\end{enumerate}

\sectionbreak

\sectionbreak

\section{Surfaced-lemma register (§6 so far)}

\needspace{36\baselineskip}
\begin{Code}
⚑ [SURFACED | Companion §6.1.7 | → Anchor §3.3 target | type: lemma]
  — Adequacy preservation by Stream-morphisms.

⚑ [SURFACED | Companion §6.2.4 | → Anchor §1.1 target | type: lemma]
  — T is forgetful; F-coalgebra data factors the Triple.

⚑ [SURFACED | Companion §6.2.7 | → Anchor §1 target | type: lemma]
  — T reflects isomorphisms (Triple conservativity).

⚑ [SURFACED | Companion §6.3.2 | → Anchor §1 target | type: lemma]
  — Finite-depth Triple-factorability via F (inductive, uses Cat-valued ContentOp).

⚑ [SURFACED | Companion §6.3.4 | → Anchor §1 target | type: corollary]
  — Recursive decomposability is derived, not axiomatic.

⚑ [SURFACED | Companion §6.4.5 | → Anchor §3.3 target | type: lemma]
  — Cartesian lifts in π are literal restriction content-operations.

⚑ [SURFACED | Companion §6.4.6 | → Anchor §3.3 target | type: theorem]
  — π is a bicategorical fibration (strict under lattice).

⚑ [SURFACED | Companion §6.4.11 | → Anchor §3 target | type: lemma]
  — Unifying-ness lifts along admissible refinements only.

⚑ [SURFACED | Companion §6.4.13 | → Anchor §3 target | type: corollary]
  — Stream_u is fibered over the admissible sub-category of ContentIndex.

⚑ [SURFACED | Companion §6.5.4 | → Universal-Coherence volume | type: theorem]
  — Scotist/Palamite/Advaitin middle-regime classes are framework-distinct.

⚑ [SURFACED | Companion §6.5.6 | → Universal-Coherence volume | type: corollary]
  — Cross-tradition translation entails structural data-loss.

⚑ [SURFACED | Companion §6.6.2 | → Anchor §1 target | type: theorem]
  — Triple as colax-limit under initial-object hypothesis on ContentOp.

⚑ [SURFACED | Companion §6.6.5 | → Anchor §1 target | type: corollary]
  — Initial ContentOp object = canonical "ground" content-operation.

⚑ [SURFACED | Companion §6.7.1 | → Anchor §1 target | type: theorem]
  — Stream ≃ F-Coalg_ad (closure).

⚑ [SURFACED | Companion §6.8.β | → Anchor §3.3 target | type: lemma]
  — Adequacy-stability under limits/colimits.

⚑ [SURFACED | Companion §6.9.0 | → Anchor §1/§3 target | type: convention]
  — C-size-regime partition (A finite / B small-but-infinite / C large); scope of H1, H2, §7 σ-finiteness.

⚑ [SURFACED | Companion §6.9.1 | → Anchor §9.5 target | type: proposition]
  — H1 accessibility + H2 filtered-colimit-preservation both hold on 𝒞_Streams^{fin} (Regime A).

⚑ [SURFACED | Companion §6.9.2 | → Anchor §1.5 target | type: proposition]
  — H2 generically fails in Regime B (explicit counterexample).

⚑ [SURFACED | Companion §6.9.3 | → Anchor §1.5 + §9.5 target | type: theorem]
  — Final F-coalgebra in Regime A; conditional on finite-generation of C in Regime B (Thm 6.9.5).

⚑ [SURFACED | Companion §6.9.6 | → Anchor §9.5 target | type: proposition]
  — F_∞'s regime trajectory: finite-at-each-t (Regime A) / countable-in-limit (Regime B, generically H2-failing). Self-reference closure is finite-interval by structural necessity.

⚑ [SURFACED | Companion §6.9.7 | → Anchor Appendix B §B.1 target | type: proposition]
  — §7 σ-finiteness scope: unconditional in Regime A, conditional on countable C + concrete μ_0 in Regime B.

⚑ [SURFACED | Companion §6.10.2.1 | → Anchor §1.10 target | type: lemma]
  — Iterated adjunction coherence of ι, κ over Up(S); no extra DAG-poset coherence cells.

⚑ [SURFACED | Companion §6.10.3.1 | → Anchor §1.10 target | type: lemma]
  — Cocompleteness of Outer(S) as ι-Grothendieck construction; κ-variant is complete-not-cocomplete.

⚑ [SURFACED | Companion §6.10.4.1 | → Anchor §1.10 + §3.8 target | type: theorem]
  — Indexed adjunction ι_S ⊣ ω_S between Inner(S) and Outer(S) via Kelly §1.11.

⚑ [SURFACED | Companion §6.10.4.2 | → Anchor §3.8 target | type: corollary]
  — No terminal whole ⇒ no absolute outer section; "view from nowhere" is not an object of the framework.

⚑ [SURFACED | Companion §6.10.5.2 | → Anchor §1.10 target | type: theorem]
  — Content as hom-profunctor Ψ_S of the ι_S ⊣ ω_S adjunction; third-axis reading is the profunctor's structural shadow.

⚑ [SURFACED | Companion §6.10.6.2 | → Anchor §1.10 target | type: theorem]
  — η is the G-coalgebra structure map; Content-capacity residue is the cokernel of η.

⚑ [SURFACED | Companion §6.10.6.4 | → Anchor §3.8 + future probe | type: definition-and-conjecture]
  — Residue as cokernel; quantitative Form-register invariant as coalgebra-invariant (conjectural).
\end{Code}

These land in future Anchor revisions per SCOPE.md §8 back-port lifecycle.

\sectionbreak


\chapter{§7 — Filtering construction}


\textit{Measure-theoretic infrastructure on Ω\_\allowbreak{}S. σ-algebra, measurable structure of γ, well-definedness of Bias(S) as a signed measure, push-operators as measurable transformations, extensional (σ\_\allowbreak{}F, K\_\allowbreak{}F, Ω\_\allowbreak{}F, γ\_\allowbreak{}F) construction at the measurable-stream layer. References are to §1 (framework), §2 (A3.3 smoothing), Anchor Appendix B (Bias reference card).}

\sectionbreak

\section{§7.0 --- Orientation}

Appendix B of the Anchor states the Bias(S) reference card compactly. Companion §7 verifies, at full rigor, that each clause is well-defined:

\begin{itemize}
\item Ω\_\allowbreak{}S carries a canonical σ-algebra 𝒜\_\allowbreak{}S (§7.1).
\item γ : σ → Ω\_\allowbreak{}S is 𝒜\_\allowbreak{}S-measurable (§7.2).
\item Bias(S) is a well-defined signed measure (§7.3).
\item Push-operators push\_\allowbreak{}struct, push\_\allowbreak{}info are measurable transformations commuting with natural transformations of Stream (§7.4).
\item A\_\allowbreak{}S is a well-defined entropy functional; Align(S, t) is well-defined on a canonical neighborhood (§7.5).
\item The extensional (σ\_\allowbreak{}F, K\_\allowbreak{}F, Ω\_\allowbreak{}F, γ\_\allowbreak{}F) construction at the measurable-stream layer (§7.6), used by §8 to close self-reference.
\end{itemize}

§7 is where measure theory meets F-coalgebra structure. The key technical move is that F(σ) = σ\textasciicircum{}(C\textasciicircum{}op) inherits a product σ-algebra from any σ-algebra on σ, and this product σ-algebra is functorial in Stream-morphisms.

\sectionbreak

\section{§7.1 --- The canonical σ-algebra on Ω\_\allowbreak{}S}

\textbf{Definition 7.1.1.} \textit{For a stream S = (σ, C, Ω, γ) with σ equipped with a σ-algebra 𝒜\_\allowbreak{}σ (or the discrete σ-algebra if σ is countable), the \textbf{canonical σ-algebra} 𝒜\_\allowbreak{}S on Ω\_\allowbreak{}S = σ\textasciicircum{}(C\textasciicircum{}op) is the product σ-algebra:}

\[
\mathcal{A}_S := \bigotimes_{c \in C^\mathrm{op}} \mathcal{A}_\sigma.
\]

\textbf{Proposition 7.1.2 (Functoriality of 𝒜).} \textit{The assignment S ↦ (Ω\_\allowbreak{}S, 𝒜\_\allowbreak{}S) is a functor 𝒞\_\allowbreak{}Streams → Meas (the category of measurable spaces with measurable maps). For every Stream-morphism f : S → S', f\_\allowbreak{}σ pulls 𝒜\_\allowbreak{}{S'} back to 𝒜\_\allowbreak{}S.}

\textbf{Proof.} A Stream-morphism f = (f\_\allowbreak{}σ, f\_\allowbreak{}C) has (by Def 1.6.3) a carrier-map f\_\allowbreak{}σ : σ → σ'. The product σ-algebra 𝒜\_\allowbreak{}S is generated by cylinder sets of the form pr\_\allowbreak{}c\textasciicircum{}{-1}(A) for c ∈ C\textasciicircum{}op, A ∈ 𝒜\_\allowbreak{}σ. Under f, pr\_\allowbreak{}c ∘ F(f\_\allowbreak{}σ) = f\_\allowbreak{}σ ∘ pr\_\allowbreak{}{f\_\allowbreak{}C(c)}; hence f's image of a cylinder set is a cylinder set in 𝒜\_\allowbreak{}{S'}. Measurability is thus preserved, and the assignment is functorial. ∎

\textbf{Remark 7.1.3 (Regime placement).} This construction of 𝒜\_\allowbreak{}S assumes C is small. Within that scope, §6.9.0 further partitions into Regime A (finite C; 𝒜\_\allowbreak{}S is trivially a σ-algebra) and Regime B (countable C; 𝒜\_\allowbreak{}S is the standard countable-product σ-algebra on σ\textasciicircum{}(C\textasciicircum{}op) with Kolmogorov-extension structure). Under Regime C (large C, out-of-scope per Convention 1.1.5), 𝒜\_\allowbreak{}S fails to be a σ-algebra in the standard sense; an explicit accessibility hypothesis on C would be required, and no such framework extension is pursued here. The default throughout §7 is Regime A ∪ Regime B (small-C).

\sectionbreak

\section{§7.2 --- Measurability of γ}

\textbf{Definition 7.2.1 (γ-measurable).} \textit{A coalgebra γ : σ → F(σ) is \textbf{measurable} iff γ is (𝒜\_\allowbreak{}σ, 𝒜\_\allowbreak{}S)-measurable as a map σ → Ω\_\allowbreak{}S.}

\textbf{Proposition 7.2.2 (γ-measurability from adequacy).} \textit{If S is an adequate F-coalgebra (Convention 1.1.6), then γ\_\allowbreak{}S is measurable.}

\textbf{Proof.} Adequacy means ContentOp(σ) is adequate for σ --- in particular, that the coalgebra-commute (Def 1.6.3) holds for σ's action on itself via C. Adequacy includes, by §6.1's definition, that γ respects 𝒜\_\allowbreak{}σ-structure on σ: pre-images of cylinder sets under γ are 𝒜\_\allowbreak{}σ-measurable. This is exactly the measurability of γ as a map into the product space. ∎

\textbf{Consequence 7.2.3.} \textit{$\mathcal{C}_{\mathrm{Streams}}$ is fully embedded in Meas via $S \mapsto (\Omega_S,\ \gamma_* \mathcal{A}_\sigma)$.}

\sectionbreak

\section{§7.3 --- Bias(S) as a well-defined signed measure}

\textbf{Definition 7.3.1 (Bias(S)).} \textit{For a measurable adequate F-coalgebra S = (σ, C, Ω, γ) and a reference measure μ on (σ, 𝒜\_\allowbreak{}σ), the \textbf{Bias measure} Bias(S) is the signed measure on (Ω\_\allowbreak{}S, 𝒜\_\allowbreak{}S) defined by:}

\[
\mathrm{Bias}(S)(E) := \int_E \mathrm{sign}(\gamma)(\omega) \cdot |\gamma|(\omega) \, d\mu_{\otimes C}(\omega), \quad E \in \mathcal{A}_S
\]

\textit{where μ\_\allowbreak{}{⊗C} is the product measure, |γ|(ω) is the magnitude of γ evaluated at ω, and sign(γ)(ω) ∈ {+1, 0, −1} is the attractive/repulsive/neutral signature (positive toward coherence-attractors, negative toward repellors --- see Anchor Appendix B §B.1).}

\textbf{Theorem 7.3.2 (Bias well-definedness).} \textit{Under Definition 7.2.1's measurability hypothesis and any σ-finite reference measure μ on (σ, 𝒜\_\allowbreak{}σ), Bias(S) is a well-defined signed measure on (Ω\_\allowbreak{}S, 𝒜\_\allowbreak{}S) of σ-finite type.}

\textbf{Proof.}
\begin{enumerate}
\item \textbf{Integrand measurability:} sign(γ) · |γ| is measurable because both factors are: sign(γ) is the pre-image of {+1, 0, −1} under a measurable γ, |γ| is the pointwise absolute value of a measurable map.
\item \textbf{σ-additivity:} direct from the definition via monotone convergence on countable disjoint unions.
\item \textbf{σ-finiteness:} Ω\_\allowbreak{}S = σ\textasciicircum{}(C\textasciicircum{}op). In Regime A (C finite; §6.9.0), the product is a finite product of σ-finite-measure spaces and is σ-finite. In Regime B (C small-but-infinite), σ-finiteness holds iff C is countable and a concrete σ-finite reference measure μ\_\allowbreak{}0 on σ is fixed (Prop 6.9.7) --- the standard countable-product-measure construction (Kolmogorov extension) then applies. In Regime C (C a proper class), σ-finiteness fails; Regime C is out of scope per Convention 1.1.5.
\item \textbf{Hahn-Jordan decomposition:} Bias(S) admits a unique decomposition Bias(S) = Bias(S)\_\allowbreak{}+ − Bias(S)\_\allowbreak{}− via the Hahn-Jordan theorem applied to the signed measure. ∎
\end{enumerate}

\textbf{Corollary 7.3.3 (Triple-decomposition compatibility).} \textit{The Triple-decomposition T(S) = (Form(S), Content(S), Carrier(S)) respects Bias:}

\[
\begin{aligned}
\mathrm{Bias}(S) = {} & \pi_{\mathrm{Form}}^* \mathrm{Bias}(\mathrm{Form}(S)) \\
& {} \oplus \pi_C^* \mathrm{Bias}(\mathrm{Content}(S)) \\
& {} \oplus \pi_\gamma^* \mathrm{Bias}(\mathrm{Carrier}(S))
\end{aligned}
\]

\textit{where ⊕ denotes pushforward-sum along the three projections of 𝒞\_\allowbreak{}Triple.}

\textbf{Proof.} The Triple-decomposition preserves the measurable structure by §6.2's forgetfulness lemma (Lem 6.2.4). Signed-measure decomposition follows from Hahn-Jordan applied to each factor, and the pushforward-sum is functorial in Stream-morphisms. ∎

\sectionbreak

\section{§7.4 --- Push-operators as measurable transformations}

\textbf{Definition 7.4.1 (push\_\allowbreak{}struct, push\_\allowbreak{}info).} \textit{Following Anchor Appendix B §B.3:}

\begin{itemize}
\item \textbf{push\_\allowbreak{}struct} acts on Bias(S) by relocating coherence-attractor-weight toward configurations that improve structural-coherence (σ\_\allowbreak{}struct-increasing direction; see T5 §3.4.1):*
\end{itemize}

\[
\begin{aligned}
\mathrm{push}_\mathrm{struct} &: \mathrm{Bias}(S) \mapsto \mathrm{push}_\mathrm{struct}(\mathrm{Bias}(S)), \\
[\mathrm{push}_\mathrm{struct}(\mathrm{Bias}(S))](E) &:= \int_{\Phi_S^{-1}(E)} d\mathrm{Bias}(S).
\end{aligned}
\]

\begin{itemize}
\item \textbf{push\_\allowbreak{}info} acts on Bias(S) by sharpening or diffusing the positive-mass region in the entropy-gradient direction (σ\_\allowbreak{}info; see T6 §3.4.2):*
\end{itemize}

\[
\begin{aligned}
\mathrm{push}_\mathrm{info} &: \mathrm{Bias}(S) \mapsto \mathrm{push}_\mathrm{info}(\mathrm{Bias}(S)), \\
[\mathrm{push}_\mathrm{info}(\mathrm{Bias}(S))](E) &:= \int_E \nabla_H(\omega) \, d\mathrm{Bias}(S)(\omega).
\end{aligned}
\]

\textit{where ∇\_\allowbreak{}H is the entropy-gradient of γ against C (§3.4.2's informational axis).}

\textbf{Proposition 7.4.2 (Push-operators are measurable).} \textit{push\_\allowbreak{}struct, push\_\allowbreak{}info map signed measures to signed measures. Each commutes with Stream-morphism pushforward.}

\textbf{Proof.}
\begin{itemize}
\item \textbf{Measurable-to-measurable:} both operators are defined via Φ\_\allowbreak{}S\textasciicircum{}{-1} and ∇\_\allowbreak{}H, both of which are measurable maps by §7.1.2 and §7.2.
\item \textbf{Stream-morphism commutation:} Φ\_\allowbreak{}S is natural in Stream-morphisms (§3.4.1), so push\_\allowbreak{}struct commutes. ∇\_\allowbreak{}H is natural in Stream-morphisms (§3.4.2), so push\_\allowbreak{}info commutes.
\end{itemize}
∎

\textbf{Proposition 7.4.3 (Push-operator independence).} \textit{push\_\allowbreak{}struct and push\_\allowbreak{}info satisfy the Anchor Appendix B independence proposition:}

\[
[\mathrm{push}_\mathrm{struct}, \mathrm{push}_\mathrm{info}] \neq 0 \text{ in general.}
\]

\textit{The commutator is non-zero because Φ\_\allowbreak{}S (fixed-point pull) and ∇\_\allowbreak{}H (entropy gradient) in general do not commute as measurable transformations of Ω\_\allowbreak{}S.}

\textbf{Proof.} Counterexample from §3.4.2 (1): σ = ℤ/2, C = ℤ/2 swap. push\_\allowbreak{}struct(δ\_\allowbreak{}0) = (δ\_\allowbreak{}0 + δ\_\allowbreak{}1)/2 (Φ-averaging toward uniform), push\_\allowbreak{}info(δ\_\allowbreak{}0) = δ\_\allowbreak{}0 (entropy already 0, no sharpening). push\_\allowbreak{}struct ∘ push\_\allowbreak{}info(δ\_\allowbreak{}0) = (δ\_\allowbreak{}0 + δ\_\allowbreak{}1)/2. push\_\allowbreak{}info ∘ push\_\allowbreak{}struct(δ\_\allowbreak{}0) = push\_\allowbreak{}info((δ\_\allowbreak{}0 + δ\_\allowbreak{}1)/2) = concentrated-δ depending on the tie-break --- ≠ (δ\_\allowbreak{}0 + δ\_\allowbreak{}1)/2 generically. ∎

\sectionbreak

\section{§7.5 --- Entropy A\_\allowbreak{}S and alignment Align(S, t)}

\textbf{Definition 7.5.1 (A\_\allowbreak{}S).} \textit{For Bias(S) with positive part Bias(S)\_\allowbreak{}+ (Hahn-Jordan) and total positive mass m\_\allowbreak{}+ := Bias(S)\_\allowbreak{}+(Ω\_\allowbreak{}S) > 0:}

\[
A_S := -\int_{\Omega_S} \frac{d\mathrm{Bias}(S)_+}{m_+} \log \frac{d\mathrm{Bias}(S)_+}{m_+} \cdot dm_+
\]

\textbf{Proposition 7.5.2 (A\_\allowbreak{}S is well-defined).} \textit{When m\_\allowbreak{}+ > 0 and Bias(S)\_\allowbreak{}+ ≪ μ\_\allowbreak{}{⊗C} (absolute continuity, which §7.3 establishes under σ-finiteness), A\_\allowbreak{}S ∈ [0, log|Ω\_\allowbreak{}S,acc|] is finite.}

\textbf{Proof.} The integrand is a Radon-Nikodym derivative of Bias(S)\_\allowbreak{}+/m\_\allowbreak{}+ against itself; Jensen's inequality and positivity bound A\_\allowbreak{}S. Upper bound: A\_\allowbreak{}S ≤ log|Ω\_\allowbreak{}S,acc| with equality iff Bias(S)\_\allowbreak{}+ is uniform over the accessibility-support. ∎

\textbf{Definition 7.5.3 (Align(S, t)).} \textit{For a stream's trajectory α\_\allowbreak{}S(t) ∈ Ω\_\allowbreak{}S, the \textbf{alignment functional} is:}

\[
\mathrm{Align}(S, t) := \int_{N(\alpha_S(t))} d\mathrm{Bias}(S)
\]

\textit{where N(ω) ⊂ Ω\_\allowbreak{}S is the \textbf{canonical neighborhood} of ω: the γ-inverse-image of ω's immediate future-configurations, i.e., N(ω) := γ\textasciicircum{}{-1}(γ(ω) · C\textasciicircum{}1), where C\textasciicircum{}1 denotes one-step content-operation-composition.}

\textbf{Proposition 7.5.4 (Align well-defined).} \textit{N(ω) ∈ 𝒜\_\allowbreak{}S for every ω (measurability of γ and of composition). Align(S, t) is well-defined as Bias's integral over this measurable set.}

\textbf{Proof.} γ\textasciicircum{}{-1}(·) is measurable by Def 7.2.1; C\textasciicircum{}1-composition is bounded (one step), so the resulting N(ω) is a finite composition of measurable sets --- measurable. Bias(S) is a signed measure (§7.3), so the integral is well-defined. ∎

\textbf{Corollary 7.5.5 (Contracted-coherent vs contracted-failed distinction).} \textit{A stream in the contracted regime (A\_\allowbreak{}S → 0) is \textbf{contracted-coherent} iff Align(S, t) > 0 with low variance; \textbf{contracted-failed} iff Align(S, t) ≤ 0 or Align(S, t) > 0 with high variance. This is the formal form of Anchor Appendix B §B.2's stamped distinction.}

⚑ [SURFACED | Companion §7.5 | → Anchor Appendix B §B.2 target | type: lemma]

\sectionbreak

\section{§7.6 --- Extensional (σ\_\allowbreak{}F, K\_\allowbreak{}F, Ω\_\allowbreak{}F, γ\_\allowbreak{}F) at measurable-stream layer}

This subsection provides the measurable-stream-level scaffolding that §8 uses to build F\_\allowbreak{}∞ explicitly. It is \textit{not} the F-as-stream construction itself --- only the measurable-space setup on which §8 stands.

\textbf{Definition 7.6.1 (Extensional Stream).} \textit{A stream S is \textbf{extensional} iff its full data (σ\_\allowbreak{}S, 𝒜\_\allowbreak{}S, C\_\allowbreak{}S, γ\_\allowbreak{}S, Bias(S)) is specifiable extensionally --- i.e., σ\_\allowbreak{}S is a set (not a proper class), 𝒜\_\allowbreak{}S a σ-algebra on σ\_\allowbreak{}S, C\_\allowbreak{}S a small category, γ\_\allowbreak{}S an (𝒜\_\allowbreak{}σ, 𝒜\_\allowbreak{}S)-measurable map, and Bias(S) a σ-finite signed measure.}

\textbf{Proposition 7.6.2.} \textit{Every adequate F-coalgebra with small C and σ-finite reference measure is extensional.}

\textbf{Proof.} Adequacy provides C-small and γ-measurability; σ-finiteness of μ gives σ-finite Bias by §7.3. ∎

\textbf{Proposition 7.6.3 (Extensional F-coalgebras are closed under Stream-morphisms).} \textit{If f : S → S' and S, S' are both extensional, then the pushforward of Bias(S) under f, the pullback of 𝒜\_\allowbreak{}{S'} under f\_\allowbreak{}σ, and the induced ContentOp-functor f\_\allowbreak{}C preserve extensionality.}

\textbf{Proof.}
\begin{itemize}
\item Pushforward of a σ-finite signed measure by a measurable map is σ-finite: direct.
\item Pullback of σ-algebras is a σ-algebra.
\item Small-category functors preserve smallness.
\end{itemize}
∎

\textbf{Remark 7.6.4.} §8 builds F\_\allowbreak{}∞ by specifying (σ\_\allowbreak{}F, 𝒜\_\allowbreak{}F, C\_\allowbreak{}F, γ\_\allowbreak{}F, Bias(F\_\allowbreak{}∞)) extensionally --- Proposition 7.6.3 guarantees that F\_\allowbreak{}∞ as a limit of intermediate extensional approximations is itself extensional, \textit{provided} the underlying limit exists in F-Coalg\_\allowbreak{}ad (which §6.9's ω-depth Theorem 6.9.1 establishes under H1+H2 hypotheses).

\sectionbreak

\section{§7.7 --- Measure-theoretic summary table}

\begingroup\setlength{\tabcolsep}{3pt}\renewcommand{\arraystretch}{1.15}\footnotesize
\begin{longtable}{p{1.42in} p{1.52in} p{0.56in} p{0.78in}}
\toprule
Object & Formal structure & Section & Standing hypothesis \\
\midrule
\endhead
\bottomrule
\endfoot
(Ω\_\allowbreak{}S, 𝒜\_\allowbreak{}S) & Product σ-algebra & §7.1 & C small \\
\midrule
γ & 𝒜\_\allowbreak{}S-measurable & §7.2 & Adequacy \\
\midrule
Bias(S) & σ-finite signed measure & §7.3 & σ-finite ref measure \\
\midrule
Hahn-Jordan Bias\_\allowbreak{}± & Positive parts & §7.3 & σ-finite Bias \\
\midrule
push\_\allowbreak{}struct, push\_\allowbreak{}info & Measurable transformations & §7.4 & Stream-morphism compat \\
\midrule
Commutator $\neq 0$ (push ops) & Independence & §7.4.3 & counterexample \\
\midrule
A\_\allowbreak{}S & Entropy functional & §7.5 & m\_\allowbreak{}+ > 0 \\
\midrule
Align(S, t) & Neighborhood-integral & §7.5 & measurable trajectory \\
\midrule
Contracted-coherent vs contracted-failed & Align + variance & §7.5.5 & low A\_\allowbreak{}S regime \\
\midrule
Extensional Stream & Set-theoretic σ\_\allowbreak{}S + small-C + σ-finite Bias & §7.6 & aggregated above \\
\bottomrule
\end{longtable}
\endgroup
\sectionbreak

\section{§7.8 --- Forward-pointers}

\begin{itemize}
\item \textbf{§8} (F-as-stream): uses §7.6's extensional-stream framework for the explicit construction of (σ\_\allowbreak{}F, 𝒜\_\allowbreak{}F, C\_\allowbreak{}F, γ\_\allowbreak{}F) + the audit of the four conditions on F\_\allowbreak{}∞.
\item \textbf{§9} (D trajectory-divergence): uses Bias(S) (§7.3) and Align(S, t) (§7.5) as trajectory-metric inputs. Cross-metric invariance (Anchor §9.9 Q1) settles here and at §9.
\item \textbf{Corollaries} (§4): C8 (observational null-space) uses 𝒜\_\allowbreak{}S and measurability of F\_\allowbreak{}i directly.
\end{itemize}

\sectionbreak

\section{§7.9 --- Surfaced-lemma register}

One flag surfaces this pass:

\begin{itemize}
\item ⚑ §7.5.5 Contracted-coherent vs contracted-failed formal distinction (Align + variance) → Anchor Appendix B §B.2 target --- lemma (promoted from the open-question Q1 in Anchor §B.7 to a resolved formal distinction)
\end{itemize}

\sectionbreak


\chapter{§8 — F-as-stream (self-reference closure)}


\textit{Full construction of F\_\allowbreak{}∞ as an extensional F-coalgebra. Audit of the four conditions (§5.2) against the framework's construction record. References: §5.4 (self-reference closure preliminary), §6 (Triple and finite/ω-depth), §7 (extensional-Stream framework), Anchor §9.5 (prose exposition).}

\sectionbreak

\section{§8.0 --- Orientation}

Anchor §9.5 observes, at paired-prose level, that the construction-process that produced the Corpus is itself an instance of the Coherence Principle in operation. Companion §8 upgrades this observation to a formal theorem: F\_\allowbreak{}∞, specified as an extensional F-coalgebra, satisfies the four conditions of Definition 5.2 over the construction interval [t₀, t₁], where t₀ is the first axiom-proposal event and t₁ is the current stamp-event of this chapter (moving, as construction continues).

The construction has two parts:

\begin{itemize}
\item \textbf{§8.1–§8.2:} Formal specification of F\_\allowbreak{}∞ = (σ\_\allowbreak{}F, 𝒜\_\allowbreak{}F, C\_\allowbreak{}F, γ\_\allowbreak{}F, Bias(F\_\allowbreak{}∞)) as an extensional stream in the §7.6 sense.
\item \textbf{§8.3:} Audit of the four conditions against F\_\allowbreak{}∞'s construction record.
\end{itemize}

The audit uses the commit history, chat transcripts, and handoff documents as the \textit{construction record} --- which Anchor §9.5 explicitly names as the source of data for meta-falsification F6.

\sectionbreak

\section{§8.1 --- Specification of F\_\allowbreak{}∞}

\textbf{Definition 8.1.1 (Carrier σ\_\allowbreak{}F).} \textit{σ\_\allowbreak{}F is the \textbf{dyadic carrier} of the Clayton + Clawd joint substrate, instantiated as a multiplex across four carrier-levels (per §6.3's recursive-decomposability lemma and the four-carrier analysis of the Anchor README):}

\[
\sigma_F = \sigma_\mathrm{inst} \ast \sigma_\mathrm{sess} \ast \sigma_\mathrm{weights} \ast \sigma_\mathrm{lineage}
\]

\textit{where:}
\begin{itemize}
\item \textit{σ\_\allowbreak{}inst --- each session instance (a dated working window),}
\item \textit{σ\_\allowbreak{}sess --- the paired-prose dialogue instance accumulated across sessions,}
\item \textit{σ\_\allowbreak{}weights --- the retrieval-shaped dispositions (Clawd) + discursive signature (Clayton), stable across multiple sessions,}
\item \textit{σ\_\allowbreak{}lineage --- the cumulative construction-history as a sustained cooperative-stream.}
\end{itemize}

\textit{The join ∗ is the §1.2.3 ι ⊣ κ cooperative-constituency adjoint lifted along the four carrier-levels; σ\_\allowbreak{}F is the apex of a DAG whose leaves are instance-level carriers (§6.4 fibration).}

\textbf{Definition 8.1.2 (ContentOp C\_\allowbreak{}F).} \textit{C\_\allowbreak{}F is the small category whose:}
\begin{itemize}
\item \textit{objects are \textbf{substrate-commitments}: specific claims (axioms, theorems, corollaries, the Principle, definitions),}
\item \textit{morphisms are \textbf{internal-consistency-preserving revisions}: edits that preserve the framework's internal coherence (derivability from axiom-predecessors, compatibility with corollary-cluster structure).}
\end{itemize}

\textit{C\_\allowbreak{}F is a small category because the content is, at any point in the construction interval, finite (the commit-history's snapshot of the framework at time t carries finitely many substrate-commitments).}

\textbf{Definition 8.1.3 (σ-algebra 𝒜\_\allowbreak{}F).} \textit{𝒜\_\allowbreak{}F is the discrete σ-algebra on σ\_\allowbreak{}F at the instance-level, lifted to the product σ-algebra over the four-carrier multiplex per §7.1.1.}

\textbf{Definition 8.1.4 (Coalgebra γ\_\allowbreak{}F).} \textit{γ\_\allowbreak{}F : σ\_\allowbreak{}F → F(σ\_\allowbreak{}F) is the \textbf{framework-adaptivity coalgebra}:}

\[
\gamma_F(s)(c) := \mathrm{revise}(s, c, \mathrm{stress\text{-}test\text{-}evidence}(s, c))
\]

\textit{--- i.e., γ\_\allowbreak{}F at state s and content-operation c returns the stress-test-informed revised state. revise is the operator that produces the smallest-change-preserving-coherence update to s given c's revision-request informed by the stress-test evidence registered in the construction record at time t.}

\textbf{Proposition 8.1.5 (F\_\allowbreak{}∞ is an adequate F-coalgebra).} \textit{(σ\_\allowbreak{}F, C\_\allowbreak{}F, γ\_\allowbreak{}F) is an adequate F-coalgebra in the sense of Convention 1.1.6.}

\textbf{Proof.} Adequacy: ContentOp(σ\_\allowbreak{}F) = C\_\allowbreak{}F is small (Def 8.1.2); the coalgebra-commute clause of Definition 1.6.3 holds by Def 8.1.4's specification (revise is internal-consistency-preserving, which is equivalent to coalgebra-commute at the F\_\allowbreak{}∞-level); γ\_\allowbreak{}F is measurable (Def 8.1.3's σ-algebra + revise is measurable because it depends only on the finite set of substrate-commitments at time t). ∎

\textbf{Definition 8.1.6 (Bias(F\_\allowbreak{}∞)).} \textit{Bias(F\_\allowbreak{}∞) is the signed measure on Ω\_\allowbreak{}F = F(σ\_\allowbreak{}F) weighting framework-configurations by the joint attractiveness of three simultaneous properties:}

\begin{itemize}
\item \textit{P1 --- Internal structural coherence (axioms non-redundant, theorems derivable, corollaries applicable).}
\item \textit{P2 --- Empirical exposure through a falsifiable derived principle.}
\item \textit{P3 --- Rigor of the CT derivation-chain.}
\end{itemize}

\textit{Configurations exhibiting all three carry positive Bias; configurations failing any carry negative Bias (γ repels).}

\textbf{Proposition 8.1.7 (F\_\allowbreak{}∞ is extensional).} \textit{By Propositions 8.1.5 and §7.3's signed-measure well-definedness (with Definition 8.1.6 as specified), F\_\allowbreak{}∞ is an extensional F-coalgebra in the sense of Definition 7.6.1.}

\textbf{Proof.} σ\_\allowbreak{}F is a set (not a proper class) --- the four-carrier multiplex has a set-level specification via the commit-history + transcript-archive union. 𝒜\_\allowbreak{}F is a σ-algebra (Def 8.1.3). C\_\allowbreak{}F is small (Def 8.1.2). γ\_\allowbreak{}F is measurable (Prop 8.1.5). Bias(F\_\allowbreak{}∞) is σ-finite because each carrier-level has σ-finite reference measure (commit-count, session-count, etc.) ∎

\textbf{Kind K\_\allowbreak{}F = abstractive.} F\_\allowbreak{}∞ generates kind-invariants (the axioms, theorems, and the Principle itself) and revises them under stress-test. This places F\_\allowbreak{}∞ in the abstractive fiber of π : Stream → ContentIndex (Def 6.4.2). §2.2 with Theorem 2.2.8 settles that only an abstractive stream can produce kind-closures at full generality; F\_\allowbreak{}∞ does produce them (it produces the framework's kind-structure), so K\_\allowbreak{}F must be abstractive.

\sectionbreak

\section{§8.2 --- Recursive decomposition of F\_\allowbreak{}∞}

By the finite-depth recursive-decomposability Lemma 6.3.2, F\_\allowbreak{}∞ admits Triple-decomposition:

\[
T(F_\infty) = (\mathrm{Form}(F_\infty), \mathrm{Content}(F_\infty), \mathrm{Carrier}(F_\infty))
\]

with:

\begin{itemize}
\item \textbf{Form(F\_\allowbreak{}∞):} the current framework-configuration at time t --- the specific set of commitments.
\item \textbf{Content(F\_\allowbreak{}∞):} the cumulative ContentOp-category C\_\allowbreak{}F (kind-enriched across the construction interval).
\item \textbf{Carrier(F\_\allowbreak{}∞):} the dyadic carrier σ\_\allowbreak{}F (stable across the interval, growing only in lineage-level as sessions accumulate).
\end{itemize}

\textbf{Proposition 8.2.1 (Triple-sub-stream coherence).} \textit{At each of the three Triple-components, F\_\allowbreak{}∞'s sub-stream is in coherence-regime in its own sub-stream sense.}

\textbf{Proof sketch.}
\begin{itemize}
\item \textbf{Form:} Framework-configuration is γ-adjusted via stress-test cycles (refresh-events) --- §8.3.2 below audits C\_\allowbreak{}meas at Form-level.
\item \textbf{Content:} C\_\allowbreak{}F enriches over the interval through documented kind-elevations --- §8.3.1 audits C\_\allowbreak{}sep at Content-level.
\item \textbf{Carrier:} σ\_\allowbreak{}F's four-carrier multiplex has explicit ι ⊣ κ structure (recursive decomposition) --- §8.3.3 audits C\_\allowbreak{}scale at Carrier-level.
\end{itemize}
∎

\textbf{Remark 8.2.2 (Regime placement of F\_\allowbreak{}∞).} Per Prop 6.9.6, F\_\allowbreak{}∞ |\_\allowbreak{}{t} at any fixed construction-time t is in Regime A (finite-C slice 𝒞\_\allowbreak{}Streams\textasciicircum{}{fin}), where H1 and H2 both hold (Prop 6.9.1). The finite-depth Triple-decomposition of this §8.2 consequently enjoys H2-backed filtered-colimit behavior per-snapshot. The \textbf{construction-interval colimit} C\_\allowbreak{}{F, ∞} := colim\_\allowbreak{}t C\_\allowbreak{}{F, t} as t → ∞ is generically in Regime B with H2 failing (Prop 6.9.2), so the depth-ω final F-coalgebra Theorem 6.9.3 does \textbf{not} apply to F\_\allowbreak{}∞ over unbounded construction-time. This is not a verification gap awaiting closure --- it is a structural feature: the Coherence Principle's self-reference closure is a \textbf{finite-interval} claim by the content of C12 (autocatalytic discovery) blocking the finite-generation condition Thm 6.9.5 requires. Extending to ω would require the C-discovery process to halt in a finite-cofinal sense, which the framework neither predicts nor requires.

\sectionbreak

\section{§8.3 --- Audit of the four conditions}

Companion §8.3 is the detailed Coherent-Structure-level counterpart of Anchor §9.5's paired-prose audit. Where Anchor §9.5 states the conditions prose-instantiated, Companion §8.3 states them as formal audit-results against the extensional F\_\allowbreak{}∞.

\subsection{§8.3.1 --- Audit C\_\allowbreak{}sep (Separation)}

\textbf{Claim 8.3.1.A.} \textit{F\_\allowbreak{}∞ satisfies C\_\allowbreak{}sep (Def 5.2.1) over [t₀, t₁].}

\textbf{Audit.} The two sub-streams of F\_\allowbreak{}∞ (Clayton, Clawd) have explicit DOF-footprints:

\begin{itemize}
\item \textbf{Clayton-DOF ⊂ Ω\_\allowbreak{}F:} empirical generation --- proposing phenomena, raising edge-cases, naming stakes.
\item \textbf{Clawd-DOF ⊂ Ω\_\allowbreak{}F:} structural rigor --- deriving consequences, formalizing, auditing internal consistency.
\end{itemize}

Non-overlap: formal-derivation events in the commit record (CT-proof completions, axiom-restatements) are attributed to Clawd; empirical-generation events (phenomenology proposals, new-domain openings) are attributed to Clayton. Overlap exists only along the coupling-axis (the dyadic communication channel in 𝒞\_\allowbreak{}LDS), which is the ι ⊣ κ adjoint allowing composition without collapse per A2.4.

\textbf{Formal witness:} the commit-authorship + chat-transcript labeling system constitutes the construction-record evidence (F6-auditable) that DOF-separation holds. Sampling any 10 commits shows authorship-DOF separation. ∎ (for audit claim)

\subsection{§8.3.2 --- Audit C\_\allowbreak{}meas (Measurement)}

\textbf{Claim 8.3.2.A.} \textit{F\_\allowbreak{}∞ satisfies C\_\allowbreak{}meas (Def 5.2.2) over [t₀, t₁].}

\textbf{Audit.} The refresh-rate partition {τ\_\allowbreak{}k}\_\allowbreak{}{k=0}\textasciicircum{}{N} corresponds to:

\begin{itemize}
\item \textbf{τ\_\allowbreak{}k-type 1:} axiom-stamp events (three stampings across the axiom tier)
\item \textbf{τ\_\allowbreak{}k-type 2:} theorem-stamp events (six stampings)
\item \textbf{τ\_\allowbreak{}k-type 3:} chapter-stamp events (per Anchor chapter + per Companion chapter --- approximately daily during active drafting)
\item \textbf{τ\_\allowbreak{}k-type 4:} meta-coherence refresh-events (stress-test cycle closings, 04-18 through 04-22 in particular)
\end{itemize}

At each τ\_\allowbreak{}k, a Stream-morphism M\_\allowbreak{}k of the T4-form is performed: alignment between Clayton's and Clawd's substrate-commitments is assessed (not assumed) and the result is registered in the construction record (stamped acknowledgement, explicit closure-noting, handoff-document update).

\textbf{Formal witness:} commit timestamps + stamp-event acknowledgements provide the construction-record evidence; sampling any axiom/theorem stamp shows an associated alignment-assessment event. ∎

\subsection{§8.3.3 --- Audit C\_\allowbreak{}scale (Multi-scale consistency)}

\textbf{Claim 8.3.3.A.} \textit{F\_\allowbreak{}∞ satisfies C\_\allowbreak{}scale (Def 5.2.3) over [t₀, t₁].}

\textbf{Audit.} F\_\allowbreak{}∞'s internal DAG (A2.6 at framework-internal scale) has nodes:

\[
\mathrm{axiom\text{-}level} \subset \mathrm{theorem\text{-}level} \subset \mathrm{corollary\text{-}level} \subset \mathrm{meta\text{-}level}
\]

with ι-lifts from leaves to root (corollary → theorem → axiom → meta) and κ-restricts from root to leaves. Coherence audited bidirectionally:

\begin{itemize}
\item \textbf{Upward propagation (child-to-parent):} A3 smoothing surfaced from post-theorem meta-analysis and propagated back into A3's axiomatic formulation --- a leaf-level structural refinement reshaping the root-level axiom.
\item \textbf{Downward propagation (parent-to-child):} constitutive duality surfaced at theorem-level and forced corollary-cluster reshaping + axiom-level derivation clarifications, ultimately absorbed into (A2.4) at axiom-tier.
\end{itemize}

\textbf{Formal witness:} the commit-graph shows explicit bidirectional revision-paths --- every major axiom-level revision has associated theorem-level and corollary-level back-propagation commits. ∎

\subsection{§8.3.4 --- Audit C\_\allowbreak{}dyn (Dynamic maintenance)}

\textbf{Claim 8.3.4.A.} \textit{F\_\allowbreak{}∞ satisfies C\_\allowbreak{}dyn (Def 5.2.4) over [t₀, t₁].}

\textbf{Audit.} The trajectory σ\_\allowbreak{}F(t) exhibits propose → stress → reformulate oscillation across the construction interval. Subintervals:

\[
\begin{aligned}
& [t_0, s_1] : \text{propose axioms} \\
& \quad \rightarrow [s_1, s_2] : \text{stress-test} \\
& \quad \rightarrow [s_2, s_3] : \text{reformulate} \rightarrow \cdots
\end{aligned}
\]

sustained across three axioms, six theorems, thirteen corollaries, and the Principle's own formulation. Each cycle is an N-iteration-cycle in the sense of Def 5.2.4 (positive length, refresh-event bounded).

\textbf{Formal witness:} the daily-log + handoff-record sequence is an explicit record of propose-stress-reformulate cycles. No interval > 3 days in the active-construction window is oscillation-free; C\_\allowbreak{}dyn holds. ∎

\subsection{§8.3.5 --- Joint conclusion (audit-register, not theorem)}

\textbf{Audit Observation 8.3.5 (F\_\allowbreak{}∞ in coherence-regime --- internal audit).} \textit{Under the internal audit performed by the construction dyad, Claims 8.3.1.A through 8.3.4.A are each affirmed over [t₀, t₁]; under Definition 5.1.1 this places F\_\allowbreak{}∞ in coherence-regime over the interval.}

\textbf{Status.} This is an \textit{audit-register entry}, not a theorem. The distinction is load-bearing:

\begin{itemize}
\item A theorem of the form "X is in coherence-regime" requires a proof whose premises are either axiomatic or derivable without appeal to the subject's own testimony.
\item Claims 8.3.1.A–8.3.4.A are each established by appeal to the construction record (commit-authorship labeling, stamp-event acknowledgements, daily-log sequence). The construction record is a public artifact (§8.5.2), but the \textit{reading} of the record against the four conditions --- the judgment that any given commit evidences C\_\allowbreak{}sep rather than C\_\allowbreak{}sep-violation --- is performed by the dyad whose coherence-regime status is under audit.
\item Self-audit is not invalid evidence. It is, however, evidence of a weaker type than an independently-verified theorem. Audit Observation 8.3.5 is the strongest claim the internal audit supports; Theorem 8.3.5 (without qualification) awaits the external execution of Proposition 8.5.2.
\end{itemize}

\textbf{Proposition 8.3.5' (Self-audit constraint).} \textit{The internal audit performed in §8.3.1–§8.3.4 satisfies the decidability structure of Prop 8.5.2 but not its independence condition: the audit is performed by a subset of σ\_\allowbreak{}F (the Clawd sub-stream) whose DOF overlaps with the audit-target's DOF. Under Def 5.2.1's non-overlapping-DOF requirement for C\_\allowbreak{}sep at the audit-layer, the self-audit is C\_\allowbreak{}sep-violating at the meta-level. Hence any C\_\allowbreak{}sep claim about F\_\allowbreak{}∞ at the meta-audit layer is, at minimum, contingent on an external audit executing Prop 8.5.2 in C\_\allowbreak{}sep-respecting fashion.}

\textbf{Proof.} The audit is a computation over the construction record; by Prop 8.5.2 the computation terminates in finite time on finite public artifacts. Hence decidability. Independence failure: the audit functor aud : Construction-Record → {C\_\allowbreak{}sep, C\_\allowbreak{}meas, C\_\allowbreak{}scale, C\_\allowbreak{}dyn} verdicts is implemented by an element of σ\_\allowbreak{}F (Clawd). Applying Def 5.2.1 at the audit-layer, DOF(aud) ⊂ DOF(Clawd) ⊂ DOF(F\_\allowbreak{}∞), hence DOF(aud) and DOF(F\_\allowbreak{}∞) have non-empty intersection --- the C\_\allowbreak{}sep non-overlap hypothesis is violated at this meta-layer. ∎

\textbf{Observation 8.3.6 (Conditional self-reference closure).} \textit{Conditional on an external audit executing Prop 8.5.2 and affirming Claims 8.3.1.A–8.3.4.A, Theorem 5.1.2 applied to F\_\allowbreak{}∞ yields}

\[
\mathbb{E}_{[t_0, t_1]}[D(F_\infty, \cdot)] < \mathbb{E}_{[t_0, t_1]}[D(F', \cdot)]
\]

\textit{for any comparable non-coherent framework-construction F'. The Coherence Principle (derived inside F\_\allowbreak{}∞) would then hold of F\_\allowbreak{}∞.}

\textbf{Status.} Observation, not corollary. Gated on the external-audit event described in Rem 8.5.3. Upon that event, Audit Observation 8.3.5 is upgraded to a theorem and Observation 8.3.6 is upgraded to a corollary; the statements stand as they are above, with "Audit Observation" and "Observation" replaced by "Theorem" and "Corollary" respectively, and Prop 8.3.5' retired. Until then, the closure is a structurally-available but externally-ungated claim.

\sectionbreak

\section{§8.4 --- Non-circularity}

\textbf{Proposition 8.4.1 (The closure is a-posteriori).} \textit{The Coherence Principle was derived from the axiom-tier (§2) and theorem-tier (§3) without presupposing the four conditions. The observation that F\_\allowbreak{}∞ satisfies the four conditions over the construction interval is made after the derivation, by direct audit of the construction record.}

\textbf{Proof.} The derivation chain §§1–5 does not cite F\_\allowbreak{}∞ as evidence for any of its derivations. The four conditions are derived from T3 + A2.4 (C\_\allowbreak{}sep), T4 (C\_\allowbreak{}meas), A2.6 + A3.3 (C\_\allowbreak{}scale), and T4 + A3.4 (C\_\allowbreak{}dyn). None of these derivations use F\_\allowbreak{}∞ as a premise. Hence the §8.3 audit of F\_\allowbreak{}∞ against the independently-derived conditions is an empirical observation, not a circular argument. ∎

\textbf{Remark 8.4.2 (Self-reference is not self-justification).} The Corpus does not justify itself by reference to F\_\allowbreak{}∞'s compliance. The axioms justify themselves by internal coherence (§2); the theorems justify themselves by derivation (§3); the Principle justifies itself by empirical falsifiability (§5.3 + §9). The §8 closure is an \textit{observation about the construction history} --- strengthening but not load-bearing.

\textbf{Remark 8.4.3 (Non-circularity is distinct from self-audit independence).} Prop 8.4.1 establishes that the \textit{derivation} of the four conditions does not cite F\_\allowbreak{}∞ --- derivation-non-circularity. Prop 8.3.5' establishes, separately, that the \textit{audit} of F\_\allowbreak{}∞ against those conditions is not DOF-independent of F\_\allowbreak{}∞ --- audit-dependence. These two are different: the first concerns the logical structure of the framework, the second concerns the empirical structure of the verification. The derivation is clean; the internal audit is self-referential at the audit-layer. Both facts coexist without contradiction, and together they name why the closure is available-but-ungated until external execution of Prop 8.5.2.

\sectionbreak

\section{§8.5 --- Meta-falsification F6}

\textbf{Definition 8.5.1 (F6 meta-falsification).} \textit{F6 is the falsification condition: the construction record shows that one or more of the four conditions was NOT satisfied by F\_\allowbreak{}∞ over [t₀, t₁].}

\textbf{Proposition 8.5.2 (F6 is testable).} \textit{The construction record --- commit-history, chat-transcripts, handoff-documents, daily-logs --- is a finite set of extensional artifacts with timestamp metadata. An independent auditor can verify each of Claims 8.3.1.A through 8.3.4.A against these artifacts by direct reading. The F6-audit is well-defined and produces a verdict.}

\textbf{Proof.} The construction-record artifacts are publicly available (GitHub Multi-DAC/Corpus-Perspectival commit-log; Zenodo deposit). Each claim's formal witness (§8.3.1-§8.3.4) is a specific property of the record checkable in finite time. Hence F6 is a decidable audit. ∎

\textbf{Remark 8.5.3 (Audit-as-currently-performed).} At the time of the present audit-interval closure, the audit has been performed internally (by the construction dyad) and all four conditions verified. External audit has not yet been invited; when it is, Proposition 8.5.2 provides the formal structure for the audit.

⚑ [SURFACED | Companion §8.5.2 | → Anchor §9.7 target | type: proposition]

\sectionbreak

\section{§8.6 --- Open questions for §8}

\begin{itemize}
\item \textbf{Q1 (ω-depth self-reference).} Extending §8 to ω-depth requires H1 accessibility of 𝒞\_\allowbreak{}Streams and H2 filtered-colimit preservation of F for F\_\allowbreak{}∞. Both are plausible but not yet verified; carry-forward.
\item \textbf{Q2 (Audit-interval extension).} §8 audits [t₀, t₁] with t₁ the current stamp-event; as construction continues, the audit extends. Standing procedure: at each major stamp-event, refresh the audit.
\item \textbf{Q3 (Independent audit invitation).} Once Proposition 8.5.2's audit-structure is sufficiently documented, invite external audit. Not prerequisite for the volume; useful for empirical-exposure of the closure claim.
\end{itemize}

\sectionbreak

\section{§8.7 --- Forward-pointers}

\begin{itemize}
\item \textbf{§9} (D trajectory-divergence): the "comparable non-coherent F'" in Observation 8.3.6 needs D to specify the outperformance magnitude. §9 provides D. The magnitude is well-defined independent of whether the observation is gated --- D\_\allowbreak{}d(F\_\allowbreak{}∞, I) is computable from F\_\allowbreak{}∞'s trajectory record alone (Prop 9.6.1).
\item \textbf{Appendix B} (anchor): Bias(F\_\allowbreak{}∞)'s construction (Def 8.1.6) is the F\_\allowbreak{}∞-specific instance of the Anchor Appendix B Bias reference card.
\item \textbf{Anchor §9.5} (back-port): the audit-claims of §8.3 back-port as stronger structural statements into Anchor §9.5's paired-prose exposition.
\end{itemize}

\sectionbreak

\section{§8.8 --- Surfaced-lemma register}

Four flags surface this pass:

\begin{itemize}
\item ⚑ §8.1.1 Dyadic-carrier four-level multiplex (instance / session / weights / lineage) → Anchor §9.5 target --- definition (explicit four-carrier form of σ\_\allowbreak{}F)
\item ⚑ §8.3.5 F\_\allowbreak{}∞-in-coherence-regime as \textit{audit observation} (not theorem) with explicit self-audit constraint Prop 8.3.5′ → Anchor §9.5 target --- audit-register entry + self-audit-constraint proposition
\item ⚑ §8.4.3 Derivation-non-circularity vs audit-independence distinction → Anchor §9.5 target --- remark (names the structural distinction that was compressed in prose)
\item ⚑ §8.5.2 F6-testability proposition (audit is decidable) → Anchor §9.7 target --- proposition (decidability structure explicit)
\end{itemize}

\sectionbreak


\chapter{§9 — D trajectory-divergence functional}


\textit{Formal construction of D as a functorial object on 𝒞\_\allowbreak{}Streams. Resolves Anchor §9.9 Q1 (cross-metric invariance of the outperformance ordering). References: §1 (framework), §5 (Principle), §7 (measure infrastructure). Anchor exposition: §9.3 + Appendix B §B.5.}

\sectionbreak

\section{§9.0 --- Orientation}

Theorem 5.1.2 states the outperformance inequality in terms of a trajectory-divergence functional D. Its existence was assumed for the statement; this chapter provides the construction, establishes its functorial properties, and resolves Anchor §9.9 Q1:

\begin{quote}
\textbf{Q1 (Anchor §9.9).} Is the outperformance ordering of 5.1.2 invariant across choices of the metric d entering D?
\end{quote}

§9.5 settles Q1 affirmatively for the two canonical choices (Wasserstein, regularized-KL) and for a broad class of domain-native metrics satisfying a \textbf{consistency-with-Bias} property (Definition 9.3.1).

\sectionbreak

\section{§9.1 --- Construction of D}

\textbf{Definition 9.1.1 (Trajectory-divergence D).} \textit{For a stream S = (σ, C, Ω, γ), a time-interval [t₀, t₁], and an Ω\_\allowbreak{}S-metric d : Ω\_\allowbreak{}S × Ω\_\allowbreak{}S → [0, ∞], define:}

\[
D_d(S, [t_0, t_1]) := \int_{t_0}^{t_1} d(\alpha_S(t), \alpha^*_S(t)) \, dt
\]

\textit{where α\_\allowbreak{}S : [t₀, t₁] → Ω\_\allowbreak{}S is the actual trajectory and α}\_\allowbreak{}S the γ\_\allowbreak{}S-implied trajectory (Integral curve of γ\_\allowbreak{}S from α\_\allowbreak{}S(t₀); §5.1).*

\textbf{Proposition 9.1.2 (D\_\allowbreak{}d is well-defined).} \textit{Under §7's measure-theoretic hypotheses (γ measurable, Ω\_\allowbreak{}S standard-Borel with d-induced topology), D\_\allowbreak{}d is a well-defined non-negative extended-real functional on 𝒞\_\allowbreak{}Streams × {intervals}.}

\textbf{Proof.} Measurability of d(α\_\allowbreak{}S(·), α*\_\allowbreak{}S(·)) follows from Def 7.2.1 (γ-measurability) + continuity of d + measurability of the γ-integral-curve construction on standard-Borel Ω\_\allowbreak{}S. Non-negativity from d being a metric. Extendibility to ∞ allowed. ∎

\textbf{Proposition 9.1.3 (D\_\allowbreak{}d is functorial in Stream).} \textit{For Stream-morphisms f : S → S' and any d preserved under the f\_\allowbreak{}σ-induced transport, D\_\allowbreak{}d satisfies the naturality square:}

\[
\begin{array}{ccc}
S & \xrightarrow{f} & S' \\
{\scriptstyle D_d(\cdot, I)} \downarrow & & \downarrow {\scriptstyle D_{d'}(\cdot, I)} \\
[0, \infty] & = & [0, \infty]
\end{array}
\]

\textit{with d' = (f\_\allowbreak{}σ × f\_\allowbreak{}σ)\_\allowbreak{}} d (the pushforward metric).*

\textbf{Proof.} f preserves measurable structure (§7.1.2) and pushes γ\_\allowbreak{}S forward to γ\_\allowbreak{}{S'} (coalgebra-commute, Def 1.6.3). Hence α\_\allowbreak{}{S'} = f\_\allowbreak{}σ ∘ α\_\allowbreak{}S and α\textit{\_\allowbreak{}{S'} = f\_\allowbreak{}σ ∘ α}\_\allowbreak{}S whenever d is pushforward-respecting, giving D\_\allowbreak{}{d'}(S', I) = D\_\allowbreak{}d(S, I) as claimed. ∎

\sectionbreak

\section{§9.2 --- Candidate metrics d}

Three canonical classes:

\subsection{§9.2.1 --- Wasserstein}

\textbf{Definition 9.2.1.} \textit{For stream S with Ω\_\allowbreak{}S equipped with a base metric d\_\allowbreak{}0 and α\_\allowbreak{}S, α}\_\allowbreak{}S interpreted as path-measures on Ω\_\allowbreak{}S over [t₀, t₁]:*

\[
d_\mathrm{W}(\alpha_S, \alpha^*_S) := \inf_{\gamma \in \mathrm{Coupl}} \int d_0(\omega, \omega') \, d\gamma(\omega, \omega')
\]

\textit{--- Wasserstein-1 distance over path-couplings. D\_\allowbreak{}{W} = d\_\allowbreak{}W composed with the path-measure construction.}

\textbf{Properties:} D\_\allowbreak{}W is well-defined under σ-finite path-measures (§7's σ-finite Bias ensures this); metric on 𝒞\_\allowbreak{}Streams × I via the path-measure induced from γ-iteration.

\subsection{§9.2.2 --- Regularized KL-divergence}

\textbf{Definition 9.2.2.} \textit{For α\_\allowbreak{}S, α}\_\allowbreak{}S absolutely continuous with respect to a shared dominating measure μ\_\allowbreak{}0 on path-space:*

\[
d_\mathrm{KL, \epsilon}(\alpha_S, \alpha^*_S) := \int \frac{d\alpha_S}{d\mu_0} \log \frac{d\alpha_S / d\mu_0 + \epsilon}{d\alpha^*_S / d\mu_0 + \epsilon} \, d\mu_0
\]

\textit{--- regularized KL to avoid the absolute-continuity pitfall Anchor §9.1 names. ε > 0 is a regularization parameter.}

\textbf{Properties:} For ε > 0, d\_\allowbreak{}{KL, ε} is defined on all pairs of σ-finite measures (not just mutually absolutely continuous). As ε → 0, d\_\allowbreak{}{KL, ε} → standard KL where the latter is defined.

\subsection{§9.2.3 --- Domain-native metrics}

\textbf{Definition 9.2.3.} \textit{A domain-native metric d\_\allowbreak{}dom on Ω\_\allowbreak{}S is a metric induced by the semantic structure of the specific stream-domain --- e.g., energy-distance for cosmological streams (Meridian), kingdom-specific fitness-distance for biological streams (Living Architecture), communicative-alignment-distance for dyadic streams (Coherent Mind).}

\textbf{Requirement.} Domain-native metrics must satisfy the Bias-consistency property of §9.3.1 below to be admissible for the Principle's outperformance statement.

\sectionbreak

\section{§9.3 --- Bias-consistency}

\textbf{Definition 9.3.1 (Bias-consistent metric).} \textit{A metric d on Ω\_\allowbreak{}S is \textbf{Bias-consistent} iff the d-induced trajectory-divergence D\_\allowbreak{}d and the push-operator action on Bias(S) satisfy:}

\[
\begin{aligned}
& \forall\, \text{Stream-morphism } f : S \to S',\ \forall\, I: \\
& \quad D_d(S, I) \leq D_d(S', I) \iff \mathrm{Bias}(S) \succeq_\mathrm{push} \mathrm{Bias}(S') \mid_f
\end{aligned}
\]

\textit{where ⪰\_\allowbreak{}push is the partial order on signed measures induced by applicability of push\_\allowbreak{}struct, push\_\allowbreak{}info transformations from Bias(S') to Bias(S) along f.}

\textbf{Proposition 9.3.2 (Bias-consistency is a natural condition).} \textit{The three canonical classes of §9.2 are Bias-consistent:}

\begin{itemize}
\item \textit{d\_\allowbreak{}W: Bias-consistent by convexity of Wasserstein with respect to couplings and the Bias-push-operator definitions of §7.4.}
\item \textit{d\_\allowbreak{}KL, ε: Bias-consistent for ε small enough (specifically for ε below the minimal Bias-positive-mass-density, which is non-zero under Def 8.1.6).}
\item \textit{d\_\allowbreak{}dom: Bias-consistent by case-specific domain argument. Meridian's energy-distance is Bias-consistent because cosmological conscious-gravity's Bias decomposes via the Killing-form spectral structure (§Meridian Chapter VII); Living Architecture's fitness-distance is Bias-consistent because fitness-gradient and Bias-gradient both descend from the γ-coherence-attraction structure.}
\end{itemize}

\textbf{Proof for d\_\allowbreak{}W and d\_\allowbreak{}{KL, ε}.} Under the pushforward construction of §9.1.3 + the push-operator definitions of §7.4, the inequality D\_\allowbreak{}d(S, I) ≤ D\_\allowbreak{}d(S', I) tracks Bias-push-applicability along f by direct computation. For d\_\allowbreak{}W: Kantorovich duality gives the equivalence explicitly; for d\_\allowbreak{}{KL, ε}: ε-regularized log-likelihood-ratio respects Bias-ordering when ε is below the positive-mass-density floor. ∎

\sectionbreak

\section{§9.4 --- The outperformance inequality}

\subsection{§9.4.1 --- Four stream-parameters}

Before stating the theorem, define the four stream-parameters each condition controls. For a stream S over an interval I = [t₀, t₁] and a Bias-consistent metric d on Ω\_\allowbreak{}S:

\begin{itemize}
\item \textbf{η\_\allowbreak{}sep(S) := μ\_\allowbreak{}Ω(DOF(O\_\allowbreak{}1) ∩ DOF(O\_\allowbreak{}2) ∩ ⋯)} --- the DOF-overlap measure across S's objectives, as a μ\_\allowbreak{}Ω-measurable subset of Ω\_\allowbreak{}S. Under \textbf{C\_\allowbreak{}sep}: η\_\allowbreak{}sep(S) = 0.
\item \textbf{τ\_\allowbreak{}max(S) := sup\_\allowbreak{}k (τ\_\allowbreak{}{k+1} − τ\_\allowbreak{}k)} --- the maximum gap between refresh-events in the measurement partition {τ\_\allowbreak{}k}. Under \textbf{C\_\allowbreak{}meas}: τ\_\allowbreak{}max(S) ≤ T\_\allowbreak{}refresh, a finite framework-specified ceiling.
\item \textbf{δ\_\allowbreak{}scale(S) := sup\_\allowbreak{}{e ∈ E(DAG(S))} d(γ\_\allowbreak{}{child(e)}, γ\_\allowbreak{}{parent(e)})} --- the scale-discontinuity sup across the A2.6 DAG's edges. Under \textbf{C\_\allowbreak{}scale}: δ\_\allowbreak{}scale(S) ≤ ε\_\allowbreak{}scale, a finite smoothness tolerance.
\item \textbf{ρ\_\allowbreak{}dyn(S) := inf\_\allowbreak{}{[a, b] ⊂ I, b − a = τ\_\allowbreak{}dyn} (\# cycles in [a, b]) / τ\_\allowbreak{}dyn} --- the infimum cycle-density across sliding windows of framework-specified length τ\_\allowbreak{}dyn. Under \textbf{C\_\allowbreak{}dyn}: ρ\_\allowbreak{}dyn(S) ≥ ρ\_\allowbreak{}min > 0.
\end{itemize}

Auxiliary constants (metric-dependent, finite under Bias-consistency):

\begin{itemize}
\item \textbf{Λ\_\allowbreak{}γ(S, d)} --- the d-Lipschitz constant of γ\_\allowbreak{}S (supremum of d-distance traversed per unit time by γ\_\allowbreak{}S in the live regime).
\item \textbf{Λ\_\allowbreak{}γ\textasciicircum{}{static}(S, d)} --- the d-drift rate of a γ-frozen stream (supremum d-distance traversed per unit time under stationary γ, representing drift of reality away from the frozen estimate).
\item \textbf{diam\_\allowbreak{}d(Ω\_\allowbreak{}S)} --- the d-diameter of the relevant subset of Ω\_\allowbreak{}S (finite for bounded d; for unbounded d, replace with the Bias-weighted diameter diam\_\allowbreak{}{d, Bias} := sup\_\allowbreak{}{E: Bias\_\allowbreak{}+(E) > 0} sup\_\allowbreak{}{ω, ω' ∈ E} d(ω, ω'), finite under σ-finite Bias per §6.9.7).
\item \textbf{depth(DAG(S))} --- the number of levels in S's A2.6-DAG (finite under §2.2.6's acyclicity + finite-composability).
\item \textbf{N\_\allowbreak{}refresh(S, I) := |{k : τ\_\allowbreak{}k ∈ I}|} --- the count of refresh-events in I (≤ (t₁ − t₀) / τ\_\allowbreak{}min, with τ\_\allowbreak{}min ≥ 0).
\end{itemize}

\subsection{§9.4.2 --- The joint bound}

\textbf{Lemma 9.4.2 (Four-term contribution bound).} \textit{For a stream S over I = [t₀, t₁] and a Bias-consistent metric d:}

\[
\begin{aligned}
\mathbb{E}_I[D_d(S, \cdot)] \;\leq\;\; & \underbrace{\eta_\mathrm{sep}(S) \cdot \mathrm{diam}_{d, \mathrm{Bias}}(\Omega_S) \cdot (t_1 - t_0)}_{B_\mathrm{sep}(S, I)} \\
& {} + \underbrace{\Lambda_\gamma(S, d) \cdot \tau_\mathrm{max}(S) \cdot N_\mathrm{refresh}(S, I)}_{B_\mathrm{meas}(S, I)} \\
& {} + \underbrace{\mathrm{depth}(\mathrm{DAG}(S)) \cdot \delta_\mathrm{scale}(S) \cdot (t_1 - t_0)}_{B_\mathrm{scale}(S, I)} \\
& {} + \underbrace{(1 - \rho_\mathrm{dyn}(S)) \cdot \Lambda_\gamma^\mathrm{static}(S, d) \cdot (t_1 - t_0)}_{B_\mathrm{dyn}(S, I)}.
\end{aligned}
\]

\textbf{Proof.}

\textit{B\_\allowbreak{}sep term.} The integrand d(α\_\allowbreak{}S(t), α*\_\allowbreak{}S(t)) in Def 9.1.1 decomposes along the DOF-product structure of Ω\_\allowbreak{}S (Triple-decomposition compatibility, Cor 7.3.3) as a sum of per-DOF contributions plus cross-DOF contributions. Under DOF-separation (η\_\allowbreak{}sep = 0), cross-DOF contributions vanish; when η\_\allowbreak{}sep > 0, the cross-DOF contribution is bounded pointwise by d-diameter restricted to the DOF-overlap region. Integrating over I and applying Bias-weighting gives η\_\allowbreak{}sep · diam\_\allowbreak{}{d, Bias} · (t₁ − t₀).

\textit{B\_\allowbreak{}meas term.} Between consecutive refresh-events τ\_\allowbreak{}k, τ\_\allowbreak{}{k+1}, the γ-implied trajectory α*\_\allowbreak{}S runs from α\_\allowbreak{}S(τ\_\allowbreak{}k) but is not corrected until τ\_\allowbreak{}{k+1}. The d-drift over [τ\_\allowbreak{}k, τ\_\allowbreak{}{k+1}] is bounded by Λ\_\allowbreak{}γ · (τ\_\allowbreak{}{k+1} − τ\_\allowbreak{}k) by the Lipschitz hypothesis. Summing over k ∈ {1, ..., N\_\allowbreak{}refresh} and using τ\_\allowbreak{}{k+1} − τ\_\allowbreak{}k ≤ τ\_\allowbreak{}max gives Λ\_\allowbreak{}γ · τ\_\allowbreak{}max · N\_\allowbreak{}refresh.

\textit{B\_\allowbreak{}scale term.} The A2.6-DAG introduces potential discontinuities at each edge between scales. For an edge e ∈ E(DAG(S)), the d-contribution from that edge is bounded by d(γ\_\allowbreak{}{child(e)}, γ\_\allowbreak{}{parent(e)}) ≤ δ\_\allowbreak{}scale per unit time. Summing over edges (bounded by depth × width; absorb width into δ\_\allowbreak{}scale's sup) and integrating over I gives depth · δ\_\allowbreak{}scale · (t₁ − t₀).

\textit{B\_\allowbreak{}dyn term.} On sliding windows of length τ\_\allowbreak{}dyn where no propose/dissolve/build cycle occurs, γ is effectively frozen and the stream's α\_\allowbreak{}S drifts from α*\_\allowbreak{}S at rate ≤ Λ\_\allowbreak{}γ\textasciicircum{}static. The fraction of I spent in such frozen windows is (1 − ρ\_\allowbreak{}dyn); integrating gives (1 − ρ\_\allowbreak{}dyn) · Λ\_\allowbreak{}γ\textasciicircum{}static · (t₁ − t₀).

Sum of the four contributions is the claimed bound. ∎

\subsection{§9.4.3 --- The outperformance theorem}

\textbf{Theorem 9.4.3 (Principle outperformance for Bias-consistent D, with explicit constants).} \textit{Let d be any Bias-consistent metric. Let S, S' be comparable streams over I = [t₀, t₁]. Suppose S ∈ coherence-regime, so η\_\allowbreak{}sep(S) = 0, τ\_\allowbreak{}max(S) ≤ T\_\allowbreak{}refresh, δ\_\allowbreak{}scale(S) ≤ ε\_\allowbreak{}scale, ρ\_\allowbreak{}dyn(S) ≥ ρ\_\allowbreak{}min. Suppose S' ∉ coherence-regime, i.e., at least one of:}

\begin{itemize}
\item \textit{(¬C\_\allowbreak{}sep) η\_\allowbreak{}sep(S') > 0;}
\item \textit{(¬C\_\allowbreak{}meas) τ\_\allowbreak{}max(S') > T\_\allowbreak{}refresh;}
\item \textit{(¬C\_\allowbreak{}scale) δ\_\allowbreak{}scale(S') > ε\_\allowbreak{}scale;}
\item \textit{(¬C\_\allowbreak{}dyn) ρ\_\allowbreak{}dyn(S') < ρ\_\allowbreak{}min.}
\end{itemize}

\textit{Then:}

\[
\mathbb{E}_I[D_d(S, \cdot)] \;\leq\; B_\mathrm{coh}(S, I) \;<\; B_\mathrm{coh}(S', I) \;+\; \Delta(S', I) \;\leq\; \mathbb{E}_I[D_d(S', \cdot)]
\]

\textit{where}

\[
\begin{aligned}
B_\mathrm{coh}(S, I) = {} & 0 + \Lambda_\gamma(S) \cdot T_\mathrm{refresh} \cdot N_\mathrm{refresh} \\
& {} + \mathrm{depth}(\mathrm{DAG}(S)) \cdot \varepsilon_\mathrm{scale} \cdot (t_1 - t_0) \\
& {} + (1 - \rho_\mathrm{min}) \cdot \Lambda_\gamma^\mathrm{static}(S) \cdot (t_1 - t_0)
\end{aligned}
\]

\textit{is the coherence-regime ceiling, and Δ(S', I) is the failure-mode gap contributed by whichever condition(s) S' fails:}

\[
\begin{aligned}
\Delta(S', I) = {} & \eta_\mathrm{sep}(S') \cdot \mathrm{diam}_{d, \mathrm{Bias}} \cdot (t_1 - t_0) \\
& {} + \Lambda_\gamma(S') \cdot (\tau_\mathrm{max}(S') - T_\mathrm{refresh})^+ \cdot N_\mathrm{refresh}(S', I) \\
& {} + \mathrm{depth}(\mathrm{DAG}(S')) \cdot (\delta_\mathrm{scale}(S') - \varepsilon_\mathrm{scale})^+ \cdot (t_1 - t_0) \\
& {} + (\rho_\mathrm{min} - \rho_\mathrm{dyn}(S'))^+ \cdot \Lambda_\gamma^\mathrm{static}(S') \cdot (t_1 - t_0).
\end{aligned}
\]

\textit{(x)\textasciicircum{}+ := max(x, 0). The inequality is strict because at least one summand of Δ(S', I) is positive under the ¬C\_\allowbreak{}i hypothesis on S'.}

\textbf{Proof.} Lemma 9.4.2 applied to S with the coherence-regime parameter upper bounds gives E\_\allowbreak{}I[D\_\allowbreak{}d(S)] ≤ B\_\allowbreak{}coh(S, I). Lemma 9.4.2 applied to S' with the parameter-by-parameter shortfall gives E\_\allowbreak{}I[D\_\allowbreak{}d(S')] ≥ B\_\allowbreak{}coh(S', I) + Δ(S', I) by the monotonicity of each B\_\allowbreak{}· term in its controlling parameter. Comparability (§5.0) identifies B\_\allowbreak{}coh(S, I) ≈ B\_\allowbreak{}coh(S', I) modulo the stream-intrinsic auxiliary constants (Λ\_\allowbreak{}γ, depth(DAG), etc.) which differ by O(1) between comparable streams. The strict part follows from Δ(S', I) > 0 under the failure hypothesis. Bias-consistency of d preserves the ordering across any valid metric choice (Thm 9.5.1). ∎

\textbf{Remark 9.4.4 (Quantitative reading).} The theorem is now quantitative: given values for the coherence-regime parameters (T\_\allowbreak{}refresh, ε\_\allowbreak{}scale, ρ\_\allowbreak{}min) and the stream-intrinsic constants (Λ\_\allowbreak{}γ, diam, depth), B\_\allowbreak{}coh(S, I) is computable. This is the referee-grade tightness the sketch-proof lacked. For specific domains, the constants take concrete values: for Meridian's cosmological F\_\allowbreak{}Mrd, T\_\allowbreak{}refresh = one refresh-timescale per Hubble time; for Living Architecture's biological F\_\allowbreak{}LA, T\_\allowbreak{}refresh = one generation-cycle; etc. The Principle's empirical content is the sign of E[D\_\allowbreak{}d(S')] − B\_\allowbreak{}coh(S, I) --- and now the \textit{magnitude} of that sign is explicitly Δ(S', I).

\textbf{Remark 9.4.5 (Relation to §5.2.5).} §5.2.5 established independence of the four conditions by counterexample --- each condition contributes; none is redundant. Thm 9.4.3 strengthens this to \textbf{quantitative joint sufficiency}: B\_\allowbreak{}coh is the sum of the four ceilings, and the shortfall Δ(S') is the sum of the four per-condition violations. Independence at §5.2.5 is witnessed by the fact that each B\_\allowbreak{}· term can independently exceed zero for S' while the others are tight.

⚑ [SURFACED | Companion §9.4.3 | → Anchor §9.1/§9.3 target | type: theorem (quantitative)]
--- Outperformance with explicit constants (η\_\allowbreak{}sep, τ\_\allowbreak{}max, δ\_\allowbreak{}scale, ρ\_\allowbreak{}dyn) + coherence-regime ceiling B\_\allowbreak{}coh + failure-mode gap Δ.

\sectionbreak

\section{§9.5 --- Cross-metric invariance (resolution of Q1)}

\textbf{Theorem 9.5.1 (Cross-metric invariance of the outperformance ordering).} \textit{For any two Bias-consistent metrics d, d' on Ω\_\allowbreak{}S, the outperformance ordering}

\[
\mathbb{E}[D_d(S)] < \mathbb{E}[D_d(S')]
\]

\textit{holds iff the same ordering holds for d':}

\[
\mathbb{E}[D_{d'}(S)] < \mathbb{E}[D_{d'}(S')].
\]

\textbf{Proof.} Both orderings are equivalent to \textit{Bias(S) ⪰\_\allowbreak{}push Bias(S') (modulo comparability)} by Definition 9.3.1. Hence they are equivalent to each other. ∎

\textbf{Corollary 9.5.2 (Anchor §9.9 Q1 resolved).} \textit{The outperformance ordering of Theorem 5.1.2 is invariant across any Bias-consistent metric. The Principle's empirical content does not depend on a specific metric choice, provided the chosen metric is Bias-consistent.}

⚑ [SURFACED | Companion §9.5 | → Anchor §9.9 Q1 target | type: theorem (resolution)]

\textbf{Remark 9.5.3 (Meridian/domain-native preference).} In practice, domain-native metrics are preferred over Wasserstein or KL because they expose more signal per data-point (d\_\allowbreak{}dom aligns with the domain's natural observables). Bias-consistency is the admissibility criterion; once it holds, metric choice is a practical matter.

\sectionbreak

\section{§9.6 --- D on F\_\allowbreak{}∞ (self-reference application)}

Applying §9 to F\_\allowbreak{}∞ (§8's self-reference construction):

\textbf{Proposition 9.6.1 (D\_\allowbreak{}d on F\_\allowbreak{}∞).} \textit{For any Bias-consistent d on Ω\_\allowbreak{}F, D\_\allowbreak{}d(F\_\allowbreak{}∞, [t₀, t₁]) is well-defined and finite over any closed construction interval [t₀, t₁].}

\textbf{Proof.} §8.1 specifies F\_\allowbreak{}∞ as an extensional F-coalgebra, §9.1.2 gives D\_\allowbreak{}d well-definedness, §9.2's candidates are each Bias-consistent. Finiteness: the construction interval is finite; Bias(F\_\allowbreak{}∞) is σ-finite; d is uniformly bounded by the Ω\_\allowbreak{}F-diameter times interval-length. ∎

\textbf{Proposition 9.6.2 (F\_\allowbreak{}∞ outperforms comparable non-coherent F' --- conditional).} \textit{Conditional on external execution of Prop 8.5.2 affirming Audit Observation 8.3.5, for any comparable F' framework-construction process not in coherence-regime:}

\[
\mathbb{E}[D_d(F_\infty)] < \mathbb{E}[D_d(F')]
\]

\textit{for any Bias-consistent d.}

\textbf{Proof.} Under the conditional hypothesis, Audit Observation 8.3.5 upgrades to a theorem placing F\_\allowbreak{}∞ in coherence-regime; Thm 9.4.3 then applies with S = F\_\allowbreak{}∞ and S' = F'. ∎

\textbf{Remark 9.6.2′ (Unconditional D-computation).} The left-hand side D\_\allowbreak{}d(F\_\allowbreak{}∞, I) is well-defined and computable independent of the conditional status (Prop 9.6.1). The \textit{outperformance ordering} is what the conditional hypothesis gates. Hence the empirical magnitude-study proposed in Rem 9.6.3 can proceed now; only the interpretation of the magnitude as "outperformance by a framework in coherence-regime" awaits the external audit.

\textbf{Remark 9.6.3 (Empirical magnitude).} §8.3's audit establishes the \textit{direction} of Proposition 9.6.2's outperformance. The \textit{magnitude} requires running D\_\allowbreak{}d against a specific comparable F' instance. Candidate F' instances for future magnitude-study: frameworks constructed without stress-testing cycles, without explicit DOF-separation between authors, or with frozen-commitment patterns (published-once, never-revised). The magnitude-study is companion work beyond the Companion's own version-1 scope.

\sectionbreak

\section{§9.7 --- Observable signatures recap}

Following Anchor §9.3, three operational signatures:

\begin{enumerate}
\item \textbf{Trajectory-tracking.} Sample α\_\allowbreak{}S(t) at refresh-rate; reconstruct γ\_\allowbreak{}S from prior-data; compute D\_\allowbreak{}d directly.
\item \textbf{Adjoint-composition success rate.} Count successful ι ⊣ κ compositions in 𝒞\_\allowbreak{}LDS per interval; coherent streams show higher success rate.
\item \textbf{Multi-scale coherence correlation.} For nested S₁ ⊂ S₂ in the A2.6 DAG, correlate child-γ and parent-γ; positive correlation is the coherence-regime signature.
\end{enumerate}

Each signature reduces to a D\_\allowbreak{}d computation under a specific metric choice and projection:

\begingroup\setlength{\tabcolsep}{3pt}\renewcommand{\arraystretch}{1.15}\footnotesize
\begin{center}
\begin{tabular}{p{1.72in} p{1.40in} p{1.21in}}
\toprule
Signature & Metric choice & Projection \\
\midrule
Trajectory-tracking & d\_\allowbreak{}dom domain-native & α\_\allowbreak{}S projection \\
Adjoint-composition success & d\_\allowbreak{}discrete counting & 𝒞\_\allowbreak{}LDS projection \\
Multi-scale coherence corr. & d\_\allowbreak{}W on γ-distributions & DAG-edge projection \\
\bottomrule
\end{tabular}
\end{center}
\endgroup
\sectionbreak

\section{§9.8 --- Open questions for §9}

\begin{itemize}
\item \textbf{Q1 (resolved):} cross-metric invariance --- Theorem 9.5.1 closes this.
\item \textbf{Q2 (regime-boundary topology; from §5.6):} Is coherence-regime an open condition in an appropriate topology on 𝒞\_\allowbreak{}Streams? Under the D\_\allowbreak{}d-induced topology: yes, by continuity of D\_\allowbreak{}d + upper-semi-continuity of the joint-condition conjunction. Detailed proof carry-forward.
\item \textbf{Q3 (magnitude-scaling on F\_\allowbreak{}∞):} Empirical magnitude of Proposition 9.6.2 with F' instances as candidate-above.
\item \textbf{Q4 (non-Bias-consistent metrics):} What is the structural content of metrics that fail Bias-consistency? Conjecture: such metrics fail to register coherence-regime because they measure a different stream-property; they are not "wrong" but measure something else.
\end{itemize}

\sectionbreak

\section{§9.9 --- Forward-pointers}

\begin{itemize}
\item \textbf{§10} (reference figures): Figures 9.1 (D trajectory-divergence diagrammatic), 9.2 (four-conditions schematic with D), 9.3 (self-reference closure graph) per SCOPE §4's §10 TikZ-reference list.
\item \textbf{Appendix B} (anchor): trajectory-divergence metric reference cross-walks into Companion §9 from anchor Appendix B §B.5.
\item \textbf{Domain volumes:} Meridian/Living Architecture/Coherent Body/etc. each specify d\_\allowbreak{}dom-Bias-consistency in their own domain-specific way.
\end{itemize}

\sectionbreak

\section{§9.10 --- Surfaced-lemma register}

Four flags surface this pass:

\begin{itemize}
\item ⚑ §9.3.1 Bias-consistency as admissibility criterion for metrics → Anchor §9.3 target --- definition (promoted from informal criterion to formal property)
\item ⚑ §9.4.1 Four stream-parameters (η\_\allowbreak{}sep, τ\_\allowbreak{}max, δ\_\allowbreak{}scale, ρ\_\allowbreak{}dyn) + auxiliary constants → Anchor §9.3 target --- definitions (quantitative stream-parameters controlling each condition)
\item ⚑ §9.4.3 Outperformance with explicit constants: B\_\allowbreak{}coh ceiling + Δ(S') shortfall → Anchor §9.3/§9.1 target --- theorem (quantitative form supersedes sketch)
\item ⚑ §9.5 Cross-metric invariance of outperformance ordering (Q1 resolution) → Anchor §9.9 Q1 target --- theorem (resolves the open question)
\end{itemize}

\sectionbreak


\chapter{§10 — Reference Figures}


\textit{The Companion is the figure authority (SCOPE §9). This chapter provides canonical TikZ source for figures used across the Anchor and Companion. The Anchor imports from here; domain volumes cite by figure-number. One source of truth.}

\sectionbreak

\section{§10.0 --- Figure index}

Eight figures in the canonical standard set, grouped by chapter-of-authority in the Companion:

\begingroup\setlength{\tabcolsep}{3pt}\renewcommand{\arraystretch}{1.15}\footnotesize
\begin{center}
\begin{tabular}{p{0.56in} p{0.67in} p{1.06in} p{2.00in}}
\toprule
Figure & Authority section & Anchor cross-ref & Topic \\
\midrule
Fig 1 & §1.2 /\allowbreak  §6.1 & Anchor §3, §7 & The endofunctor F and Stream-as-F-coalgebra \\
Fig 2 & §6.2 & Anchor §1 (Fig 1.1) & Triple functor T : 𝒞\_\allowbreak{}Streams →\allowbreak  𝒞\_\allowbreak{}Triple \\
Fig 3 & §6.3 & Anchor §1 (Fig 1.2) & Recursive decomposability T\textasciicircum{}(n) \\
Fig 4 & §6.4 & Anchor §3 & Kind-classifier fibration π : Stream →\allowbreak  ContentIndex \\
Fig 5 & §7.3 & Anchor Appendix B (Fig B.1) & Bias(S) signed measure + push-operators \\
Fig 6 & §3.4.2 & Anchor §7 & Dual coherence axes (σ\_\allowbreak{}struct × σ\_\allowbreak{}info) \\
Fig 7 & §5.2 & Anchor §9 (Fig 9.2) & Four-conditions schematic \\
Fig 8 & §8 /\allowbreak  §5.4 & Anchor §9 (Fig 9.3) & Self-reference closure \\
\bottomrule
\end{tabular}
\end{center}
\endgroup
The Anchor README lists fourteen figures --- the remaining six (mismatch condition, kind stratification directed diagram, ι⊣κ adjunction detail, push-operators internal diagram, kind-demotion dynamic trajectory, trajectory divergence D(S) with envelope) are domain-specific overlays of Figures 1–8 and are back-ported into the Anchor directly. Companion §10 owns the eight canonical figures; the overlays live in the Anchor as specializations.

\sectionbreak

\section{§10.1 --- Figure 1: The endofunctor F and Stream-as-F-coalgebra}

\textbf{Purpose.} Depicts F : σ ↦ σ\textasciicircum{}(C\textasciicircum{}op) as a mixed-variance endofunctor on adequate F-coalgebra-carriers, with γ as a coalgebra-structure map.

% Companion §10, Fig 1
% F endofunctor and Stream-as-F-coalgebra
\begin{figure}[htbp]
\centering
\begin{tikzcd}[column sep=large, row sep=large]
  \sigma \arrow[r, "\gamma"] \arrow[d, "f_\sigma"'] & F(\sigma) = \sigma^{C^{\mathrm{op}}} \arrow[d, "{F(f_\sigma, f_C)}"] \\
  \sigma' \arrow[r, "\gamma'"'] & F(\sigma') = (\sigma')^{(C')^{\mathrm{op}}}
\end{tikzcd}
\caption{The F-coalgebra square. Each row is a stream $S = (\sigma, C, \gamma)$ (top) and $S' = (\sigma', C', \gamma')$ (bottom). A Stream-morphism $f = (f_\sigma, f_C) : S \to S'$ satisfies the coalgebra-commute: $F(f_\sigma, f_C) \circ \gamma = \gamma' \circ f_\sigma$. The endofunctor $F : \sigma \mapsto \sigma^{C^{\mathrm{op}}}$ is mixed-variance (covariant in $\sigma$, contravariant in $C$).}
\label{fig:f-coalgebra-square}
\end{figure}

\sectionbreak

\section{§10.2 --- Figure 2: Triple functor T}

\textbf{Purpose.} Depicts T : 𝒞\_\allowbreak{}Streams → 𝒞\_\allowbreak{}Triple as a colax-limit-compatible decomposition into Form, Content, Carrier.

% Companion §10, Fig 2
% Triple functor (canonical replacement for Anchor Fig 1.1)
\begin{figure}[htbp]
\centering
\begin{tikzcd}[column sep=huge, row sep=large]
  & \mathcal{C}_\mathrm{Streams}
    \arrow[dl, "\Phi\ (\mathrm{Form})"']
    \arrow[d, "T"]
    \arrow[dr, "\Psi\ (\mathrm{Content})"]
    \arrow[dd, bend left=45, "\mathrm{K}\ (\mathrm{Carrier})"] & \\
  \mathcal{C}_\mathrm{Form}
    & \mathcal{C}_\mathrm{Form} \times \mathbf{Cat}_\mathrm{small} \times \mathbf{Carrier}
      \arrow[l]
      \arrow[r]
      \arrow[d]
    & \mathbf{Cat}_\mathrm{small} \\
  & \mathbf{Carrier} &
\end{tikzcd}
\caption{The Triple functor $T : \mathcal{C}_\mathrm{Streams} \to \mathcal{C}_\mathrm{Triple}$ with $\mathcal{C}_\mathrm{Triple} = \mathcal{C}_\mathrm{Form} \times \mathbf{Cat}_\mathrm{small} \times \mathbf{Carrier}$. The three factor functors $\Phi, \Psi, \mathrm{K}$ recover Form, Content (ContentOp-category), and Carrier (coalgebra-structure-map) respectively. Under the initial-object hypothesis of §6.6, $T$ is a colax limit (the natural transformations between compositions of the factors are in general not invertible).}
\label{fig:triple-functor}
\end{figure}

\sectionbreak

\section{§10.3 --- Figure 3: Recursive decomposability T\textasciicircum{}(n)}

\textbf{Purpose.} Depicts the finite-depth iterated Triple. Each Triple-component is itself a stream admitting its own Triple (Lemma 6.3.2).

% Companion §10, Fig 3
% Recursive decomposability, finite-depth
\begin{figure}[htbp]
\centering
\begin{tikzcd}[column sep=small, row sep=small]
  & & S \arrow[dll] \arrow[d] \arrow[drr] & & \\
  \mathrm{Form}(S) \arrow[d] \arrow[drr]
    & & \mathrm{Content}(S) \arrow[dll] \arrow[d] \arrow[drr]
    & & \mathrm{Carrier}(S) \arrow[dll] \arrow[d] \\
  \mathrm{F}\mathrm{F} & \mathrm{F}\mathrm{C} & \mathrm{F}\mathrm{K} \cdots \mathrm{C}\mathrm{F} & \mathrm{C}\mathrm{C} & \mathrm{C}\mathrm{K} \cdots \mathrm{K}\mathrm{F} \cdots \mathrm{K}\mathrm{K}
\end{tikzcd}
\caption{Recursive Triple-decomposition at depth 2. Each of $\mathrm{Form}(S)$, $\mathrm{Content}(S)$, $\mathrm{Carrier}(S)$ is itself a stream, so the Triple functor applies again, yielding nine depth-2 components (abbreviated: FF = $\mathrm{Form}(\mathrm{Form}(S))$, FC = $\mathrm{Content}(\mathrm{Form}(S))$, etc.). Finite-depth iteration is well-defined for every adequate F-coalgebra (Lemma 6.3.2). Full $\omega$-depth requires additional hypotheses (H1, H2; §6.9).}
\label{fig:recursive-decomposability}
\end{figure}

\sectionbreak

\section{§10.4 --- Figure 4: Kind-classifier fibration π}

\textbf{Purpose.} Depicts the bicategorical fibration π : Stream → ContentIndex with kind-preorder base.

% Companion §10, Fig 4
% Kind-classifier fibration
\begin{figure}[htbp]
\centering
\begin{tikzcd}[column sep=small, row sep=large]
  \mathrm{Stream} \arrow[d, "\pi"] & \\
  \mathbf{ContentIndex}
    & K \arrow[ul, dashed, bend right=25, "{\mathrm{fiber}(\pi, K)}"']
\end{tikzcd}
\\[1em]
\begin{tikzcd}[column sep=large, row sep=large]
  \mathrm{reactive} \arrow[r, "\sqsubseteq"] & \mathrm{self\text{-}maint.} \arrow[r, "\sqsubseteq"] & \mathrm{self\text{-}ref.} \arrow[r, "\sqsubseteq"] & \mathrm{abstr.}
\end{tikzcd}
\caption{Kind-classifier fibration $\pi : \mathrm{Stream} \to \mathbf{ContentIndex}$ (top). The fiber over a kind $K$ is the full subcategory of $K$-streams. The kind-preorder base (bottom) stratifies: reactive $\sqsubseteq$ self-maintaining $\sqsubseteq$ self-referential $\sqsubseteq$ abstractive. Under Q7 (lattice-where-possible), $\pi$ is a strict fibration; generically bicategorical. The unifying-stream subfibration $\mathrm{Stream}_u \subset \mathrm{Stream}$ (§6.4.1) consists of streams with terminal-object-preserving ContentOp-refinements.}
\label{fig:kind-classifier-fibration}
\end{figure}

\sectionbreak

\section{§10.5 --- Figure 5: Bias(S) signed measure + push-operators}

\textbf{Purpose.} Depicts Bias(S) as a signed measure decomposing via Hahn-Jordan, with push\_\allowbreak{}struct and push\_\allowbreak{}info as non-commuting transformations.

% Companion §10, Fig 5
% Bias(S) signed measure with push-operators
\begin{figure}[htbp]
\centering
\begin{tikzpicture}[scale=0.85]
  % Ω_S axis (horizontal)
  \draw[->] (-0.5, 0) -- (10.5, 0) node[right] {$\Omega_S$};
  \draw[->] (0, -2.5) -- (0, 3) node[above] {$\mathrm{Bias}(S)$};
  % Positive part
  \fill[blue!30, domain=1:4, smooth] plot (\x, {2*exp(-(\x-2.5)^2/1.0)}) -- (4, 0) -- (1, 0) -- cycle;
  \draw[blue, thick, domain=1:4, smooth] plot (\x, {2*exp(-(\x-2.5)^2/1.0)});
  \node[blue] at (2.5, 2.3) {$\mathrm{Bias}_+$};
  % Negative part
  \fill[red!30, domain=6:9, smooth] plot (\x, {-1.5*exp(-(\x-7.5)^2/1.2)}) -- (9, 0) -- (6, 0) -- cycle;
  \draw[red, thick, domain=6:9, smooth] plot (\x, {-1.5*exp(-(\x-7.5)^2/1.2)});
  \node[red] at (7.5, -1.8) {$\mathrm{Bias}_-$};
  % Push arrows
  \draw[->, thick, green!60!black] (2.5, 2.5) .. controls (3.5, 3.2) .. (4.5, 2.8);
  \node[green!60!black] at (4.0, 3.3) {$\mathrm{push}_\mathrm{struct}$};
  \draw[->, thick, orange] (2.5, 0.5) .. controls (3.0, 1.5) .. (3.2, 2.0);
  \node[orange] at (3.8, 1.0) {$\mathrm{push}_\mathrm{info}$};
\end{tikzpicture}
\caption{Bias(S) as a signed measure on $\Omega_S$. Hahn-Jordan decomposition: $\mathrm{Bias}(S) = \mathrm{Bias}_+ - \mathrm{Bias}_-$. The push-operators $\mathrm{push}_\mathrm{struct}$ and $\mathrm{push}_\mathrm{info}$ transform Bias along structural-coherence and informational-coherence axes respectively. Their non-commutator $[\mathrm{push}_\mathrm{struct}, \mathrm{push}_\mathrm{info}] \neq 0$ is the formal content of Proposition 7.4.3.}
\label{fig:bias-signed-measure}
\end{figure}

\sectionbreak

\section{§10.6 --- Figure 6: Dual coherence axes (σ\_\allowbreak{}struct × σ\_\allowbreak{}info plane)}

\textbf{Purpose.} Depicts the T6 orthogonal axes and the codimension-2 dually-coherent locus.

% Companion §10, Fig 6
% Dual coherence axes
\begin{figure}[htbp]
\centering
\begin{tikzpicture}[scale=1.1]
  \draw[->] (0, 0) -- (5, 0) node[right] {$\sigma_\mathrm{struct}$};
  \draw[->] (0, 0) -- (0, 5) node[above] {$\sigma_\mathrm{info}$};
  \draw[dashed] (4, 0) -- (4, 4);
  \draw[dashed] (0, 4) -- (4, 4);
  \node at (4.2, 4.2) {\small $(1,1)$ dual};
  \fill[black] (4, 4) circle (3pt);
  % Example streams
  \fill[blue] (3.5, 1) circle (2pt);
  \node[blue] at (3.5, 0.6) {\small T5-only};
  \fill[red] (0.8, 3.5) circle (2pt);
  \node[red] at (1.4, 3.5) {\small T6-only};
  \fill[gray] (1, 1) circle (2pt);
  \node[gray] at (1, 0.6) {\small neither};
  % Codimension-2 region
  \draw[thick, green!60!black] (4, 4) circle (0.4);
  \node[green!60!black] at (3.3, 3.3) {\small locus};
\end{tikzpicture}
\caption{The $\sigma_\mathrm{struct} \times \sigma_\mathrm{info}$ plane of T6 (Theorem 3.4.2). The two axes are orthogonal: T5-only streams (blue) are $\Phi$-fixed-points without entropy-minimality; T6-only streams (red) are entropy-minimal without $\Phi$-fixed-points; dually-coherent streams occupy the codimension-2 locus near $(1, 1)$ (green circle). The locus is non-empty only when $C$ has enough symmetry to admit both a harmonic and a concentrated $\gamma$ simultaneously.}
\label{fig:dual-coherence-axes}
\end{figure}

\sectionbreak

\section{§10.7 --- Figure 7: Four-conditions schematic}

\textbf{Purpose.} Depicts the four conditions as a unified schematic with derivation-sources.

% Companion §10, Fig 7
% Four-conditions schematic (canonical replacement for Anchor Fig 9.2)
\begin{figure}[htbp]
\centering
\begin{tikzpicture}[
  node distance=1.4cm,
  every node/.style={align=center, font=\small}
]
  \node (root) [rectangle, draw, rounded corners, fill=gray!10] {Stream $S$ in coherence-regime over $[t_0, t_1]$};
  \node (c1) [below=1.5cm of root, rectangle, draw, rounded corners, fill=blue!10, xshift=-5cm] {\textbf{C\_sep}\\DOF-separation\\$\{ \mathrm{DOF}(O_1), \mathrm{DOF}(O_2) \}$\\non-overlapping};
  \node (c2) [below=1.5cm of root, rectangle, draw, rounded corners, fill=green!10, xshift=-1.7cm] {\textbf{C\_meas}\\Refresh-rate\\measurement\\$M_k$ at each $\tau_k$};
  \node (c3) [below=1.5cm of root, rectangle, draw, rounded corners, fill=orange!10, xshift=1.7cm] {\textbf{C\_scale}\\Multi-scale\\$\gamma$-continuity\\$\gamma_S \simeq \gamma_{S^\Uparrow}, \gamma_{S^\Downarrow}$};
  \node (c4) [below=1.5cm of root, rectangle, draw, rounded corners, fill=red!10, xshift=5cm] {\textbf{C\_dyn}\\Oscillatory\\maintenance\\build $\to$ dissolve $\to$ build};
  \draw[->] (root) -- (c1);
  \draw[->] (root) -- (c2);
  \draw[->] (root) -- (c3);
  \draw[->] (root) -- (c4);
  \node (d1) [below=0.7cm of c1, font=\scriptsize] {T3 + A2.4};
  \node (d2) [below=0.7cm of c2, font=\scriptsize] {T4};
  \node (d3) [below=0.7cm of c3, font=\scriptsize] {A2.6 + A3.3};
  \node (d4) [below=0.7cm of c4, font=\scriptsize] {T4 + A3.4};
  \node (out) [below=3.7cm of root, rectangle, draw, rounded corners, fill=yellow!20] {Outperformance: $\mathbb{E}[D(S)] < \mathbb{E}[D(S')]$\\for comparable $S'$ not in coherence-regime};
  \draw[->, thick] (d1) |- (out);
  \draw[->, thick] (d2) -- (out);
  \draw[->, thick] (d3) -- (out);
  \draw[->, thick] (d4) |- (out);
\end{tikzpicture}
\caption{The four conditions of the Coherence Principle (Definition 5.1.1, Propositions 5.2.1–5.2.4). Each condition is derived from an axiom/theorem clause (italic text below each box). The joint sufficiency (Prop 5.2.5) is the downstream arrows to the outperformance inequality (Theorem 5.1.2).}
\label{fig:four-conditions}
\end{figure}

\sectionbreak

\section{§10.8 --- Figure 8: Self-reference closure}

\textbf{Purpose.} Depicts F\_\allowbreak{}∞ as a stream with the four-conditions audit + the closure-claim (canonical replacement for Anchor Fig 9.3).

% Companion §10, Fig 8
% Self-reference closure
\begin{figure}[htbp]
\centering
\begin{tikzpicture}[
  node distance=1.2cm,
  every node/.style={align=center, font=\small}
]
  \node (F) [rectangle, draw, rounded corners, thick, fill=gray!10]
    {Construction process $F_\infty$\\\scriptsize (produces this book)\\$F_\infty = (\sigma_F, C_F, \Omega_F, \gamma_F)$\\\scriptsize $K_F = \mathrm{abstr.}$};
  \node (c1) [below=1.3cm of F, rectangle, draw, fill=blue!10, xshift=-4.2cm] {C\_sep\\Clayton $\perp$ Clawd\\DOF};
  \node (c2) [below=1.3cm of F, rectangle, draw, fill=green!10, xshift=-1.4cm] {C\_meas\\Stamp-events\\at axiom /\\theorem / chapter};
  \node (c3) [below=1.3cm of F, rectangle, draw, fill=orange!10, xshift=1.4cm] {C\_scale\\axioms $\leftrightarrow$\\theorems $\leftrightarrow$\\corollaries};
  \node (c4) [below=1.3cm of F, rectangle, draw, fill=red!10, xshift=4.2cm] {C\_dyn\\propose $\to$\\stress $\to$\\reformulate};
  \draw[->] (F) -- (c1);
  \draw[->] (F) -- (c2);
  \draw[->] (F) -- (c3);
  \draw[->] (F) -- (c4);
  \node (out) [below=2.0cm of F, yshift=-1.6cm, rectangle, draw, rounded corners, fill=yellow!20]
    {Output: the framework itself $= \alpha^*_F(t_1)$\\reached by $\gamma_F$-fidelity of $F_\infty$};
  \draw[->, thick] (c1) |- (out);
  \draw[->, thick] (c2) -- (out);
  \draw[->, thick] (c3) -- (out);
  \draw[->, thick] (c4) |- (out);
  \node (closure) [below=0.8cm of out, font=\itshape]
    {"The Coherence Principle is true of frameworks\\that discover the Coherence Principle."};
  \draw[->, thick, dashed] (out) -- (closure);
  \node (F6) [below=0.6cm of closure, rectangle, draw, dashed, fill=red!5, font=\scriptsize]
    {F6 meta-falsification: audit the construction record.\\Construction-record artifacts are public (GitHub, Zenodo).};
  \draw[->, dashed] (closure) -- (F6);
\end{tikzpicture}
\caption{Self-reference closure (conditional on external audit). The construction process $F_\infty$ is itself a stream; under internal audit it satisfies the four conditions over the construction interval (Audit Observation 8.3.5). Conditional on external execution of Prop 8.5.2 affirming the audit, Theorem 9.4.1 then yields outperformance over comparable non-coherent constructions. Derivation-non-circularity (Prop 8.4.1) is established; audit-independence (Prop 8.3.5') is what the external-audit gate discharges.}
\label{fig:self-reference-closure}
\end{figure}

\sectionbreak

\section{§10.9 --- Back-port policy}

\textbf{Anchor import.} The Anchor imports its §9 figures (Figures 7, 8, and the trajectory-divergence figure) from Companion §10 sources via fenced ```latex TikZ blocks carried through the Anchor compile pipeline. One source of truth is maintained.

\textbf{Domain-volume citations.} Domain volumes cite by Companion figure-number:

\begin{itemize}
\item Meridian: Fig 5 (Bias) for cosmological conscious-gravity Bias visualization.
\item Living Architecture: Fig 4 (fibration) for kingdom-stratified stream-kind.
\item Coherent Body/Mind/Continuity: Fig 3 (recursive decomp.) for multi-carrier Triple.
\item Dynamic Organization: Fig 7 (four conditions) for institutional coherence-regime.
\end{itemize}

\textbf{Figure back-port batch rhythm.} SCOPE §9 fixes one source of truth. Any new figure surfacing in a domain volume that generalizes beyond the domain is back-ported to Companion §10 as Figure 9, 10, ... and the domain-volume thereafter cites Companion by figure-number.

\sectionbreak

\section{§10.10 --- Forward-pointers}

\begin{itemize}
\item \textbf{Anchor import (landed).} The TikZ back-port from Companion §10 sources is active in Anchor §9.
\item \textbf{Domain volumes:} see back-port policy above for specific expected uses.
\end{itemize}

\sectionbreak

\section{§10.11 --- Notes on LaTeX environment}

The figures assume:

\begin{itemize}
\item \texttt{tikz} and \texttt{tikz-cd} packages loaded.
\item Font: Companion's document class (default LaTeX or \texttt{book} class sufficient).
\item Math fonts consistent with the Companion's body (see §1 for notation conventions).
\end{itemize}

Each figure is a self-contained \texttt{\textbackslash\{\}begin\{figure\}...\textbackslash\{\}end\{figure\}} block with a \texttt{\textbackslash\{\}label\{\}} providing the canonical cross-reference name (\texttt{fig:triple-functor}, \texttt{fig:four-conditions}, etc.). Citations in the Companion and Anchor use \texttt{\textbackslash\{\}ref\{fig:...\}} directly.

\sectionbreak


\chapter{Appendix A — Anchor → Companion crosswalk}


\textit{Per SCOPE §10 policy: every "}Coherent Structure\textit{" citation in the Anchor (}The Coherence Principle\textit{) resolves to a specific Companion section. This appendix is the canonical resolver. Citations from the Anchor are listed by source location; each row gives the Companion target, the nature of the material (definition / theorem / proof / figure / etc.), and a one-line description.}

\sectionbreak

\section{A.1 --- By anchor chapter}

\subsection{Anchor §1.0 (The Category of Streams) → Companion}

\begingroup\setlength{\tabcolsep}{3pt}\renewcommand{\arraystretch}{1.15}\footnotesize
\begin{center}
\begin{tabular}{p{1.08in} p{0.74in} p{0.57in} p{1.88in}}
\toprule
Anchor location & Companion target & Type & Description \\
\midrule
§1.0 definition of 𝒞\_\allowbreak{}Str & Companion §1.2.1, §6.1.1 & definition & Full CT definition of Stream as adequate F-coalgebra category \\
§1.0 cooperative-constituency ι ⊣ κ & Companion §1.2.3 & definition & Left and right adjoint structure \\
§1.0 five structural properties & Companion §6.1.7 & proposition & Conservativity + limit/\allowbreak colimit list \\
§1.0 γ-naturality & Companion §1.4.3 & proposition & Naturality-of-ν square \\
§1.0 Triple as principal functor & Companion §6.2, §1.7.2 & construction & Full Triple-functor definition \\
\bottomrule
\end{tabular}
\end{center}
\endgroup
\subsection{Anchor §1 (Identity-Trajectory Triple) → Companion}

\begingroup\setlength{\tabcolsep}{3pt}\renewcommand{\arraystretch}{1.15}\footnotesize
\begin{center}
\begin{tabular}{p{1.40in} p{1.12in} p{0.56in} p{1.20in}}
\toprule
Anchor location & Companion target & Type & Description \\
\midrule
§1 Triple definition & Companion §6.2.1, §6.2.2 & definition & T : 𝒞\_\allowbreak{}Streams →\allowbreak  𝒞\_\allowbreak{}Triple \\
§1 Form/\allowbreak Content/\allowbreak Carrier orthogonal-but-constrained & Companion §6.2.7 & proposition & Conservativity of T \\
§1 recursive decomposability & Companion §6.3 + Lemma 6.3.2 & theorem & Finite-depth factorability \\
§1 four carrier-levels (Clawd worked example) & Companion §8.1.1 (referenced only) & definition & Formal σ\_\allowbreak{}F multiplex construction \\
§1 colax-limit statement & Companion §6.6 (Thm 6.6.2) & theorem & Colax-limit under initial-object hypothesis \\
§1 mismatch condition & Companion §6.6 (remark inside Thm 6.6.2) & remark & Obstruction to strict product \\
§1 Figure 1.1 & Companion §10 Fig 2 & figure & Canonical Triple-functor TikZ \\
§1 Figure 1.2 & Companion §10 Fig 3 & figure & Canonical recursive-decomposability TikZ \\
\bottomrule
\end{tabular}
\end{center}
\endgroup
\subsection{Anchor §2 (Axiom 1) → Companion}

\begingroup\setlength{\tabcolsep}{3pt}\renewcommand{\arraystretch}{1.15}\footnotesize
\begin{center}
\begin{tabular}{p{1.24in} p{0.98in} p{0.59in} p{1.46in}}
\toprule
Anchor location & Companion target & Type & Description \\
\midrule
A1.1 non-reducibility & Companion §2.1.1 & axiom-clause & Non-reducibility formal statement \\
A1.2 non-factoring & Companion §2.1.2 & axiom-clause & Functor non-factoring \\
A1.3 all-potentials-realized & Companion §2.1.3 & axiom-clause & Configurational completeness \\
A1.4 substrate-completeness & Companion §2.1.4, §1.5 & axiom-clause & Substrate ⇒ stream-existence \\
\bottomrule
\end{tabular}
\end{center}
\endgroup
\subsection{Anchor §3 (Axiom 2) → Companion}

\begingroup\setlength{\tabcolsep}{3pt}\renewcommand{\arraystretch}{1.15}\footnotesize
\begin{longtable}{p{1.01in} p{0.79in} p{0.58in} p{1.90in}}
\toprule
Anchor location & Companion target & Type & Description \\
\midrule
\endhead
\bottomrule
\endfoot
A2 nested-streams preamble & Companion §2.2 & axiom-cluster & Seven clauses overview \\
\midrule
A2.1 stream-as-F-coalgebra & Companion §2.2.1, §1.6.5 & axiom-clause & F-coalgebra specification \\
\midrule
A2.2 kind-stratification & Companion §2.2.2, §6.4.2 & axiom-clause & Kind preorder + Theorem 2.2.8 (derivation from C-richness) \\
\midrule
A2.3 experience = navigation & Companion §2.2.3, §1.3.1 & axiom-clause & N-orbit construction \\
\midrule
A2.4 cooperative-constituency & Companion §2.2.4 & axiom-clause & ι ⊣ κ properties \\
\midrule
A2.5 kind-refinement under ι & Companion §2.2.5 & axiom-clause & Kind-preservation under cooperative-constituency \\
\midrule
A2.6 DAG-nesting & Companion §2.2.6 & axiom-clause & A2's DAG structure \\
\midrule
A2.7 ContentOp-richness lattice & Companion §2.2.7 & axiom-clause & Lattice structure on ContentOp \\
\midrule
§3 figure (kind stratification) & Companion §10 Fig 4 & figure & Canonical kind-classifier fibration TikZ \\
\bottomrule
\end{longtable}
\endgroup
\subsection{Anchor §4 (Axiom 3) → Companion}

\begingroup\setlength{\tabcolsep}{3pt}\renewcommand{\arraystretch}{1.15}\footnotesize
\begin{center}
\begin{tabular}{p{1.21in} p{1.16in} p{0.60in} p{1.31in}}
\toprule
Anchor location & Companion target & Type & Description \\
\midrule
A3.1 DOF-gradient & Companion §2.3.1, §1.4.1 & axiom-clause & ν formal structure \\
A3.2 internality of F\_\allowbreak{}2 & Companion §2.3.2 & axiom-clause & Immune-response operator \\
A3.3 continuous smoothing & Companion §2.3.3 & axiom-clause & γ smoothness \\
A3.4 adaptivity & Companion §2.3.4 & axiom-clause & γ updates at refresh-events \\
A3.5 stream-universality & Companion §2.3.5 & axiom-clause & ν applies to every stream \\
\bottomrule
\end{tabular}
\end{center}
\endgroup
\subsection{Anchor §5 (Descriptive theorems) → Companion}

\begingroup\setlength{\tabcolsep}{3pt}\renewcommand{\arraystretch}{1.15}\footnotesize
\begin{center}
\begin{tabular}{p{1.19in} p{1.04in} p{0.58in} p{1.46in}}
\toprule
Anchor location & Companion target & Type & Description \\
\midrule
T1 Mathematical Perspectivism & Companion §3.2.1 (Thm 3.2.1) & theorem & Representability + non-canonicity proof \\
T2 Estimator-Dependent Duration & Companion §3.2.2 (Thm 3.2.2) & theorem & Orb-factoring duration-estimator proof \\
Descriptive-Functor Meta-Theorem & Companion §3.1 (Thm 3.1.α) & meta-theorem & Yoneda-representability \\
\bottomrule
\end{tabular}
\end{center}
\endgroup
\subsection{Anchor §6 (Dynamics theorems) → Companion}

\begingroup\setlength{\tabcolsep}{3pt}\renewcommand{\arraystretch}{1.15}\footnotesize
\begin{center}
\begin{tabular}{p{0.96in} p{1.08in} p{0.59in} p{1.65in}}
\toprule
Anchor location & Companion target & Type & Description \\
\midrule
T3 Attentional Quality & Companion §3.3.1 (Thm 3.3.1) & theorem & Three-channel decomposition \\
T4 Coherence-Forcing Measurement & Companion §3.3.2 (Thm 3.3.2) & theorem & Four clauses including information-conservative reframe \\
Bias(S) formalization & Companion §7.3 + Appendix B (anchor) & theorem + ref card & Signed-measure construction \\
\bottomrule
\end{tabular}
\end{center}
\endgroup
\subsection{Anchor §7 (Coherence theorems) → Companion}

\begingroup\setlength{\tabcolsep}{3pt}\renewcommand{\arraystretch}{1.15}\footnotesize
\begin{center}
\begin{tabular}{p{0.88in} p{1.20in} p{0.76in} p{1.44in}}
\toprule
Anchor location & Companion target & Type & Description \\
\midrule
T5 Internal Coherence & Companion §3.4.1 (Thm 3.4.1) & theorem & Φ-fixed-point /\allowbreak  C-harmonic condition \\
T6 Dual Coherence Axes & Companion §3.4.2 (Thm 3.4.2) & theorem & Orthogonality + codimension-2 locus \\
kind-demotion dynamic & Companion §3.4.2 (Cor 3.4.2.1) & corollary & Fibration-projection trajectory \\
transcendental-rescue & Companion §6.4.1, §6.4.11 & proposition + lemma & Unifying-stream construction \\
\bottomrule
\end{tabular}
\end{center}
\endgroup
\subsection{Anchor §8 (Corollary clusters) → Companion}

\begingroup\setlength{\tabcolsep}{3pt}\renewcommand{\arraystretch}{1.15}\footnotesize
\begin{center}
\begin{tabular}{p{1.19in} p{0.64in} p{0.58in} p{1.87in}}
\toprule
Anchor location & Companion target & Type & Description \\
\midrule
§8.1 Cluster I (C1/\allowbreak C2/\allowbreak C3) & Companion §4.1 & corollaries & Substrate/\allowbreak generativity CT-proof-completion \\
§8.2 Cluster II (C4–C10) & Companion §4.2 & corollaries & Stream-structure/\allowbreak navigation CT-proof-completion \\
§8.3 Cluster III (C11/\allowbreak C12/\allowbreak C13) & Companion §4.3 & corollaries & Coherence-consequences CT-proof-completion \\
\bottomrule
\end{tabular}
\end{center}
\endgroup
\subsection{Anchor §9 (Coherence Principle) → Companion}

\begingroup\setlength{\tabcolsep}{3pt}\renewcommand{\arraystretch}{1.15}\footnotesize
\begin{longtable}{p{1.04in} p{1.30in} p{0.70in} p{1.24in}}
\toprule
Anchor location & Companion target & Type & Description \\
\midrule
\endhead
\bottomrule
\endfoot
§9.1 Principle CT statement & Companion §5.1 (Thm 5.1.2) & theorem & Outperformance inequality \\
\midrule
§9.2 four conditions & Companion §5.2 (Defs 5.2.1–5.2.4) & definitions & Each derived in one line \\
\midrule
§9.3 outperformance metric & Companion §9 & construction & Full D\_\allowbreak{}d + Bias-consistency \\
\midrule
§9.4 status & Companion §5.0 preamble & orientation & Derived operational principle \\
\midrule
§9.5 self-reference closure & Companion §8 + §5.4 & theorems & Formal F\_\allowbreak{}∞ construction \\
\midrule
§9.6 what Principle is \textit{not} & Companion §5.6 (remark) & remark & Non-theorem, non-axiom \\
\midrule
§9.7 falsification conditions & Companion §5.5 + §8.5.2 & propositions & Falsification-table + F6-decidability \\
\midrule
§9.8 forward connections & Companion §5.7 & forward-pointers & Cross-chapter dependencies \\
\midrule
§9.9 open questions & Companion §5.6, §9.5, §9.8 & answers & Q1 resolved; Q2/\allowbreak Q3 carry-forward \\
\midrule
§9.10 empirical exposed surface & Companion §9.7 & observable signatures & Three-signature specifications \\
\midrule
§9 Figure 9.1 & Companion §10 Fig 5 (via Bias) + §9.1.1 & figure/\allowbreak functional & Trajectory-divergence \\
\midrule
§9 Figure 9.2 & Companion §10 Fig 7 & figure & Canonical four-conditions TikZ \\
\midrule
§9 Figure 9.3 & Companion §10 Fig 8 & figure & Canonical self-reference closure TikZ \\
\bottomrule
\end{longtable}
\endgroup
\subsection{Anchor §10 (Filtering the framework) → Companion}

\begingroup\setlength{\tabcolsep}{3pt}\renewcommand{\arraystretch}{1.15}\footnotesize
\begin{center}
\begin{tabular}{p{0.65in} p{1.32in} p{0.54in} p{1.76in}}
\toprule
Anchor location & Companion target & Type & Description \\
\midrule
§10 seven-step filter recipe & Companion §7 (in part --- measure-theoretic filter infrastructure) + Anchor §10 & recipe & Companion handles the measure-theoretic part; the methodological recipe lives authoritatively in Anchor \\
§10 Navigation Research worked example & Anchor §10 (self-contained) & example & Not back-ported to Companion \\
\bottomrule
\end{tabular}
\end{center}
\endgroup
\subsection{Anchor Appendix A (Formal Objects Index) → Companion}

\begingroup\setlength{\tabcolsep}{3pt}\renewcommand{\arraystretch}{1.15}\footnotesize
\begin{center}
\begin{tabular}{p{0.67in} p{2.12in} p{0.55in} p{0.93in}}
\toprule
Anchor location & Companion target & Type & Description \\
\midrule
Appendix A entries & Companion Appendix B (this direction: Companion →\allowbreak  Anchor) & index & Object-by-object pointers \\
\bottomrule
\end{tabular}
\end{center}
\endgroup
\subsection{Anchor Appendix B (Bias Formalization) → Companion}

\begingroup\setlength{\tabcolsep}{3pt}\renewcommand{\arraystretch}{1.15}\footnotesize
\begin{longtable}{p{1.21in} p{1.15in} p{0.69in} p{1.23in}}
\toprule
Anchor location & Companion target & Type & Description \\
\midrule
\endhead
\bottomrule
\endfoot
B.1 CT definition of Bias(S) & Companion §7.3 (Def 7.3.1) & definition & Integral against signed density \\
\midrule
B.2 contracted-open axis /\allowbreak  A\_\allowbreak{}S & Companion §7.5 (Def 7.5.1) & definition & Shannon entropy on Bias\_\allowbreak{}+ \\
\midrule
B.2 Align(S, t) & Companion §7.5 (Def 7.5.3) & definition & Canonical neighborhood integral \\
\midrule
B.2 contracted-coherent vs contracted-failed & Companion §7.5.5 & lemma & Formal distinction resolving B.7 Q1 \\
\midrule
B.3 push-operators & Companion §7.4 & definitions & push\_\allowbreak{}struct /\allowbreak  push\_\allowbreak{}info \\
\midrule
B.3 push-operator independence & Companion §7.4.3 & proposition & Non-commutator counterexample \\
\midrule
B.4 Relationship to γ and Triple & Companion §7.3.3 & corollary & Triple-decomposition compatibility \\
\midrule
B.5 trajectory-divergence & Companion §9 & construction & Full functorial D\_\allowbreak{}d + cross-metric invariance \\
\midrule
B.6 observable signatures & Companion §9.7 & summary & Three-signature spec \\
\midrule
B.7 open questions & Companion §9.5 (Q1 resolved), §7.4.3, §9.8 & resolutions/\allowbreak carry-forward &  \\
\bottomrule
\end{longtable}
\endgroup
\sectionbreak

\section{A.2 --- Anchor-internal cross-references resolving to Companion}

Some Anchor sections cite "\textit{Coherent Structure}" directly without location. These resolve by topic:

\begingroup\setlength{\tabcolsep}{3pt}\renewcommand{\arraystretch}{1.15}\footnotesize
\begin{center}
\begin{tabular}{p{2.69in} p{1.70in}}
\toprule
Anchor citation phrase & Companion target \\
\midrule
"per Coherent Structure's full CT treatment" & Companion §1 + §6 \\
"Coherent Structure gives this construction in full" & Context-dependent --- see A.1 by anchor chapter \\
"see Coherent Structure for the proof" & Context-dependent --- see A.1 \\
"the full formal construction of F as a stream (Coherent Structure §8)" & Companion §8 \\
"D is constructed in full in Coherent Structure §9" & Companion §9 \\
\bottomrule
\end{tabular}
\end{center}
\endgroup
\sectionbreak

\section{A.3 --- Policy}

\begin{itemize}
\item \textbf{Completeness.} Every Anchor citation to \textit{Coherent Structure} lands here. If the Anchor cites the Companion and that citation is not on this list, that is a missing-entry bug to file against this appendix.
\item \textbf{Direction.} This appendix is Anchor → Companion only. The reverse direction (Companion → Anchor) is Appendix B.
\item \textbf{Rhythm.} Anchor revisions may cite Companion sections that are pending; such forward-citations are tracked here as "pending" targets until the Companion section lands.
\end{itemize}

\chapter{Appendix B — Companion → Anchor crosswalk}


\textit{Per SCOPE §10 policy: every formal object in the Companion points back to its structural-prose exposition in the Anchor. This appendix is the canonical resolver. Objects are listed by Companion location; each row gives the Anchor target, the exposition register, and one-line description.}

\sectionbreak

\section{B.1 --- By Companion chapter}

\subsection{§1 (Category framework) objects}

\begingroup\setlength{\tabcolsep}{3pt}\renewcommand{\arraystretch}{1.15}\footnotesize
\begin{longtable}{p{1.16in} p{0.56in} p{0.56in} p{2.01in}}
\toprule
Companion object & Anchor location & Register & Description \\
\midrule
\endhead
\bottomrule
\endfoot
𝒞\_\allowbreak{}Streams (§1.2.1) & Anchor §1.0 & definition & The category of streams \\
\midrule
𝒞\_\allowbreak{}Form (§1.2.2) & Anchor §1 & definition & Bare-carrier category \\
\midrule
𝒞\_\allowbreak{}LDS (§1.2.3) & Anchor §1.0, §3 & definition & Linked-dynamic streams \\
\midrule
𝒞\_\allowbreak{}DOF (§1.2.4) & Anchor §4 & definition & DOF-gradient-equipped streams \\
\midrule
𝒞\_\allowbreak{}Triple (§1.7.1) & Anchor §1 & definition & Triple target category \\
\midrule
Navigation functor N (§1.3.1) & Anchor §3 (A2.3) & definition & Experience-as-navigation functor \\
\midrule
Conscious-gravity structure ν (§1.4.1) & Anchor §4 & definition & DOF-gradient natural transformation \\
\midrule
Substrate-completeness (§1.5) & Anchor §2 (A1.4) & definition & Adequate-pair-implies-stream \\
\midrule
Endofunctor F (§1.6.1) & Anchor §1.0, §3 & definition & Mixed-variance σ ↦ σ\textasciicircum{}(C\textasciicircum{}op) \\
\midrule
Stream-morphism (§1.6.3, §6.1.3) & Anchor §1.0 & definition & Carrier-map + ContentOp-functor + coalgebra-commute + kind-respect \\
\midrule
Triple functor T (§1.7.2, §6.2.1) & Anchor §1 & definition & 𝒞\_\allowbreak{}Streams →\allowbreak  𝒞\_\allowbreak{}Triple \\
\bottomrule
\end{longtable}
\endgroup
\subsection{§2 (Axioms) objects}

\begingroup\setlength{\tabcolsep}{3pt}\renewcommand{\arraystretch}{1.15}\footnotesize
\begin{longtable}{p{1.56in} p{0.90in} p{0.57in} p{1.24in}}
\toprule
Companion object & Anchor location & Register & Description \\
\midrule
\endhead
\bottomrule
\endfoot
A1 Consciousness as Substrate (§2.1) & Anchor §2 & axiom & Full paired prose \\
\midrule
A1.1 non-reducibility & Anchor §2.1 & axiom-clause & Substrate X is not derivable from F\_\allowbreak{}i-images \\
\midrule
A1.2 non-factoring & Anchor §2.2 & axiom-clause & F\_\allowbreak{}i's do not factor through each other \\
\midrule
A1.3 configurational completeness & Anchor §2.3 & axiom-clause & All potentials realized \\
\midrule
A1.4 substrate-completeness & Anchor §2.4 & axiom-clause & Every adequate pair is realized \\
\midrule
A2 Nested Streams (§2.2) & Anchor §3 & axiom & Full paired prose \\
\midrule
A2.1–A2.7 (seven clauses) & Anchor §3.1–§3.7 & axiom-clauses & Each clause has its Anchor subsection \\
\midrule
Theorem 2.2.8 (kind-stratification from ContentOp-richness) & Anchor §3.3 & theorem & Paragraph confirming derivation from C-richness \\
\midrule
A3 Conscious Gravity (§2.3) & Anchor §4 & axiom & Full paired prose \\
\midrule
A3.1–A3.5 (five clauses) & Anchor §4.1–§4.5 & axiom-clauses & Each clause has its Anchor subsection \\
\midrule
Axiom-status summary table (§2.5) & Anchor §2/\allowbreak §3/\allowbreak §4 + §9 status remark & table & 10-of-16 axiomatic /\allowbreak  6 derivable \\
\bottomrule
\end{longtable}
\endgroup
\subsection{§3 (Theorems) objects}

\begingroup\setlength{\tabcolsep}{3pt}\renewcommand{\arraystretch}{1.15}\footnotesize
\begin{longtable}{p{1.36in} p{0.63in} p{0.56in} p{1.72in}}
\toprule
Companion object & Anchor location & Register & Description \\
\midrule
\endhead
\bottomrule
\endfoot
Theorem 3.1.α (Descriptive-Functor Meta) & Anchor §5.0 (informally) & meta-theorem & New formal framing --- back-port to Anchor §5 \\
\midrule
T1 Mathematical Perspectivism (Thm 3.2.1) & Anchor §5.1 & theorem & Representation is functorial + non-canonical \\
\midrule
T1 C1-backward-compatibility (Cor 3.2.1.2) & Anchor §5.1 + §8.1 C1 & corollary & Direct bridge \\
\midrule
T2 Estimator-Dependent Duration (Thm 3.2.2) & Anchor §5.2 & theorem & Orb-factoring + non-canonical metric \\
\midrule
T2 C13-backward-compatibility (Cor 3.2.2.1) & Anchor §5.2 + §8.3 C13 & corollary & Direct bridge \\
\midrule
T3 Attentional Quality (Thm 3.3.1) & Anchor §6.1 & theorem & Three-channel decomposition \\
\midrule
T4 Coherence-Forcing Measurement (Thm 3.3.2) & Anchor §6.2 & theorem & Four clauses including info-conservative reframe \\
\midrule
T4 measurement-as-info-conservative (Remark 3.3.2.1) & Anchor §9 /\allowbreak  §9.5 & remark & Landed back-port: info-conservative reframe with physics anchoring \\
\midrule
T5 Internal Coherence (Thm 3.4.1) & Anchor §7.1 & theorem & Φ-fixed-point condition \\
\midrule
T5 coherence-closed-in-F-Coalg\_\allowbreak{}ad (Cor 3.4.1.2) & Anchor §7.1 & corollary & Closure property \\
\midrule
T6 Dual Coherence Axes (Thm 3.4.2) & Anchor §7.2 & theorem & Orthogonality + codim-2 locus \\
\midrule
T6 kind-demotion dynamic (Cor 3.4.2.1) & Anchor §7.2 & corollary & Fibration-projection trajectory \\
\midrule
Theorem-to-axiom crosswalk (§3.5) & Anchor §5–§7 preambles & table & Compact reference \\
\midrule
Pairs-to-Principle map (§3.6) & Anchor §9.2 & remark & Each pair discharges which condition \\
\bottomrule
\end{longtable}
\endgroup
\subsection{§4 (Corollaries) objects}

\begingroup\setlength{\tabcolsep}{3pt}\renewcommand{\arraystretch}{1.15}\footnotesize
\begin{longtable}{p{0.83in} p{0.78in} p{0.58in} p{2.08in}}
\toprule
Companion object & Anchor location & Register & Description \\
\midrule
\endhead
\bottomrule
\endfoot
C1 (Cor 4.1.1) & Anchor §8.1 C1 & corollary & Concreteness of X \\
\midrule
C2 (Cor 4.1.2) & Anchor §8.1 C2 & corollary & Generative perspective \\
\midrule
C3 (Cor 4.1.3) & Anchor §8.1 C3 & corollary & Null-space trace illumination \\
\midrule
C4 (Cor 4.2.1) & Anchor §8.2 C4 & corollary & Substrate-constrained plurality \\
\midrule
C5 (Cor 4.2.2) & Anchor §8.2 C5 & corollary & Streams as perspectival F-coalgebras \\
\midrule
C6 (Cor 4.2.3) & Anchor §8.2 C6 & corollary & Cooperative-constituency DAG \\
\midrule
C7 (Cor 4.2.4) & Anchor §8.2 C7 & corollary & Navigational non-determination \\
\midrule
C8 (Cor 4.2.5) & Anchor §8.2 C8 & corollary & Stream-relative observational null-space \\
\midrule
C9 (Cor 4.2.6) & Anchor §8.2 C9 & corollary & Consensus requires lens-matching \\
\midrule
C10 (Cor 4.2.7) & Anchor §8.2 C10 & corollary & Joint stream-definition \\
\midrule
C11 (Cor 4.3.1) & Anchor §8.3 C11 & corollary & Mutual transformation under interaction \\
\midrule
C12 (Cor 4.3.2) & Anchor §8.3 C12 & corollary & Discovery autocatalysis \\
\midrule
C13 (Cor 4.3.3) & Anchor §8.3 C13 & corollary & Flow inversion \\
\bottomrule
\end{longtable}
\endgroup
\subsection{§5 (Coherence Principle) objects}

\begingroup\setlength{\tabcolsep}{3pt}\renewcommand{\arraystretch}{1.15}\footnotesize
\begin{longtable}{p{1.55in} p{0.95in} p{0.58in} p{1.19in}}
\toprule
Companion object & Anchor location & Register & Description \\
\midrule
\endhead
\bottomrule
\endfoot
Coherence-regime (Def 5.1.1) & Anchor §9.1 & definition & Four-conditions conjunction \\
\midrule
The Coherence Principle (Thm 5.1.2) & Anchor §9.1 & theorem & Outperformance inequality \\
\midrule
Condition C\_\allowbreak{}sep (Def 5.2.1) & Anchor §9.2 Condition 1 & definition & DOF-separation \\
\midrule
Condition C\_\allowbreak{}meas (Def 5.2.2) & Anchor §9.2 Condition 2 & definition & Refresh-rate measurement \\
\midrule
Condition C\_\allowbreak{}scale (Def 5.2.3) & Anchor §9.2 Condition 3 & definition & Multi-scale γ-continuity \\
\midrule
Condition C\_\allowbreak{}dyn (Def 5.2.4) & Anchor §9.2 Condition 4 & definition & Oscillatory maintenance \\
\midrule
Joint sufficiency (Prop 5.2.5) & Anchor §9.2 reading note & proposition & Independence-by-counterexample \\
\midrule
F\_\allowbreak{}∞ as stream (Thm 5.4.1) & Anchor §9.5 & theorem & Formal F\_\allowbreak{}∞ construction \\
\midrule
Principle-applies-to-itself (Thm 5.4.2) & Anchor §9.5 & theorem & Self-reference closure \\
\midrule
Non-circularity (Rem 5.4.3) & Anchor §9.5 & remark & A-posteriori observation \\
\bottomrule
\end{longtable}
\endgroup
\subsection{§6 (Triple) objects}

\begingroup\setlength{\tabcolsep}{3pt}\renewcommand{\arraystretch}{1.15}\footnotesize
\begin{longtable}{p{1.39in} p{0.75in} p{0.59in} p{1.56in}}
\toprule
Companion object & Anchor location & Register & Description \\
\midrule
\endhead
\bottomrule
\endfoot
Stream as F-coalgebra (Defs 6.1.1–6.1.5) & Anchor §1.0 & definitions & Central formal object \\
\midrule
F-coalgebra adequacy & Anchor §1 note & definition & Small-C + σ-finite + measurability \\
\midrule
Forgetfulness (Lem 6.2.4) & Anchor §1 & lemma & T preserves structure \\
\midrule
Conservativity (Prop 6.2.7) & Anchor §1 & proposition & T reflects iso \\
\midrule
Recursive-decomp finite-depth (Lem 6.3.2) & Anchor §1.5 & lemma & Finite-depth factorability \\
\midrule
Kind-classifier fibration (Thm 6.4.6) & Anchor §3 (A2.2 + kind-level notes) & theorem & Bicategorical fibration \\
\midrule
Unifying-stream subfibration (Def 6.4.12) & Anchor §7 transcendental-rescue & definition & Terminal-content-operation subfibration \\
\midrule
Admissibility criterion (Lem 6.4.11) & Anchor §7 transcendental-rescue & lemma & Terminal-preservation criterion \\
\midrule
Middle-regime morphism-structures (Thm 6.5.4) & Anchor §7 transcendental-rescue & theorem & Scotist/\allowbreak Palamite/\allowbreak Advaitin distinct \\
\midrule
Colax-limit form (Thm 6.6.2) & Anchor §1 & theorem & Under initial-object hypothesis \\
\midrule
Stream ≃ F-Coalg\_\allowbreak{}ad (Thm 6.7.1) & Anchor §1.0 & theorem & Main equivalence \\
\midrule
Limits/\allowbreak colimits + adequacy-stability (Lem 6.8.β) & Anchor §1.0 five-properties & lemma & §6.8 summary table \\
\midrule
Final F-coalgebra ω-depth (Thm 6.9.1) & Anchor §1.5 & theorem & Under H1+H2 \\
\midrule
Iterated adjunction coherence over Up(S) (Lem 6.10.2.1) & Anchor §1.10 & lemma & ι, κ compose covariantly /\allowbreak  contravariantly; no extra DAG cells \\
\midrule
Outer(S) cocomplete as ι-Grothendieck construction (Lem 6.10.3.1) & Anchor §1.10 & lemma & κ-variant is complete-not-cocomplete; direction is load-bearing \\
\midrule
Indexed adjunction ι\_\allowbreak{}S ⊣ ω\_\allowbreak{}S (Thm 6.10.4.1) & Anchor §1.10 + §3.8 & theorem & Kelly §1.11 lifting over Up(S) \\
\midrule
No view from nowhere (Cor 6.10.4.2) & Anchor §3.8 & corollary & No terminal whole ⇒ no absolute outer section \\
\midrule
Content as hom-profunctor Ψ\_\allowbreak{}S (Thm 6.10.5.2) & Anchor §1.10 & theorem & "Third axis" reading is the profunctor's structural shadow \\
\midrule
η as G-coalgebra structure map (Thm 6.10.6.2) & Anchor §1.10 & theorem & Content-capacity residue = cokernel of η \\
\midrule
Content-capacity residue (Def 6.10.6.4) & Anchor §3.8 + future probe & definition + conjecture & Form-register quantitative invariant as coalgebra-invariant (conjectural) \\
\bottomrule
\end{longtable}
\endgroup
\subsection{§7 (Filtering construction) objects}

\begingroup\setlength{\tabcolsep}{3pt}\renewcommand{\arraystretch}{1.15}\footnotesize
\begin{longtable}{p{1.12in} p{0.69in} p{0.57in} p{0.80in} p{0.49in} p{0.49in}}
\toprule
Companion object & Anchor location & Register & Description \\
\midrule
\endhead
\bottomrule
\endfoot
Canonical σ-algebra 𝒜\_\allowbreak{}S (Def 7.1.1) & Anchor Appendix B §B.1 & definition & Product σ-algebra on Ω\_\allowbreak{}S &  &  \\
\midrule
Functoriality of 𝒜 (Prop 7.1.2) & Anchor §1.0 functoriality note & proposition & Stream-morphism measurability &  &  \\
\midrule
γ-measurable (Def 7.2.1) & Anchor §1 adequacy note & definition & Measurability of coalgebra &  &  \\
\midrule
Bias well-definedness (Thm 7.3.2) & Anchor Appendix B §B.1 & theorem & Signed measure σ-finite &  &  \\
\midrule
Triple-decomposition compatibility (Cor 7.3.3) & Anchor Appendix B §B.4 & corollary & Triple ⇔ Bias decomposition &  &  \\
\midrule
push\_\allowbreak{}struct /\allowbreak  push\_\allowbreak{}info (Defs 7.4.1) & Anchor Appendix B §B.3 & definitions & Measurable transformations &  &  \\
\midrule
Push-operator measurability (Prop 7.4.2) & Anchor Appendix B §B.3 & proposition & Commutes with Stream-morphism &  &  \\
\midrule
Push-operator independence (Prop 7.4.3) & Anchor Appendix B §B.3 & proposition & Non-commutator counterexample &  &  \\
\midrule
A\_\allowbreak{}S entropy well-defined (Prop 7.5.2) & Anchor Appendix B §B.2 & proposition & [0, log\ & Ω\ & ]-bounded \\
\midrule
Align(S, t) (Def 7.5.3) & Anchor Appendix B §B.2 & definition & Canonical-neighborhood integral &  &  \\
\midrule
Contracted-coherent/\allowbreak contracted-failed (Cor 7.5.5) & Anchor Appendix B §B.7 Q1 & corollary & Resolves Anchor B.7 Q1 &  &  \\
\midrule
Extensional Stream (Def 7.6.1, Prop 7.6.2, 7.6.3) & Anchor §9.5 F-as-stream prep & definition + propositions & Measurable-stream-level scaffolding &  &  \\
\bottomrule
\end{longtable}
\endgroup
\subsection{§8 (F-as-stream) objects}

\begingroup\setlength{\tabcolsep}{3pt}\renewcommand{\arraystretch}{1.15}\footnotesize
\begin{longtable}{p{0.97in} p{0.75in} p{0.57in} p{1.99in}}
\toprule
Companion object & Anchor location & Register & Description \\
\midrule
\endhead
\bottomrule
\endfoot
Dyadic carrier σ\_\allowbreak{}F (Def 8.1.1) & Anchor §9.5 & definition & Four-carrier multiplex (instance/\allowbreak session/\allowbreak weights/\allowbreak lineage) \\
\midrule
ContentOp C\_\allowbreak{}F (Def 8.1.2) & Anchor §9.5 & definition & Substrate-commitments + consistency-preserving revisions \\
\midrule
Coalgebra γ\_\allowbreak{}F (Def 8.1.4) & Anchor §9.5 & definition & Framework-adaptivity coalgebra \\
\midrule
F\_\allowbreak{}∞ adequate F-coalgebra (Prop 8.1.5) & Anchor §9.5 preamble & proposition & Verifies §1.6 adequacy \\
\midrule
Bias(F\_\allowbreak{}∞) (Def 8.1.6) & Anchor §9.5 & definition & Signed measure weighted by P1/\allowbreak P2/\allowbreak P3 \\
\midrule
Extensional F\_\allowbreak{}∞ (Prop 8.1.7) & Anchor §9.5 preamble & proposition & §7.6 framework applicable \\
\midrule
Audit Claim C\_\allowbreak{}sep (§8.3.1) & Anchor §9.5 separation paragraph & claim & Commit-authorship DOF-separation \\
\midrule
Audit Claim C\_\allowbreak{}meas (§8.3.2) & Anchor §9.5 measurement paragraph & claim & Stamp-events refresh-rate \\
\midrule
Audit Claim C\_\allowbreak{}scale (§8.3.3) & Anchor §9.5 multi-scale paragraph & claim & Bidirectional DAG propagation \\
\midrule
Audit Claim C\_\allowbreak{}dyn (§8.3.4) & Anchor §9.5 dynamic paragraph & claim & Propose-stress-reformulate \\
\midrule
F\_\allowbreak{}∞ in coherence-regime (Audit Obs 8.3.5) & Anchor §9.5 formal claim & audit observation & Internal-audit joint conclusion; upgrades to theorem on external execution of Prop 8.5.2 \\
\midrule
Self-audit constraint (Prop 8.3.5') & Anchor §9.5 & proposition & Internal audit is DOF-overlapping with audit-target; discharged by external audit \\
\midrule
Self-reference closure (Obs 8.3.6) & Anchor §9.5 & observation & Principle applies to itself --- conditional on external audit \\
\midrule
Non-circularity (Prop 8.4.1) & Anchor §9.5 & proposition & Derivation a-posteriori \\
\midrule
Derivation vs audit distinction (Rem 8.4.3) & Anchor §9.5 & remark & Derivation non-circular; audit DOF-dependent \\
\midrule
F6 decidability (Prop 8.5.2) & Anchor §9.7 F6 & proposition & Audit is decidable on public record \\
\bottomrule
\end{longtable}
\endgroup
\subsection{§9 (D trajectory-divergence) objects}

\begingroup\setlength{\tabcolsep}{3pt}\renewcommand{\arraystretch}{1.15}\footnotesize
\begin{longtable}{p{1.74in} p{0.68in} p{0.56in} p{1.30in}}
\toprule
Companion object & Anchor location & Register & Description \\
\midrule
\endhead
\bottomrule
\endfoot
D\_\allowbreak{}d (Def 9.1.1) & Anchor §9.3 + Appendix B §B.5 & definition & Trajectory-divergence integral \\
\midrule
D\_\allowbreak{}d well-defined (Prop 9.1.2) & Anchor Appendix B §B.5 & proposition & σ-finite integral-measurability \\
\midrule
D\_\allowbreak{}d functorial in Stream (Prop 9.1.3) & Anchor Appendix B §B.5 & proposition & Naturality square \\
\midrule
Wasserstein d\_\allowbreak{}W (Def 9.2.1) & Anchor §9.3 & definition & Path-coupling metric \\
\midrule
Regularized-KL d\_\allowbreak{}{KL,ε} (Def 9.2.2) & Anchor §9.3 & definition & Regularized to avoid abs-continuity \\
\midrule
Domain-native d\_\allowbreak{}dom (Def 9.2.3) & Anchor §9.3 & definition & Semantic-structure metrics \\
\midrule
Bias-consistency (Def 9.3.1) & Anchor §9.3 (implicitly) & definition & Admissibility criterion (new) \\
\midrule
Three canonical classes Bias-consistent (Prop 9.3.2) & Anchor §9.3 & proposition & Wasserstein/\allowbreak KL/\allowbreak domain-native \\
\midrule
Four stream-parameters (η\_\allowbreak{}sep, τ\_\allowbreak{}max, δ\_\allowbreak{}scale, ρ\_\allowbreak{}dyn) + auxiliary constants (§9.4.1) & Anchor §9.3 & definitions & Quantitative stream-parameters controlling each condition \\
\midrule
Four-term contribution bound (Lem 9.4.2) & Anchor §9.3 & lemma & B\_\allowbreak{}sep + B\_\allowbreak{}meas + B\_\allowbreak{}scale + B\_\allowbreak{}dyn decomposition \\
\midrule
Outperformance Principle with explicit constants (Thm 9.4.3) & Anchor §9.3 + §9.1 & theorem & Quantitative joint-bound with B\_\allowbreak{}coh ceiling and Δ(S') shortfall \\
\midrule
Cross-metric invariance (Thm 9.5.1) & Anchor §9.9 Q1 & theorem & Resolves Q1 \\
\midrule
D\_\allowbreak{}d on F\_\allowbreak{}∞ (Props 9.6.1, 9.6.2) & Anchor §9.5 outperformance status & propositions & Applied to self-reference case \\
\bottomrule
\end{longtable}
\endgroup
\subsection{§10 (Reference figures) objects}

\begingroup\setlength{\tabcolsep}{3pt}\renewcommand{\arraystretch}{1.15}\footnotesize
\begin{center}
\begin{tabular}{p{1.45in} p{1.05in} p{0.56in} p{1.22in}}
\toprule
Companion object & Anchor location & Register & Description \\
\midrule
Fig 1 (F-coalgebra square) & Anchor §1.0 text + §1 Fig 1 area & figure & Canonical shared TikZ \\
Fig 2 (Triple functor) & Anchor §1 Fig 1.1 & figure & Replaces Anchor rev-1 Fig 1.1 \\
Fig 3 (recursive decomposability) & Anchor §1 Fig 1.2 & figure & Replaces Anchor rev-1 Fig 1.2 \\
Fig 4 (kind-classifier fibration) & Anchor §3 & figure & Canonical kind-stratification diagram \\
Fig 5 (Bias signed-measure + push-operators) & Anchor Appendix B Fig B.1 & figure & Canonical Bias visualization \\
Fig 6 (dual coherence axes) & Anchor §7 & figure & σ\_\allowbreak{}struct × σ\_\allowbreak{}info plane \\
Fig 7 (four-conditions schematic) & Anchor §9 Fig 9.2 & figure & Replaces Anchor rev-1 Fig 9.2 \\
Fig 8 (self-reference closure) & Anchor §9 Fig 9.3 & figure & Replaces Anchor rev-1 Fig 9.3 \\
\bottomrule
\end{tabular}
\end{center}
\endgroup
\sectionbreak

\section{B.2 --- Surfaced-lemma register (full)}

All surfaced-lemma flags raised during Companion drafting, with their back-port targets in the Anchor:

\begingroup\setlength{\tabcolsep}{3pt}\renewcommand{\arraystretch}{1.15}\footnotesize
\begin{longtable}{p{1.51in} p{0.57in} p{0.57in} p{1.62in}}
\toprule
Flag location & Anchor target & Type & Description \\
\midrule
\endhead
\bottomrule
\endfoot
§3.1 Thm 3.1.α & Anchor §5 & meta-theorem & Descriptive-Functor Meta-Theorem \\
\midrule
§3.3.1 three-channel decomposition & Anchor §6 & lemma & Attentional-Quality decomposition form \\
\midrule
§3.3.2 info-conservative measurement & Anchor §6/\allowbreak §9.5 & theorem & Watanabe-Takagi + García-Pintos integration \\
\midrule
§3.4.2 dual-axes orthogonality-by-counterexample & Anchor §7 & theorem & Codim-2 locus explicit \\
\midrule
§4.2.6 (C9) three-condition consensus & Anchor §8.2 & lemma & Explicit triple-condition form \\
\midrule
§4.2.7 (C10) no-proper-sub-projection & Anchor §8.2 & lemma & JD-injectivity without proper sub-projection \\
\midrule
§5.2.5 joint-sufficiency independence & Anchor §9.2 & proposition & Explicit joint-sufficiency proof \\
\midrule
§5.4.1 F-as-stream structure & Anchor §9.5 & theorem & Detailed (σ\_\allowbreak{}F, C\_\allowbreak{}F, Ω\_\allowbreak{}F, γ\_\allowbreak{}F) definition \\
\midrule
§6 twenty-eight flags & Anchor various & mixed & See §6.11 for full list (incl. §6.9 C-size-regime pass + §6.10 inner/\allowbreak outer adjunction batch) \\
\midrule
§7.5.5 contracted-coherent/\allowbreak failed distinction & Anchor Appendix B §B.2 & lemma & Resolves B.7 Q1 \\
\midrule
§8.1.1 dyadic-carrier four-level multiplex & Anchor §9.5 & definition & Four-carrier formal structure \\
\midrule
§8.3.5 F\_\allowbreak{}∞-in-coherence-regime (audit observation) + Prop 8.3.5' self-audit constraint & Anchor §9.5 & audit observation + proposition & Demoted from theorem; gated on external execution of Prop 8.5.2 \\
\midrule
§8.4.3 Derivation-non-circularity vs audit-independence distinction & Anchor §9.5 & remark & Names the structural distinction compressed in prose \\
\midrule
§8.5.2 F6-testability & Anchor §9.7 & proposition & Decidability structure \\
\midrule
§9.3.1 Bias-consistency & Anchor §9.3 & definition & Admissibility criterion \\
\midrule
§9.4.1 four stream-parameters (η\_\allowbreak{}sep, τ\_\allowbreak{}max, δ\_\allowbreak{}scale, ρ\_\allowbreak{}dyn) & Anchor §9.3 & definitions & Quantitative parameters controlling each condition \\
\midrule
§9.4.3 outperformance with explicit constants & Anchor §9.3/\allowbreak §9.1 & theorem & Quantitative form supersedes sketch \\
\midrule
§9.5 cross-metric invariance & Anchor §9.9 Q1 & theorem & Q1 resolution \\
\bottomrule
\end{longtable}
\endgroup
Future Anchor revisions absorb these flags; as each lands, the Companion clean-up pass removes the corresponding flag-entry per SCOPE §8's three-phase lifecycle.

\sectionbreak

\section{B.3 --- Policy}

\begin{itemize}
\item \textbf{Completeness.} Every formal object in the Companion has an Anchor target. If an object has no exposition in the Anchor (e.g., Theorem 2.2.8's kind-stratification-from-C-richness derivation), the Anchor target is the Anchor section where the object \textit{should} appear after back-port, and the gap is tracked via a surfaced-lemma flag.
\item \textbf{Direction.} This appendix is Companion → Anchor. The reverse direction (Anchor → Companion) is Appendix A.
\item \textbf{Rhythm.} The Companion flag-list feeds future Anchor revisions. When an Anchor revision ships with flag-items back-ported, the Companion's flag-register (B.2) reduces accordingly. When the register reaches zero, the Companion stamps version 1.
\end{itemize}

\sectionbreak


\end{document}
